\input _template

\let\angle\sphericalangle

\setlanguage{SK}

\Title{Pridávanie bodov do obrázka}

\MathcompsLink{pridavanie-bodov}

\Author{Patrik Bak}

\sec Úvod

Všetci poznáme základné metódy riešenia olympiádnej geometrie, ako sú

\begitems
\i počítanie uhlov a~tetivové štvoruholníky,
\i zhodnosť a~podobnosť trojuholníkov,
\i mocnosť bodu a~chordály.
\enditems

Len ich kombináciou vieme vyriešiť mnoho úloh. Nie vždy je však ľahké prísť na to, ktorú techniku použiť. Všeobecné techniky, ktoré nám pomáhajú, sú

\begitems
\i preformulovanie zadania,
\i rozpoznanie známych podtvrdení,
\i eliminácia bodov,
\i pridávanie bodov do obrázka.
\enditems

V~tomto materiáli sa zameriame na tú poslednú a~precvičíme si niekoľko konkrétnych situácií, ktoré sa opakujú v~úlohách všetkých úrovní obťažnosti.

\sec Všeobecné tipy

\secc Kreslenie obrázkov

Jedna z~vecí, ktorá môže rozhodnúť o tom, či úlohu vyriešite alebo nie, je dobrý obrázok. Zopár tipov:

\begitems
\i kreslite veľké obrázky,
\i nekreslite rôznostranné trojuholníky ako rovnostranné,
\i používajte rysovacie pomôcky,
\i používajte dobre zastrúhané ceruzky a~gumu,
\i používajte farebné ceruzky (napr. na vyznačenie rovnako dlhých úsečiek),
\i nekreslite kružnice, kým nemusíte; ak štyri body ležia na kružnici, stačí si to poznačiť bokom, alebo zafarbiť hrany tetivového štvoruholníka,
\i pri kreslení kolmice použite namiesto trojuholníka s~ryskou dve kružnice (mne to vždy dávalo presnejšie výsledky),
\i nekreslite priamky a~oblúky dlhšie, než je nutné (viac rušivých prvkov znamená, že vidíte menej); dvakrát si rozmyslite, kým do obrázka nejakú priamku pridáte (alebo kým ju zvýrazníte či zafarbíte),
\i pomocné čiary priebežne mažte,
\i pri kreslení presného obrázka smiete použiť tvrdenie, ktoré sa snažíte dokázať, alebo podtvrdenia, ktoré ste už objavili,
\i prekresľujte často (napr. po preformulovaní), pokojne ten istý obrázok nakreslite viackrát (napr. na overenie hypotéz),
\i ak presný obrázok nakresliť nejde (napr. úloha má záludné implicitne definované body), nakreslite presne aspoň takú jeho časť, akú viete,
\i pri kreslení trojuholníka~$ABC$ umiestnite horný vrchol tak, aby situácia vyzerala vizuálne symetricky (v~dobre postavených úlohách je to vrchol~$A$),
\i niekedy je užitočné umiestniť nejakú netradičnú priamku zvisle alebo vodorovne (napr. uhlopriečku štvoruholníka),
\i najprv nakreslite hrubý náčrt od ruky a~s~pomôckami začnite až po krátkom zoznámení sa so situáciou (pokiaľ úlohu nevyriešite už z~náčrtu); získate tak cit pre to, aký veľký má obrázok byť a~kde ho začať kresliť.
\enditems

\secc Ako rozmýšľať pri riešení úloh

Keď dostaneme úlohu a~nakreslíme obrázok, začneme analyzovať:

\begitems
\i Predpokladajme, že tvrdenie, ktoré máme dokázať, platí. Vieme z~toho odvodiť niečo zaujímavé?
\i Ktoré predpoklady o~definovaných bodoch sme použili a~ktoré nevieme uchopiť? Ide o~uhlové alebo dĺžkové podmienky?
\i Ktorá časť úlohy sa dá najťažšie uchopiť? Ktoré body pôsobia najľubovoľnejšie?
\i Predstavte si riešenie. Z~ktorého podtvrdenia by tvrdenie mohlo vyplynúť?
\i Ktoré uhly určujú konfiguráciu? Ktoré uhly vieme pomocou nich vyjadriť?
\i Vieme nejaký bod eliminovať preformulovaním situácie?
\enditems

V~praxi samotná analýza nie vždy prinesie úspech, inak by to bolo príliš ľahké~\SlightSmile{} So situáciou sa dá ešte pohrať:
\begitems
\i Naslepo počítajte nejaké uhly a~pozerajte, čo sa stane. Všímajte si, či sme náhodou nenašli tetivový štvoruholník alebo inú zaujímavú konfiguráciu.
\i Naslepo napíšte nejaké pomery (napr. z~podobných trojuholníkov) a~skúste ich kombinovať so známymi rovnosťami dĺžok. Všímajte si, či sa novozískané vzťahy dajú interpretovať ako podobnosť dvoch iných trojuholníkov.
\i Definujte nejaké pekne vyzerajúce body, ktoré by mohli situáciu objasniť alebo o~nej prezradiť niečo zaujímavé (to je to, čo si v~tomto materiáli precvičíme~\SlightSmile{}).
\i Skúste, ako situácia vyzerá v~špeciálnych prípadoch, alebo skúmajte voľné body v~limitných polohách.
\enditems

Môžeme si aj vytvárať hypotézy. Dobrá hypotéza je v~mnohých prípadoch kľúčom k~riešeniu ťažkých úloh. Autori úloh sa snažia zo všetkých síl kľúčové tvrdenia skryť a~človek ich často jednoducho musí uhádnuť. Dobre nakreslený obrázok hádanie uľahčuje. Metódy hádania:

\begitems
\i Hádajte pekné geometrické vlastnosti, najčastejšie tetivové štvoruholníky.
\i Keď máme tip, overte, či v~nakreslenom obrázku platí aspoň približne.
\i Potom overte, či ho vieme dokázať, hoci aj pomocou tvrdenia, ktoré máme dokázať; byť si istý je dôležité.
\i Predpokladajme, že uhádnuté tvrdenie platí. Čo by sme v~našej situácii vedeli odvodiť? Možno by sme dospeli k~niečomu, čo vo všeobecnosti platiť nemôže? Možno sa úloha dá pomocou neho vyriešiť v~pár krokoch?
\i Používajte symetrické úvahy: ak toto platí, tak zo symetrie musí platiť aj niečo iné\dots
\i Ak hypotézu nevieme dokázať ani vyvrátiť, zapíšte ju do čistopisu slovne. Ak platí a~je podstatná, môže vyniesť body.
\enditems

\sec Známe podtvrdenia

V~dnešných úlohách budeme potrebovať niekoľko tvrdení, ktoré sú vo všeobecnosti veľmi užitočné a~oplatí sa ich vedieť priamo použiť v~zložitom obrázku. Tu je krátky zoznam niektorých z~tých najdôležitejších, ktoré by nám mohli pomôcť aj pri riešení úloh z~tohto materiálu:

\EnvId{4G1lSZGp6WOW_d5PrWMRp-zakladne-uhly}
\Theorem{Základné uhly}{Nech $ABC$ je ostrouhlý trojuholník. Potom:

\begitems
\i ak je $H$ ortocentrum, tak $|\angle BHC| = 180^\circ - \alpha$,
\i ak je $O$ stred opísanej kružnice, tak $|\angle BOC| = 2\alpha$,
\i ak je $I$ stred vpísanej kružnice, tak $|\angle BIC| = 90^\circ + \frac\alpha2$,
\i ak je $I_a$ stred kružnice pripísanej oproti vrcholu~$A$, tak $|\angle BI_aC|=90^\circ-\frac\alpha2$.
\enditems
}{}

\EnvId{bvirmYCiJc6UXSLVpp95M-svrckov-bod}
\Theorem{Švrčkov bod}{Nech $ABC$ je trojuholník so stredom vpísanej kružnice~$I$ a~$|AB|\neq|AC|$. Nech~$I_a$ je stred kružnice pripísanej oproti vrcholu~$A$. Potom sa os uhla~$\angle BAC$ a~os úsečky~$BC$ pretnú na kružnici opísanej trojuholníku~$ABC$ v~bode, ktorý je stredom úsečky~$II_a$ a~zároveň stredom kružnice opísanej štvoruholníku $BICI_a$.}{}

\EnvId{e3oWC42hI0-yUVRoB0_8O-anti-svrckov-bod}
\Theorem{Anti-Švrčkov bod}{Nech $ABC$ je trojuholník s~$|AB|\neq|AC|$. Nech $I_b$ a~$I_c$ sú postupne stredy kružníc pripísaných oproti vrcholom~$B$ a~$C$. Potom sa vonkajšia os uhla~$\angle BAC$ a~os úsečky~$BC$ pretnú na kružnici opísanej trojuholníku~$ABC$ v~bode~$M$, ktorý je stredom úsečky~$I_bI_c$ a~zároveň stredom kružnice opísanej štvoruholníku $BCI_bI_c$.}{}

\EnvId{hXmQS3PNqrvP6cXDIywLj-spiralna-podobnost-chodi-po-dvoch}
\Theorem{Špirálna podobnosť chodí po dvoch}{Ak sú trojuholníky~$OBC$ a~$OB_0C_0$ podobné a~rovnako orientované, tak sú podobné aj trojuholníky $OBB_0$ a~$OCC_0$.}{}

\EnvId{xHpsKi1FxuaOE4TsQDYXt-zrkadlenie-ortocentra}
\Theorem{Zrkadlenie ortocentra}{Nech $ABC$ je trojuholník s~ortocentrom~$H$. Potom:
\begitems
\i zrkadlový obraz~$H$ podľa priamky~$BC$ leží na kružnici opísanej trojuholníku~$ABC$,
\i zrkadlový obraz~$H$ podľa stredu úsečky~$BC$ leží na kružnici opísanej trojuholníku $ABC$ v~bode diametrálne protiľahlom k~$A$.
\enditems}{}

\sec Pridávanie bodov do obrázka

Konečne sme pri veci~\SlightSmile{} Pri riešení úlohy sa často stáva, že nevieme poriadne uchopiť nejakú hypotézu alebo tvrdenie, ktoré máme dokázať. Dôvodom môže byť to, že v~obrázku chýbajú všetky body potrebné na vyriešenie úlohy. Autori úloh mávajú vo zvyku kľúčové body skrývať, a~tým úlohy robiť zaujímavými a~často oveľa ťažšími. Vedieť pridať tie správne body je pravé umenie. Zopár všeobecných rád:

\begitems
\i Body pridávame, aby sme lepšie porozumeli nejakej hypotéze alebo tvrdeniu, ktoré máme dokázať.
\i Pridávame body, ktoré generujú pekné vlastnosti ako podobnosť, tetivové štvoruholníky, rovnobežnosť, \dots
\i Pridávame body preto, aby sme mohli odstrániť iné, ťažšie uchopiteľné body.
\enditems

Existuje niekoľko opakujúcich sa situácií, keď sa vždy oplatí uvažovať o~pridaní nového bodu, keď nič iné nefunguje. V~tomto materiáli precvičujeme najmä tieto:

\begitems \style n

\i Ak sú v~úlohe stredy úsečiek, oplatí sa pridať ďalšie stredy, aby vznikli stredné priečky.
\i Ak je v~úlohe stred úsečky, obraz v~stredovej súmernosti podľa tohto stredu vytvorí rovnobežník.
\i Ak je v~úlohe pravouhlý trojuholník, oplatí sa doplniť ho na rovnoramenný trojuholník.
\i Ak je v~úlohe ortocentrum, oplatí sa zamyslieť nad jeho obrazmi v~zrkadlení podľa strany alebo podľa stredu strany daného trojuholníka.
\i Ak sú v~úlohe rovnobežné úsečky, oplatí sa uvažovať o~rovnoľahlostiach, ktoré zobrazia jednu úsečku na druhú, a~zrkadliť pomocou nich ďalšie body.
\i Ak dve kružnice prechádzajú spoločným bodom, zaujímavým bodom môže byť ich druhý priesečník. Vytvorí aspoň chordálu.
\i Niekedy stačí pretnúť vhodné priamky, alebo vhodnú priamku a~kružnicu. Dôvodov na to môže byť viacero, napríklad sa môže objaviť rovnoramenný trojuholník alebo tetivový štvoruholník.

\enditems

Tento zoznam ani zďaleka nie je úplný, no na najbližších 24 úloh by mal stačiť. Vo všeobecnosti sú pri pridávaní bodov do obrázka možnosti takmer neobmedzené.

\sec Úlohy

Nasledujúce úlohy sú zhruba štyroch typov:

\begitems \style n
\i Kľúčovým krokom je pridanie správneho bodu; zvyšok je ľahký.
\i Pridanie je pomerne priamočiare, ale dotiahnutie úlohy vyžaduje trochu práce.
\i Pred pridaním bodu treba situáciu preskúmať, čo potom pridanie inšpiruje.
\i To isté ako vyššie, no ani po takom pridaní to nie je ľahké.
\enditems

\secc Štandardné úlohy

Nasledujúce úlohy sú zhruba usporiadané podľa obťažnosti a~používajú zhruba povedané \uv{štandardné} triky.

\EnvId{yt0KVCE0mByjG0W5AhbYA-patuholnik-s-rovnakymi-uhlami}
\Problem{0}{}{
    Päťuholník~$ABCDE$ má všetky vnútorné uhly rovnaké. Dokážte, že os úsečky~$AB$, os úsečky~$CD$ a~os uhla~$DEA$ sa pretínajú v~jednom bode.
}{
    Úloha o~piatich bodoch. Čo viac by si človek mohol priať? Paradoxne sa to nedá bez pridania bodov. Rovnaké uhly by mali niečo naznačovať.
}{
    Rovnaké uhly sa dajú využiť pretnutím vhodných dvojíc predĺžených strán, čím vzniknú rovnoramenné trojuholníky. Osi strán zrazu majú úplne iný význam.
}{
    Kľúčom je pretnúť predĺžené strany. Rovnaké vnútorné uhly dajú rovnoramenné trojuholníky, v~ktorých sa osi úsečiek $AB$ a~$CD$ zmenia na osi uhlov.

    Súčet vnútorných uhlov päťuholníka je $540^\circ$, takže každý z~nich má veľkosť $108^\circ$. Všetky sú menšie ako $180^\circ$, päťuholník $ABCDE$ je teda konvexný, a~keďže ani jeden z~nich nie je priamy, pre každú stranu ležia zvyšné tri vrcholy striktne v~tej istej otvorenej polrovine určenej priamkou tejto strany.

    Predĺžme stranu $EA$ za bod $A$ a~stranu $CB$ za bod $B$. Body $E$ a~$C$ ležia v~tej istej polrovine určenej priamkou $AB$, takže obe tieto polpriamky vedú do opačnej polroviny, pričom s~úsečkou $AB$ zvierajú uhly $180^\circ - 108^\circ = 72^\circ$. Ich súčet je menší ako $180^\circ$, a~tak sa polpriamky pretnú v~bode $P$. V~trojuholníku $ABP$ je $|\angle PAB| = |\angle PBA| = 72^\circ$, čiže $|PA| = |PB|$ a~os úsečky $AB$ je zároveň osou uhla $\angle APB$.

    Rovnako predĺžme stranu $BC$ za bod $C$ a~stranu $ED$ za bod $D$. Body $B$ a~$E$ ležia v~tej istej polrovine určenej priamkou $CD$, obe polpriamky teda vedú do opačnej a~s~úsečkou $CD$ zvierajú uhly $72^\circ$. Pretnú sa preto v~bode $Q$, pre ktorý je $|\angle QCD| = |\angle QDC| = 72^\circ$, takže $|QC| = |QD|$ a~os úsečky $CD$ je osou uhla $\angle CQD$.

    Bod $P$ leží na polpriamke opačnej k~polpriamke $BC$ a~bod $Q$ na polpriamke opačnej k~polpriamke $CB$, takže na priamke $BC$ nasledujú body v~poradí $P$, $B$, $C$, $Q$. Bod $E$ na priamke $BC$ neleží, čiže $EPQ$ je trojuholník. Ďalej $A$ leží medzi $P$ a~$E$ a~bod $D$ medzi $Q$ a~$E$.

    Z~tohto poradia je polpriamka $PA$ polpriamkou $PE$ a~polpriamka $PB$ polpriamkou $PQ$, takže os úsečky $AB$ je vnútornou osou uhla trojuholníka $EPQ$ pri vrchole $P$. Podobne sú polpriamky $QD$, $QC$ polpriamkami $QE$, $QP$, a~tak je os úsečky $CD$ vnútornou osou uhla pri vrchole $Q$. Napokon $EA$ je polpriamkou $EP$ a~$ED$ polpriamkou $EQ$, takže os uhla $\angle DEA$ je vnútornou osou uhla pri vrchole $E$. Všetky tri priamky sú teda vnútornými osami uhlov trojuholníka $EPQ$ a~pretínajú sa v~strede jeho vpísanej kružnice.
}

\EnvId{rIhQMrZeKZ2MH-PmkjT9D-bod-na-taznici-rovnoramenny-trojuholnik}
\Problem{0}{}{
    Vo vnútri trojuholníka~$ABC$, na jeho ťažnici~$AM$, sa nachádza bod~$K$ taký, že $|CK|=|AB|$. Nech~$L$ je priesečník priamok~$CK$ a~$AB$. Dokážte, že trojuholník~$AKL$ je rovnoramenný.
}{
    Bod~$K$ na ťažnici vyzerá zvláštne. V~každom prípade potrebujeme využiť fakt, že ide o~ťažnicu, čiže že~$M$ je stred, ideálne spôsobom, ktorý nám pomôže lepšie pochopiť rovnosť $|AB|=|CK|$.
}{
    Kľúčom je zrkadliť~$A$ podľa~$M$ do~$A'$, čo dá bodu~$M$ silný význam, keďže rovnosť~$|AB|=|CK|$ sa zrazu stane~$|A'C|=|CK|$.
}{
    Rovnosť $|CK| = |AB|$ spája dve úsečky, ktoré nemajú spoločný ani bod, ani smer, a~jediné, čo o~ťažnici vieme, je to, že $M$ je stred. Stred úsečky pritom volá po stredovej súmernosti podľa neho: nech $A'$ je obraz $A$ v~súmernosti so stredom $M$. Uhlopriečky $AA'$ a~$BC$ štvoruholníka $ABA'C$ sa navzájom polia, takže ide o~rovnobežník, a~preto
    $$
    |CA'| = |AB| = |CK|, \qquad CA' \parallel AB.
    $$
    Podmienka zo zadania sa tým presunula do jedného trojuholníka.

    Bod $K$ leží vnútri trojuholníka $ABC$, teda vnútri úsečky $AM$, a~keďže $M$ je stred $AA'$, leží aj vnútri úsečky $AA'$. Polpriamka $CK$ vychádza z~vrcholu $C$ a~prechádza vnútorným bodom trojuholníka, takže pretína protiľahlú stranu $AB$ v~jej vnútornom bode; ten je práve $L$ a~bod $K$ leží medzi $C$ a~$L$.

    Bod $C$ neleží na priamke $AM$, takže trojuholník $CKA'$ nie je degenerovaný; z~$|CA'| = |CK|$ je rovnoramenný so základňou $KA'$. Keďže $K$ leží vnútri $AA'$, je polpriamka $A'K$ totožná s~$A'A$, a~teda
    $$
    |\angle CKA'| = |\angle CA'K| = |\angle CA'A|.
    $$

    Priamka $AA'$ pretína úsečku $BC$ v~jej vnútornom bode $M$, takže oddeľuje $B$ od $C$. Uhly $\angle CA'A$ a~$\angle BAA'$ sú preto striedavými uhlami pri rovnobežkách $CA'$, $AB$ a~priečke $AA'$. Polpriamka $AB$ obsahuje $L$ a~polpriamka $AA'$ obsahuje $K$, čiže
    $$
    |\angle CA'A| = |\angle BAA'| = |\angle LAK|.
    $$

    Napokon uhly $\angle AKL$ a~$\angle A'KC$ sú vrcholové: polpriamky $KA$, $KA'$ sú opačné (bod $K$ leží vnútri $AA'$) a~opačné sú aj polpriamky $KL$, $KC$ (bod $K$ leží medzi $C$ a~$L$). Spojením všetkého dostávame
    $$
    |\angle AKL| = |\angle A'KC| = |\angle LAK|,
    $$
    takže trojuholník $AKL$ je rovnoramenný so základňou $AK$ a~platí $|LA| = |LK|$.
}

\EnvId{ull-eZWewTAOKLQHeu11Z-kolmica-na-mh-cez-ortocentrum}
\Problem{0}{}{
    Nech $ABC$ je ostrouhlý trojuholník s~ortocentrom~$H$ a~nech~$M$ je stred~$BC$. Body~$P$ a~$Q$ ležia postupne na stranách~$AB$ a~$AC$ tak, že $P$, $H$, $Q$ ležia na jednej priamke kolmej na~$MH$. Dokážte, že $|HP| = |HQ|$.
}{
    V~úlohe máme~$H$, stred~$BC$ a~nejakú úlohu hrá aj priamka~$MH$. Nuž, čo urobíme s~tým ortocentrom~\SlightSmile?
}{
    Ak~$H'$ je obraz~$H$ v~stredovej súmernosti podľa~$M$, tak vieme, že dopadne na kružnicu opísanú trojuholníku~$ABC$ do bodu protiľahlého k~$A$. Objavili sa nejaké tetivové štvoruholníky?
}{
    Označme $\alpha = |\angle BAC|$. V~úlohe vystupuje ortocentrum aj stred strany, tak ortocentrum podľa tohto stredu zrkadlíme: nech $H'$ je obraz bodu $H$ v~stredovej súmernosti podľa $M$. Podľa vety \uv{Zrkadlenie ortocentra} leží $H'$ na kružnici opísanej trojuholníku $ABC$, a~to v~bode diametrálne protiľahlom k~$A$.

    Trojuholník je ostrouhlý, takže $H$ leží v~jeho vnútri; neleží teda na priamke $BC$, čiže $H \neq M$ a~$H' \neq H$. Priamka $MH$ zo zadania je tak priamkou $HH'$.

    Úsečky $BC$ a~$HH'$ majú spoločný stred $M$, takže $BHCH'$ je rovnobežník. Podľa vety \uv{Základné uhly} je $|\angle BHC| = 180^\circ - \alpha$ a~susedné uhly rovnobežníka majú súčet $180^\circ$, takže
    $$
    |\angle HBH'| = |\angle HCH'| = \alpha. \eqno(*)
    $$

    Bod $H$ neleží na priamke $AB$, a~teda $P \neq H$. Ani $H'$ na priamke $AB$ neleží: tá pretína opísanú kružnicu iba v~bodoch $A$, $B$, pričom $H' \neq A$, a~rovnosť $H' = B$ by znamenala, že $AB$ je priemer, čiže $|\angle ACB| = 90^\circ$. Bod $P$ je teda rôzny od $H$ aj od $H'$ a~môžeme uvažovať kružnicu $\omega_1$ s~priemerom $PH'$.

    Priamka $PQ$ prechádza bodom $H$ a~je kolmá na $MH = HH'$, takže $|\angle PHH'| = 90^\circ$ a~bod $H$ leží na $\omega_1$. Leží na nej aj bod $B$: ak $P = B$, je to zrejmé, a~ak $P \neq B$, tak priamka $PB$ je priamka $AB$ a~z~Tálesovej vety nad priemerom $AH'$ je $|\angle PBH'| = |\angle ABH'| = 90^\circ$. Vznikol tak tetivový štvoruholník $PBHH'$.

    Uhly $\angle HPH'$ a~$\angle HBH'$ prislúchajú v~$\omega_1$ tej istej tetive $HH'$, takže sú zhodné alebo majú súčet $180^\circ$. Prvý z~nich je ostrý, lebo trojuholník $PHH'$ má pravý uhol pri $H$, a~druhý sa podľa $(*)$ rovná $\alpha < 90^\circ$; sú teda zhodné, čiže $|\angle HPH'| = \alpha$. (V~prípade $P = B$ ide rovno o~ten istý uhol.)

    Rovnaká úvaha platí pre $Q$. Bod $H'$ neleží na priamke $AC$, lebo $H' = C$ by znamenalo, že $AC$ je priemer, čiže $|\angle ABC| = 90^\circ$. Kružnica $\omega_2$ s~priemerom $QH'$ prechádza bodom $H$ vďaka $|\angle QHH'| = 90^\circ$ a~bodom $C$ vďaka $|\angle ACH'| = 90^\circ$, takže z~tetivy $HH'$ a~$(*)$ dostávame $|\angle HQH'| = |\angle HCH'| = \alpha$.

    Trojuholníky $PHH'$ a~$QHH'$ majú pravý uhol pri $H$ a~uhol $\alpha$ pri vrchole $P$, resp. $Q$, takže zo súčtu uhlov je $|\angle PH'H| = |\angle QH'H| = 90^\circ - \alpha$. Zhodujú sa teda v~strane $HH'$ a~v~oboch uhloch pri jej krajných bodoch, čiže sú podľa vety $usu$ zhodné. Odtiaľ $|HP| = |HQ|$.
}

\EnvId{n2g_94VghWgVCIT8C-mse-spojnica-stredov-stran-a-uhlopriecky}
\Problem{0}{}{
    V~štvoruholníku~$ABCD$ zviera priamka spájajúca stredy strán~$BC$ a~$AD$ rovnaké uhly s~obidvoma uhlopriečkami. Dokážte, že uhlopriečky sú rovnako dlhé.
}{
    Uhol medzi strednou priečkou a~uhlopriečkou štvoruholníka nie je niečo, čo by som vídal každý deň. Kľúčom bude zamerať sa na samotné stredy. Čo tak pridať jeden alebo viac ďalších stredov?
}{
    Dobrým stredom na pridanie je stred~$AB$ alebo~$CD$. Vďaka rovnobežnosti pochádzajúcej zo stredných priečok vieme podmienku zo zadania pekne preformulovať. Potom to už dorazíme.
}{
    Zadanie dáva dva stredy strán, a~to je presne situácia, keď sa oplatí pridať ďalší stred a~nechať vzniknúť stredné priečky. Označme $M$, $N$ stredy strán $BC$, $AD$ a~pridajme stred $P$ strany $AB$. Body $P$, $M$ sú stredy strán $AB$, $BC$ trojuholníka $ABC$ a~body $P$, $N$ sú stredy strán $AB$, $AD$ trojuholníka $ABD$, takže sa obe uhlopriečky naraz objavia ako stredné priečky vychádzajúce z~bodu $P$:
    $$
    \gather
    PM \parallel AC, \qquad |PM| = \tfrac12|AC|, \\
    PN \parallel BD, \qquad |PN| = \tfrac12|BD|.
    \endgather
    $$
    Dokazovaná rovnosť $|AC| = |BD|$ sa tým mení na $|PM| = |PN|$, čiže na to, že trojuholník $PMN$ je rovnoramenný, a~podmienka zo zadania hovorí presne o~uhloch, ktoré priamka $MN$ zviera s~jeho ramenami.

    Overme, že $PMN$ je naozaj trojuholník. Priamky $AC$, $BD$ nie sú rovnobežné: konvexný štvoruholník má uhlopriečky, ktoré sa pretínajú vnútri, a~v~nekonvexnom leží jeden vrchol vnútri trojuholníka určeného zvyšnými tromi, pričom priamka spájajúca tento vrchol s~protiľahlým vrcholom (teda jedna z~uhlopriečok) pretína protiľahlú stranu tohto trojuholníka, čiže druhú uhlopriečku. Keby bod $P$ ležal na priamke $MN$, splývali by priamky $PM$ i~$PN$ s~priamkou $MN$ a~vyšlo by $AC \parallel MN \parallel BD$, spor.

    Nech $\varphi \in [0^\circ, 90^\circ]$ je spoločná veľkosť uhla, ktorý priamka $MN$ zviera s~oboma uhlopriečkami; vďaka rovnobežnostiam vyššie ho zviera aj s~priamkami $PM$ a~$PN$. Vnútorný uhol $|\angle PMN|$ trojuholníka $PMN$ je preto $\varphi$ alebo $180^\circ - \varphi$ a~to isté platí pre $|\angle PNM|$. Súčet dvoch vnútorných uhlov trojuholníka je menší ako $180^\circ$, čo vylučuje dvojicu $\varphi$, $180^\circ - \varphi$ so súčtom rovno $180^\circ$, aj dvojicu $180^\circ - \varphi$, $180^\circ - \varphi$ so súčtom $360^\circ - 2\varphi \geq 180^\circ$. Zostáva jediná možnosť
    $$
    |\angle PMN| = |\angle PNM| = \varphi,
    $$
    takže trojuholník $PMN$ je rovnoramenný, $|PM| = |PN|$ a~odtiaľ $|AC| = 2|PM| = 2|PN| = |BD|$.
}

\EnvId{Z8MdAzTmbIxxpEyc-UnMy-ortocentra-tetivoveho-stvoruholnika}
\Problem{0}{Česko-slovenská olympiáda 2009}{
    Nech $ABCD$ je tetivový štvoruholník. Dokážte, že priamka spájajúca ortocentrum trojuholníka~$ABC$ s~ortocentrom trojuholníka~$ABD$ je rovnobežná s~priamkou~$CD$.
}{
    Ortocentrá sú komplikované body, ak nakreslíme všetky výšky. Takú kresbu okamžite zavrhnite. Ich obrazy vo vhodných zrkadleniach dávajú lepšie pochopenie situácie. Pomáha tiež preformulovať to, čo vlastne dokazujeme, aby sme po zrkadlení ortocentier vedeli tvrdenie stále jasne vysloviť.
}{
    Po prvé, stačí ukázať, že úsečka spájajúca ortocentrá má rovnakú dĺžku ako~$CD$ (potom máme rovnobežník). Toto uvedomenie nám dodá istotu, že keď tieto ortocentrá zrkadlíme podľa osi, stále máme dobré tvrdenie na dokázanie, totiž rovnosť dvoch tetív kružnice.
}{
    Označme $\omega$ kružnicu opísanú štvoruholníku $ABCD$, $O$ jej stred, $M$ stred strany $AB$ a~$H_1$, $H_2$ postupne ortocentrá trojuholníkov $ABC$ a~$ABD$. Body $A$, $B$, $C$, $D$ ležia na $\omega$, takže žiadne tri z~nich nie sú kolineárne a~oba trojuholníky sú nedegenerované. Podstatné je, že obom je opísaná tá istá kružnica $\omega$ a~že majú spoločnú stranu $AB$.

    Výšky nekreslíme, stačí nám z~definície ortocentra to, že $CH_1 \perp AB$ a~$DH_2 \perp AB$, čiže $CH_1 \parallel DH_2$. Tvrdenie by teda bolo hotové, keby bol štvoruholník $CH_1H_2D$ rovnobežník. Rovnobežnosť jednej dvojice protiľahlých strán už máme, takže ide len o~dĺžky: potrebujeme $|H_1H_2| = |CD|$ so súhlasnou orientáciou. V~tom je celý zisk z~preformulovania, lebo dĺžky sa zrkadlením nemenia. Ortocentrá teraz smieme vymeniť za ich obrazy a~dokazované tvrdenie sa aj potom bude dať jasne vysloviť.

    Zrkadlíme podľa stredu spoločnej strany, teda podľa bodu $M$; ten je pre oba trojuholníky ten istý, a~to je jediný dôvod, prečo sa výsledné obrazy dajú porovnať. Podľa vety o~zrkadlení ortocentra je obrazom $H_1$ v~stredovej súmernosti podľa $M$ bod $C_1$ kružnice $\omega$ diametrálne protiľahlý k~$C$; tá istá veta pre trojuholník $ABD$ dáva, že obrazom $H_2$ je bod $D_1 \in \omega$ diametrálne protiľahlý k~$D$. Z~$C \neq D$ je $C_1 \neq D_1$, a~teda $H_1 \neq H_2$, takže priamka spájajúca obe ortocentrá vôbec existuje.

    Stredová súmernosť podľa $M$ zobrazuje priamku $H_1H_2$ na priamku $C_1D_1$ a~stredová súmernosť podľa $O$ zobrazuje priamku $CD$ na tú istú priamku $C_1D_1$, keďže $C_1$, $D_1$ sú obrazy $C$, $D$. Zostávajúce tvrdenie je tak rovnosť dvoch tetív kružnice $\omega$ a~tá je po tomto pozorovaní zadarmo:
    $$
    |H_1H_2| = |C_1D_1| = |CD|.
    $$
    Stredová súmernosť navyše zobrazuje každú priamku na priamku s~ňou rovnobežnú, takže $H_1H_2 \parallel C_1D_1$ aj $CD \parallel C_1D_1$, a~preto $H_1H_2 \parallel CD$. (Pri niektorých polohách vrcholov obe priamky splynú, napríklad keď je $AB$ priemer $\omega$, lebo vtedy $H_1 = C$ a~$H_2 = D$; totožné priamky považujeme za rovnobežné.)
}

\EnvId{alkX6MavufXBEbJ3d6M4z-ortocentrum-v-rovnoramennom-lichobezniku}
\Problem{0}{}{
    Nech $ABCD$ je rovnoramenný lichobežník s~$AB \parallel CD$ vpísaný do kružnice so stredom~$O$. Nech~$H$ je ortocentrum trojuholníka $ABD$ a~nech~$P$ je priesečník~$DH$ a~$AB$. Dokážte, že $CH \parallel OP$.
}{
    Máme ortocentrum. Nemusí byť hneď jasné, čo s~ním, ale čo tak zrkadliť ho podľa niečoho?
}{
    Nech~$H'$ je obraz~$H$ v~zrkadlení podľa~$AB$. Dopadne na kružnicu opísanú štvoruholníku~$ABCD$. Teraz obrázok obsahuje všetko, čo potrebujeme, a~zostáva už len pospájať pár faktov.
}{
    V~úlohe vystupuje ortocentrum, tak sa pozrime na jeho obrazy v~zrkadlení podľa strany a~podľa stredu strany trojuholníka $ABD$. Označme $\omega$ kružnicu opísanú štvoruholníku $ABCD$ a~$H'$ obraz bodu $H$ v~osovej súmernosti podľa priamky $AB$. Ukážeme, že $H'$ je bod diametrálne protiľahlý k~$C$; zvyšok už bude stredná priečka.

    Priamka $DH$ je výška z~$D$ v~trojuholníku $ABD$, takže je kolmá na $AB$ a~$P$ je jej päta. Osová súmernosť podľa $AB$ zobrazuje priamku $DH$ samu na seba, lebo je na $AB$ kolmá; bod $H'$ preto leží na priamke $DH$ a~$P$ je stred úsečky $HH'$.

    Nech $M$ je stred úsečky $AB$ a~$s$ jej os. Bod $O$ leží na $s$, lebo je rovnako vzdialený od $A$ a~$B$. Priamka $s$ je kolmá na $AB$, teda aj na $CD$, a~kolmica na tetivu vedená stredom kružnice prechádza stredom tejto tetivy, takže $s$ je zároveň osou úsečky $CD$. Súmernosť podľa $s$ preto zobrazuje $\omega$ na seba, zachováva jej stred $O$, zamieňa $A$ s~$B$ a~$C$ s~$D$.

    Podľa druhej časti vety \uv{Zrkadlenie ortocentra}, použitej na trojuholník $ABD$ a~jeho stranu $AB$, je obrazom bodu $H$ v~stredovej súmernosti podľa $M$ bod $E$ diametrálne protiľahlý k~$D$. Stredová súmernosť podľa $M$ je pritom zložením osových súmerností podľa dvoch navzájom kolmých priamok prechádzajúcich bodom $M$, teda podľa $AB$ a~podľa $s$: ak bod $H$ zrkadlíme najprv podľa $AB$ a~potom podľa $s$, dostaneme $E$. Prvý krok dá $H'$, takže $H'$ je obrazom $E$ v~súmernosti podľa $s$. Tá zobrazuje $D$ na $C$ a~$O$ na $O$, čiže bod diametrálne protiľahlý k~$D$ zobrazuje na bod diametrálne protiľahlý k~$C$. Bod $H'$ je teda diametrálne protiľahlý k~$C$; špeciálne leží na $\omega$ a~$O$ je stred úsečky $CH'$.

    Bod $C$ neleží na priamke $DH$, lebo tá je kolmá na $CD$, a~teda pretína priamku $CD$ jedine v~$D$. Rôzne sú aj body $H$ a~$H'$: keby splynuli, ležal by $H$ na $AB$, čiže $H = P$, a~zároveň na $\omega$, čiže $P$ by bol jeden z~bodov $A$, $B$. Súmernosť podľa $s$ však zobrazuje pätu kolmice z~$D$ na $AB$ na pätu kolmice z~$C$ na $AB$, takže by týmito pätami boli práve body $A$ a~$B$. Z~rovnobežnosti $CD \parallel AB$ by potom vyšlo $|CD| = |AB|$ a~$ABCD$ by bol rovnobežník, nie lichobežník.

    Body $H$, $H'$, $C$ teda tvoria trojuholník, pričom $P$ je stred jeho strany $HH'$ a~$O$ stred strany $H'C$. Úsečka $OP$ je jeho strednou priečkou, a~preto $OP \parallel CH$.
}

\EnvId{WL5FFRrQH7pTblc3sQzBH-pravy-uhol-pri-strede-usecky}
\Problem{0}{Poľsko 2002}{
    Nech $ABCD$ je konvexný štvoruholník s~$|AB|<|BC|$ a~$|AD|<|CD|$. Body $P$ a~$Q$ ležia postupne na~$BC$ a~$CD$ tak, že $|BP|=|BA|$ a~$|DQ|=|DA|$. Nech~$M$ je stred~$PQ$. Dokážte, že ak $|\angle BMD|=90^\circ$, tak $ABCD$ je tetivový štvoruholník.
}{
    Potrebujeme využiť zvláštnu podmienku $|\angle BMD|=90^\circ$. Bod~$M$ je stred úsečky, čo už niečo naznačuje, a~pravý uhol tento náznak len posilňuje. V~každom prípade sa oplatí použiť farebné ceruzky na vyznačenie rovnako dlhých úsečiek.
}{
    Kľúčom je zrkadliť napríklad~$D$ podľa~$M$. Po vyznačení všetkých rovnakých úsečiek vznikajúcich z~výsledného rovnobežníka a~rovnoramenného trojuholníka by malo niečo padnúť do oka.
}{
    Bod $M$ je stredom úsečky, takže sa žiada zrkadliť podľa neho a~vyrobiť rovnobežník; predpoklad $|\angle BMD|=90^\circ$ napovedá, ktorý bod si to zaslúži. Nech teda $D'$ je obraz bodu $D$ v~stredovej súmernosti podľa $M$ (uhol $\angle BMD$ je definovaný, takže $M \neq D$, a~teda aj $D' \neq D$).

    Zo zadania je $|BP|=|BA|<|BC|$ a~$|DQ|=|DA|<|DC|$, takže $P$ je vnútorným bodom strany $BC$ a~$Q$ vnútorným bodom strany $CD$. Štvoruholník $ABCD$ je konvexný, takže žiadne tri jeho vrcholy nie sú kolineárne; priamky $BC$ a~$CD$ sa preto pretínajú jedine v~$C$ a~$P \neq Q$.

    Úsečky $PQ$ a~$DD'$ majú spoločný stred $M$, takže $DPD'Q$ je rovnobežník, a~teda
    $$
    |PD'| = |QD| = |DA|, \qquad \overline{PD'} = \overline{DQ}.
    $$
    Priamka $BM$ je kolmá na $MD$, čiže na $DD'$, a~prechádza stredom $M$ tejto úsečky. Je to teda os úsečky $DD'$, trojuholník $BDD'$ je rovnoramenný a~
    $$
    |BD'| = |BD|.
    $$

    Trojuholníky $BPD'$ a~$BAD$ sa tak zhodujú vo všetkých troch stranách: $|BP|=|BA|$ zo zadania, $|PD'|=|AD|$ z~rovnobežníka a~$|BD'|=|BD|$ z~rovnoramennosti. Podľa $sss$ sú zhodné, pričom stranám $BP$, $PD'$ zodpovedajú strany $BA$, $AD$, takže uhly pri vrcholoch $P$ a~$A$ sa rovnajú:
    $$
    |\angle BPD'| = |\angle BAD|.
    $$

    Zostáva prečítať uhol $\angle BPD'$ v~pôvodnom obrázku. Bod $P$ leží vnútri úsečky $BC$, takže $\overline{PB}$ je kladným násobkom $\overline{CB}$; bod $Q$ leží vnútri úsečky $DC$, takže $\overline{PD'} = \overline{DQ}$ je kladným násobkom $\overline{DC}$, čiže záporným násobkom $\overline{CD}$. Ramená uhlov $\angle BPD'$ a~$\angle BCD$ sú teda rovnobežné, jedna dvojica rovnako a~druhá opačne orientovaná, odkiaľ
    $$
    |\angle BAD| + |\angle BCD| = |\angle BPD'| + |\angle BCD| = 180^\circ.
    $$
    Konvexný štvoruholník, v~ktorom je súčet protiľahlých uhlov $180^\circ$, je tetivový, takže $ABCD$ je tetivový štvoruholník.
}

\EnvId{Tax9-rYDqJvAbCiA0VqMC-paty-kolmic-a-stred-strany}
\Problem{0}{USA TST 2000}{
    Nech $ABCD$ je tetivový štvoruholník, ktorého uhlopriečky sa pretínajú v~$P$. Nech~$K$ a~$L$ sú päty kolmíc z~$P$ na~$AB$ a~$CD$ a~nech~$M$ je stred strany~$AD$. Dokážte, že $|MK|=|ML|$.
}{
    Oplatí sa nakresliť úsečku~$AD$ vodorovne (netradične), čo dá najlepší obrázok (prekreslite hneď teraz, ak ste tak neurobili). Ani v~novom obrázku nemusí byť jasné, čo sa má stať. Stred~$AD$ nevyzerá jednoducho uchopiteľne. Možno potrebuje na spoluprácu nejaké ďalšie stredy. Ideálne také, ktoré súvisia aj s~bodmi~$K$ a~$L$.
}{
    Kľúčom sú zvyšné stredy strán trojuholníka~$ADP$. Potom sa udeje mágia a~od riešenia nás delí len pár ľahkých úvah.
}{
    Bod $M$ je stred strany $AD$ a~sám osebe nedáva ani pekný uhol, ani použiteľnú dĺžku. Keď v~úlohe takto osamotene visí stred úsečky, pridávame ďalšie stredy: doplňme $M$ na trojicu stredov strán trojuholníka $ADP$, teda pridajme stred $X$ úsečky $AP$ a~stred $Y$ úsečky $DP$. Odmena príde okamžite. Bod $K$ leží na priamke $AB$ a~$PK \perp AB$, takže $|\angle AKP| = 90^\circ$, prípadne $K = A$; tak či onak $K$ leží na Tálesovej kružnici nad priemerom $AP$, ktorej stredom je práve $X$. Rovnako $L$ leží na Tálesovej kružnici nad priemerom $DP$ so stredom $Y$. Jedno pridanie tak zviazalo $K$, $L$ aj $M$.

    Žiadne tri z~bodov $A$, $B$, $C$, $D$ nie sú kolineárne, lebo ležia na kružnici. Bod $P$ leží na priamkach $AC$ aj $BD$, takže $P \neq A$, $P \neq D$ a~$P$ neleží na priamke $AB$ ani na priamke $CD$ (inak by priamka $AB$ splynula s~$AC$, resp. $CD$ s~$BD$). Špeciálne $K \neq P$ a~$L \neq P$.

    Z~Tálesových kružníc je
    $$
    |XK| = |XP| = \tfrac12|AP|, \qquad |YL| = |YP| = \tfrac12|DP|.
    $$
    Rovnoľahlosť so stredom $A$ a~koeficientom $\tfrac12$ zobrazí $P$ na $X$ a~$D$ na $M$, takže $|MX| = \tfrac12|DP|$ a~polpriamka $XM$ smeruje rovnako ako $PY$. Rovnoľahlosť so stredom $D$ a~koeficientom $\tfrac12$ zobrazí $P$ na $Y$ a~$A$ na $M$, takže $|MY| = \tfrac12|AP|$ a~polpriamka $YM$ smeruje rovnako ako $PX$. Dĺžky si teda odpovedajú krížom:
    $$
    |MX| = |YL| = \tfrac12|DP|, \qquad |XK| = |MY| = \tfrac12|AP|.
    $$
    Trojuholníky $MXK$ a~$LYM$ sa tak zhodujú v~dvoch dvojiciach strán, a~ak sa zhodnú aj v~uhloch, ktoré tieto strany zvierajú, budú podľa vety $sus$ zhodné a~vyjde $|MK| = |LM|$. Úloha sa tým zredukovala na rovnosť $|\angle MXK| = |\angle MYL|$.

    Tu sa už orientovaným uhlom nevyhneme: päty $K$, $L$ môžu podľa konfigurácie padnúť aj mimo úsečiek $AB$, $CD$ a~polpriamka $XP$ nemusí ležať vnútri uhla $MXK$, kým znamienka ustrážia všetky polohy naraz. Pri pevne zvolenej orientácii roviny označme pre polpriamky $p$, $q$ so spoločným začiatkom $\angle(p \to q)$ uhol otočenia, ktoré prevedie $p$ na $q$, braný modulo $360^\circ$; pre priamky $p$, $q$ nech $\angle(p,q)$ je orientovaný uhol modulo $180^\circ$, takže jeho dvojnásobok je dobre určený modulo $360^\circ$.

    Položme $\omega = \angle(XP \to XK)$ a~$\omega' = \angle(YP \to YL)$; sú to uhly otočení so stredmi $X$, $Y$, ktoré zobrazia $P$ na $K$, resp. na $L$. Kľúčové je, že $\omega' = -\omega$. Body $A$, $K$, $P$ ležia na kružnici so stredom $X$, takže stredový uhol je dvojnásobkom obvodového a~platí
    $$
    \omega = 2\angle(AP,AK) = 2\angle(AC,AB),
    $$
    lebo $P$ leží na priamke $AC$ a~$K$ na priamke $AB$. (Ak $K = A$, čo nastane práve pri $AB \perp AC$, je $X$ stred úsečky $KP$, čiže $\omega = 180^\circ$, a~rovnosť platí tiež.) Rovnako $\omega' = 2\angle(DB,DC)$. Zo štyroch bodov na jednej kružnici je $\angle(AB,AC) = \angle(DB,DC)$, a~preto
    $$
    \omega = -2\angle(AB,AC) = -2\angle(DB,DC) = -\omega'.
    $$

    Ostáva pospájať. Polpriamka $XM$ smeruje rovnako ako $PY$, čiže opačne než $YP$, a~polpriamka $XP$ smeruje rovnako ako $MY$, čiže opačne než $YM$; otočenie oboch ramien o~$180^\circ$ uhol nemení, takže
    $$
    \angle(XM \to XP) = \angle(YP \to YM) = -\angle(YM \to YP).
    $$
    Spolu s~$\omega' = -\omega$ dostávame
    $$
    \eqalign{
        \angle(XM \to XK) &= \angle(XM \to XP) + \omega \cr
        &= -\angle(YM \to YP) - \omega' = -\angle(YM \to YL). \cr
    }
    $$
    Orientované uhly $\angle(XM \to XK)$ a~$\angle(YM \to YL)$ sú teda opačné, čiže neorientované uhly $|\angle MXK|$ a~$|\angle MYL|$ sú rovnaké. Podľa vety $sus$ sú teda trojuholníky $MXK$ a~$LYM$ zhodné, a~preto $|MK| = |ML|$.
}

\EnvId{0TlY8HJdJ_fVKGlkNgzpO-bod-vnutri-rovnobeznika}
\Problem{0}{}{
    V~rovnobežníku $ABCD$ leží bod~$P$, pre ktorý platí $|\angle APB| + |\angle CPD| = 180^\circ$. Dokážte, že $|\angle PBC|=|\angle PDC|$.
}{
    Skúste posunúť bod~$P$ tak, aby podmienka na súčet uhlov dávala zmysel.
}{
    Kľúčovým bodom je obraz bodu~$P$ v~posunutí v~smere~$BC$ (teda~$AD$); označme ho~$P'$. Potom máme tetivový štvoruholník $CP'DP$, ako aj množstvo rovnobežníkov.
}{
    Podmienka $|\angle APB| + |\angle CPD| = 180^\circ$ viaže dva uhly pri tom istom vrchole $P$, a~tak sa s~ňou nedá nič robiť. Rovnobežník má však dve dvojice rovnobežných zhodných strán, čo volá po posunutí zobrazujúcom jednu dvojicu na druhú. Ak ním posunieme aj bod $P$, jeden z~tých dvoch uhlov sa presunie k~novému vrcholu a~zo súčtu $180^\circ$ sa stane podmienka tetivovosti.

    Nech teda $\tau$ je posunutie zobrazujúce $B$ na $C$, teda posunutie o~vektor $BC = AD$, a~nech $P' = \tau(P)$. Keďže $\tau(A) = D$ a~$\tau(B) = C$, prejde trojuholník $ABP$ na trojuholník $DCP'$, takže
    $$
    |\angle DP'C| = |\angle APB| = 180^\circ - |\angle CPD|.
    $$

    Najprv určme, kde bod $P'$ leží. Bod $P$ je vnútorným bodom rovnobežníka, takže rovnobežka s~$BC$ vedená bodom $P$ pretína stranu $AB$ vo vnútornom bode $X$ a~stranu $DC$ vo vnútornom bode $Y$, pričom $P$ leží vnútri úsečky $XY$. Bod $\tau(X)$ leží na priamke $\tau(AB) = DC$ aj na priamke $XY$ (posunutie je v~smere $XY$), a~tie sa pretínajú jedine v~$Y$, čiže $\tau(X) = Y$. Preto $|XP'| = |XP| + |XY| > |XY|$, bod $P'$ leží na polpriamke $XY$ za bodom $Y$, a~body $P$, $P'$ tak ležia v~opačných polrovinách určených priamkou $CD$.

    Body $C$, $D$, $P$ neležia na jednej priamke; označme $\omega$ kružnicu opísanú trojuholníku $CDP$. Množina bodov polroviny neobsahujúcej $P$, z~ktorých vidno úsečku $CD$ pod uhlom $180^\circ - |\angle CPD|$, je práve oblúk kružnice $\omega$ v~tejto polrovine. Bod $P'$ do nej podľa predošlých dvoch odsekov patrí, takže leží na $\omega$ a~$CP'DP$ je tetivový štvoruholník. Tetiva $CD$ oddeľuje na $\omega$ body $P'$ a~$P$, čiže v~cyklickom poradí idú body $C$, $P'$, $D$, $P$; body $P'$ a~$D$ preto ležia na tom istom oblúku určenom tetivou $CP$.

    Zostáva preniesť uhol $\angle PBC$ ku kružnici $\omega$. Posunutie $\tau$ zobrazuje $B$ na $C$ a~$P$ na $P'$, takže úsečky $BP$ a~$CP'$ sú rovnobežné, rovnako dlhé a~rovnako orientované; keďže $P$ neleží na priamke $BC$, je $BPP'C$ nedegenerovaný rovnobežník. Jeho protiľahlé uhly pri vrcholoch $B$ a~$P'$ sa rovnajú:
    $$
    |\angle PBC| = |\angle PP'C|.
    $$
    Uhly $\angle PP'C$ a~$\angle PDC$ sú obvodové uhly kružnice $\omega$ nad tetivou $CP$ z~toho istého oblúka, a~teda zhodné. Spolu
    $$
    |\angle PBC| = |\angle PP'C| = |\angle PDC|.
    $$
}

\EnvId{RV54g2hn6AMdIasT4IODZ-konvexny-sestuholnik-so-zhodnymi-stranami}
\Problem{0}{Česko-slovenská olympiáda 2026}{
    Je daný konvexný šesťuholník $ABCDEF$ taký, že $|AB| = |BC|$, $|CD| = |DE| \neq |AD|$, $|EF| = |FA|$, $|\angle BDC| + |\angle EDF| = |\angle FDB|$. Dokážte, že $|\angle CBA| + |\angle EDC| + |\angle AFE| = 360^\circ$.
}{
    Podmienka na súčet uhlov je kľúčová. Potrebujeme zaviesť~bod, ktorý využije podmienku na súčet uhlov a~urobí ju o~niečo zmysluplnejšou.
}{
    Kľúčovým bodom je zrkadlový obraz $C$ podľa priamky $BD$, alebo analogicky zrkadlový obraz $E$ podľa priamky $DF$. Je to ten istý bod. A~je to všetko, čo potrebujeme.  Nezabudnite využiť nerovnostnú časť podmienky $|CD|=|DE| \neq |AD|$. Bez nej úloha neplatí.
}{
    Podmienka $|\angle BDC| + |\angle EDF| = |\angle FDB|$ hovorí, že uhly $BDC$ a~$EDF$ presne vyplnia uhol $BDF$. To nabáda preklopiť $C$ podľa priamky $BD$ a~$E$ podľa priamky $DF$: oba obrazy padnú dovnútra uhla $BDF$ a~prvým krokom bude ukázať, že padnú na ten istý bod.

    Najprv si ujasnime konfiguráciu. Uhlopriečky z~vrcholu $D$ delia vnútorný uhol konvexného šesťuholníka pri $D$ v~tom poradí, v~akom idú vrcholy po obvode, takže polpriamky $DC$, $DB$, $DA$, $DF$, $DE$ idú v~tomto uhlovom poradí a~platí
    $$
    |\angle CDB| + |\angle BDF| + |\angle FDE| = |\angle CDE| < 180^\circ.
    $$
    Priamka $BD$ preto oddeľuje $C$ od bodov $A$, $F$, $E$ a~priamka $DF$ oddeľuje $E$ od bodov $C$, $B$, $A$. Keďže oba sčítance v~podmienke zo zadania sú kladné, je tiež $|\angle BDC| < |\angle FDB|$ a~$|\angle FDE| < |\angle FDB|$.

    Označme $P$ obraz bodu $C$ v~osovej súmernosti podľa priamky $BD$ a~$Q$ obraz bodu $E$ v~osovej súmernosti podľa priamky $DF$. Ukážeme, že $P = Q$. Bod $P$ leží na opačnej strane priamky $BD$ ako $C$, teda na tej istej ako $F$, a~$|\angle BDP| = |\angle BDC| < |\angle FDB|$, takže polpriamka $DP$ leží vnútri uhla $BDF$. Symetricky $|\angle FDQ| = |\angle FDE| < |\angle FDB|$ a~vnútri uhla $BDF$ leží aj polpriamka $DQ$. Zo zadania je $|\angle BDP| = |\angle FDB| - |\angle EDF|$, a~teda
    $$
    |\angle FDP| = |\angle FDB| - |\angle BDP| = |\angle EDF| = |\angle FDQ|.
    $$
    Polpriamky $DP$ a~$DQ$ zvierajú s~priamkou $DF$ rovnaký uhol a~ležia pri nej na tej istej strane (obe vnútri uhla $BDF$), preto splývajú. Spolu s~$|DP| = |DC| = |DE| = |DQ|$ dostávame $P = Q$.

    Jediný bod $P$ tak spĺňa naraz všetky tri dĺžkové podmienky zo zadania:
    $$
    \gather
    |BP| = |BC| = |BA|, \\
    |DP| = |DC| = |DE|, \\
    |FP| = |FE| = |FA|.
    \endgather
    $$

    Prvá a~tretia z~nich hovoria, že $P$ leží na kružnici $k_B$ so stredom $B$ a~polomerom $|BA|$ aj na kružnici $k_F$ so stredom $F$ a~polomerom $|FA|$. Tie majú rôzne stredy, takže majú najviac dva spoločné body. Osová súmernosť podľa priamky $BF$ zobrazí každú z~nich samu na seba, a~keďže vnútorný uhol šesťuholníka pri vrchole $A$ je menší ako $180^\circ$, bod $A$ na priamke $BF$ neleží a~jeho obraz $A'$ je od neho rôzny. Spoločnými bodmi kružníc $k_B$, $k_F$ sú teda práve $A$ a~$A'$ a~bod $P$ je jedným z~nich. Možnosť $P = A$ však odpadá: dávala by $|AD| = |PD| = |CD|$, čo je v~spore s~predpokladom $|CD| \neq |AD|$. Preto $P = A'$. Toto je jediné miesto, kde sa nerovnostná časť podmienky použije, a~je nenahraditeľné: keby $P$ splynulo s~$A$, ležalo by mimo trojuholníka $BDF$ a~záverečný súčet uhlov by sa rozpadol.

    V~konvexnom šesťuholníku odrezáva uhlopriečka $BF$ trojuholník $ABF$, takže priamka $BF$ oddeľuje $A$ od $D$. Bod $P = A'$ preto leží v~otvorenej polrovine s~hraničnou priamkou $BF$ obsahujúcej $D$. Zároveň polpriamka $DP$ leží vnútri uhla $BDF$ a~$P \neq D$. Prienikom vnútra uhla $BDF$ s~touto polrovinou je práve vnútro trojuholníka $BDF$, takže $P$ je vnútorným bodom trojuholníka $BDF$.

    Osová súmernosť podľa $BD$ zobrazí trojuholník $BCD$ na $BPD$, súmernosť podľa $DF$ zobrazí $DEF$ na $DPF$ a~súmernosť podľa $BF$ zobrazí $BAF$ na $BPF$, čiže
    $$
    \gather
    |\angle BPD| = |\angle BCD|, \\
    |\angle DPF| = |\angle DEF|, \\
    |\angle FPB| = |\angle FAB|.
    \endgather
    $$
    Bod $P$ leží vo vnútri trojuholníka $BDF$, takže uhly pri ňom dávajú plný uhol, a~teda
    $$
    |\angle BCD| + |\angle DEF| + |\angle FAB| = 360^\circ.
    $$
    Súčet vnútorných uhlov šesťuholníka je $720^\circ$, odkiaľ
    $$
    |\angle CBA| + |\angle EDC| + |\angle AFE| = 720^\circ - 360^\circ = 360^\circ.
    $$
}

\EnvId{L-IkptXbj4-UJ_vpMWLKS-bod-s-piatimi-zhodnymi-uhlami}
\Problem{0}{Česko-slovenská olympiáda 2024}{
    Bod $P$ ležiaci vo vnútri konvexného štvoruholníka $ABCD$ spĺňa rovnosti
    $$|\angle PAD| = |\angle ADP| = |\angle CBP| = |\angle PCB| = |\angle CPD|.$$
    Nech $O$ je stred kružnice opísanej trojuholníku $CPD$. Dokážte, že $|OA| = |OB|$.
}{
    Uhly implikujú rovnobežné priamky. Naznačujú správny ťah vpred.
}{
    Bodom na zavedenie je priesečník $X$ priamok $BC$ a $AD$, keďže $PCXD$ je teraz rovnobežník. Úlohu teraz vieme vyriešiť so správnym pohľadom na to, čo máme.
}{
    Označme $\varphi$ spoločnú veľkosť piatich uhlov zo zadania. Trojuholník $APD$ (bod $P$ leží vo vnútri štvoruholníka, na priamke $AD$ teda neleží) má pri vrcholoch $A$ a~$D$ uhly veľkosti $\varphi$, je preto rovnoramenný s~$|PA| = |PD|$ a~navyše $\varphi < 90^\circ$; rovnako je $|PB| = |PC|$.

    Zvyšné dve rovnosti hovoria o~uhloch pri priečke. Keďže $P$ leží vo vnútri konvexného štvoruholníka $ABCD$, idú polpriamky $PA$, $PB$, $PC$, $PD$ okolo bodu $P$ v~tomto cyklickom poradí a~súčet uhlov $\angle APB$, $\angle BPC$, $\angle CPD$, $\angle DPA$ je $360^\circ$. Ak teda uhly meriame od polpriamky $PD$ v~smere $D \to A \to B \to C$, leží polpriamka $PA$ pod uhlom $|\angle DPA| < 180^\circ$, kým polpriamka $PC$ pod uhlom $360^\circ - |\angle CPD| > 180^\circ$, čiže priamka $PD$ oddeľuje $A$ od $C$. Uhly $\angle ADP$ a~$\angle DPC$ sú tak striedavé uhly priamok $AD$ a~$PC$ pri priečke $DP$, a~keďže sú zhodné, je $AD \parallel PC$. Analogicky priamka $PC$ oddeľuje $B$ od $D$ a~zo zhodnosti uhlov $\angle PCB$ a~$\angle CPD$ dostaneme $BC \parallel PD$.

    Dve dvojice rovnobežiek si pýtajú priesečník. Priamky $AD$ a~$BC$ rovnobežné nie sú (inak by boli rovnobežné aj $PC$ a~$PD$, čiže by $C$, $P$, $D$ ležali na jednej priamke), takže sa pretínajú v~bode $X$. Z~$PC \parallel XD$ a~$CX \parallel PD$ je $PCXD$ rovnobežník. Dokážeme dokonca viac, než sa žiada:
    $$
    |OA| = |OX| = |OB|.
    $$

    Kľúčové je pozorovanie, že \textbf{os úsečky $PC$ je zároveň osou úsečky $AX$}.

    Nech $M$ je stred úsečky $PC$, $N$ stred úsečky $AD$ a~$Y$ stred úsečky $AX$. Posunutie zobrazujúce $P$ na $C$ zobrazuje vďaka rovnobežníku $PCXD$ bod $D$ na $X$, takže úsečka $AX$ vznikne z~úsečky $AD$ posunutím jedného konca; jej stred sa pritom posunie o~polovicu tohto posunutia, čiže posunutie zobrazujúce $P$ na $M$ zobrazuje $N$ na $Y$. Úsečka $MY$ je teda obrazom úsečky $PN$ a~má rovnaký smer. Trojuholník $APD$ je rovnoramenný, takže jeho ťažnica $PN$ je zároveň výškou, čiže $PN \perp AD$, a~preto aj $MY \perp AD$.

    Priamka $MY$ tak prechádza stredom $M$ úsečky $PC$ kolmo na $PC$ (lebo $PC \parallel AD$): je to os úsečky $PC$. Zároveň je kolmá na priamku $AD$ v~bode $Y$, čo je stred úsečky $AX$, ktorej oba krajné body na $AD$ ležia: je to teda aj os úsečky $AX$. (Body $A$ a~$X$ sú rôzne: $|\angle ADP| = \varphi$, kým $|\angle XDP| = 180^\circ - \varphi$ ako uhol rovnobežníka susedný k~$\angle DPC$, a~$\varphi \neq 90^\circ$.)

    Bod $O$ je stred kružnice opísanej trojuholníku $CPD$, leží teda na osi úsečky $PC$, a~preto $|OA| = |OX|$. Zámena $A \leftrightarrow B$, $D \leftrightarrow C$ zachováva všetky predpoklady (štvoruholník $BADC$ je ten istý konvexný štvoruholník prebehnutý opačným smerom a~pätica uhlov zo zadania prejde sama na seba) aj body $P$, $X$, $O$, takže rovnakou úvahou je os úsečky $PD$ osou úsečky $BX$ a~$|OB| = |OX|$.
}

\EnvId{aI08awqy1yrLlCuEOFHJ7-pata-vysky-a-stred-oh}
\Problem{0}{Česko-slovenská olympiáda 2011}{
    V~trojuholníku $ABC$, ktorý je ostrouhlý a~nie je rovnostranný, nech~$P$ je päta výšky z~$A$ na stranu~$BC$, $H$ ortocentrum, $O$ stred opísanej kružnice, $D$ priesečník polpriamky~$AO$ so stranou~$BC$ a~$E$ stred úsečky~$AD$. Dokážte, že priamka~$EP$ prechádza stredom úsečky~$OH$.
}{
    Situácia je pomerne zamotaná a~na prvý pohľad sa nedá objaviť veľa zaujímavého. Je v~nej však bod~$H$, s~ktorým vieme robiť rôzne zaujímavé veci. Ktorá z~nich sa sem hodí?
}{
    Zrkadlite~$H$ podľa~$BC$, čím získate~$H'$ ležiaci na kružnici opísanej trojuholníku~$ABC$. To znamená $|OA| = |OH'|$. Aj stred úsečky~$OH$ sa stane uchopiteľnejším.
}{
    Označme $\omega$ kružnicu opísanú trojuholníku $ABC$, $R$ jej polomer a~$N$ stred úsečky $OH$; trojuholník nie je rovnostranný, takže $O \neq H$.

    Ak $|AB| = |AC|$, ležia na osi strany $BC$ okrem $O$ aj body $A$, $H$ a~$P$, takže $D = P$ a~priamka $EP$ je práve táto os. Tá obsahuje $O$ aj $H$, a~teda aj stred úsečky $OH$. Ďalej preto predpokladajme $|AB| \neq |AC|$; vtedy priamka $AO$ nie je kolmá na $BC$.

    Jediným bodom konfigurácie, s~ktorým sa dá voľne narábať, je ortocentrum, a~jeho obraz $H'$ v~osovej súmernosti podľa $BC$ urobí naraz dve veci: podľa tvrdenia o~zrkadlení ortocentra leží $H'$ na $\omega$, čiže
    $$
    |OH'| = |OA| = R,
    $$
    a~zároveň je $P$ stredom úsečky $HH'$, čím sa dostane do hry aj stred úsečky $OH$. Priamka $AP$ je totiž kolmá na $BC$, v~tejto súmernosti sa teda zobrazí sama na seba a~bod $P$ zostane pevný; bod $H'$ preto leží na priamke $AP$ a~$P$ je stredom $HH'$.

    Trojuholník je ostrouhlý, takže $H$ leží v~jeho vnútri, a~teda vnútri úsečky $AP$. Bod $H'$ potom leží na polpriamke $AP$ za bodom $P$ a~body idú na výške z~$A$ v~poradí $A$, $H$, $P$, $H'$. Špeciálne $P$ je vnútorným bodom úsečky $AH'$ a~$H'$ leží v~opačnej polrovine s~hraničnou priamkou $BC$ ako $A$, takže $H' \neq A$.

    Nech $A'$ je bod diametrálne protiľahlý k~$A$. Keby bolo $H' = A'$, priamka $AA'$ by splývala s~výškou z~$A$ a~bola by kolmá na $BC$, čo sme vylúčili. Je teda $H' \neq A'$ a~z~Tálesovej vety nad priemerom $AA'$ dostávame $|\angle AH'A'| = 90^\circ$. Priamka $H'A'$ je preto kolmá na $AH'$, a~keďže $AH' \perp BC$, je
    $$
    H'A' \parallel BC.
    $$

    Uvažujme rovnoľahlosť $h$ so stredom $A$, ktorá zobrazuje $P$ na $H'$; jej koeficient $|AH'| \, / \, |AP|$ je kladný, lebo $P$ leží vnútri úsečky $AH'$. Priamku $BC$ zobrazí $h$ na rovnobežku s~$BC$ prechádzajúcu bodom $H'$, čo je práve priamka $H'A'$, a~priamku $AD$ zobrazí samu na seba. Priamka $AD$ nie je rovnobežná s~$BC$ (pretína ju v~bode $D$), takže s~priamkou $H'A'$ má jediný spoločný bod, a~tým je $A'$. Preto $h(D) = A'$, a~keďže rovnoľahlosť zachováva stredy úsečiek, zobrazí sa stred $E$ úsečky $AD$ na stred úsečky $AA'$, čiže na $O$:
    $$
    h(P) = H', \qquad h(D) = A', \qquad h(E) = O.
    $$

    Bod $E$ je vnútorným bodom úsečky $AD$, teda neleží na $BC$ a~$E \neq P$. Obrazom priamky $EP$ v~$h$ je priamka $OH'$, a~rovnoľahlosť zobrazuje priamku na priamku s~ňou rovnobežnú (prípadne na ňu samu). Priamka $EP$ je teda rovnobežná s~$OH'$ alebo s~ňou totožná.

    S~bodom $N$ naložíme rovnako. Rovnoľahlosť so stredom $H$ a~koeficientom $\frac12$ zobrazí $H'$ na stred úsečky $HH'$, čiže na $P$, a~bod $O$ na stred úsečky $HO$, čiže na $N$. Z~$|OH'| = R > 0$ je $O \neq H'$, a~teda aj $P \neq N$; priamka $PN$ je obrazom priamky $OH'$, takže je s~ňou rovnobežná alebo s~ňou totožná.

    Priamky $EP$ a~$PN$ obe prechádzajú bodom $P$ a~obe sú rovnobežné s~priamkou $OH'$ alebo s~ňou totožné. Taká priamka je bodom $P$ určená jednoznačne, takže $EP$ a~$PN$ splývajú a~bod $N$, čiže stred úsečky $OH$, leží na priamke $EP$.
}

\EnvId{M3FyEsmtSGx30V-INIVXw-dotycnica-a-stred-usecky-ap}
\Problem{0}{}{
    Nech $ABC$ je ostrouhlý trojuholník vpísaný do kružnice~$k$. Dotyčnica ku~$k$ v~bode~$A$ pretína priamku~$BC$ v~$P$. Nech~$M$ je stred~$AP$ a~nech~$Q$ je druhý priesečník priamky~$MB$ s~$k$. Dokážte, že $|\angle PQA| = |\angle AQC|$.
}{
    Stred~$AP$ je zvláštny bod, ktorý akosi nesedí. Iba že by sme ho použili na stredovú súmernosť, ktorou preneseme jeden z~uhlov, ktorých rovnosť dokazujeme.
}{
    Dáva zmysel zrkadliť~$Q$ podľa~$M$, aby sme preniesli uhol~$PQA$; zrazu sa rovná~$PQ'A$. Tvrdenie, ktoré máme dokázať, nám potom po krátkom počítaní uhlov povie, čo stačí ukázať.
}{
    Bod $M$ vstupuje do zadania jedine ako stred úsečky $AP$, a~stred úsečky si pýta stredovú súmernosť podľa neho: tá vyrobí rovnobežník. Označme $Q'$ obraz bodu $Q$ v~tejto súmernosti; uhol $PQA$ sa ňou prenesie do vrcholu $Q'$.

    Označme $t$ dotyčnicu kružnice $k$ v~bode $A$, na ktorej ležia $P$ i~$M$. Pritom $P \neq A$, lebo $A$ neleží na priamke $BC$; preto je aj $M \neq A$ a~$M \neq P$. Priamka $t$ pretína priamku $AB$ jedine v~$A$ a~priamku $BC$ jedine v~$P$, takže $M$ neleží ani na jednej z~nich; keďže $Q$ leží na priamke $MB$, dostávame $Q \neq A$ aj $Q \neq C$. Bod $Q$ preto neleží na $t$ (tá má s~$k$ spoločný jediný bod $A$), a~keďže $M$ je stredom úsečky $AP$ i~úsečky $QQ'$, je $AQPQ'$ rovnobežník. Jeho protiľahlé uhly dávajú
    $$
    |\angle PQA| = |\angle AQ'P|.
    $$

    Keďže $M$ leží na $t$ a~$M \neq A$, leží zvonka $k$ a~jeho mocnosť vzhľadom na $k$ je $|MA|^2 > 0$. Body $B$, $Q$ teda ležia na tej istej polpriamke z~$M$ a~$|MB| \cdot |MQ| = |MA|^2$. Súmernosť podľa $M$ zachováva $|MQ'| = |MQ|$ a~kladie $Q'$ na opačnú polpriamku, takže $M$ leží vnútri úsečky $BQ'$; vnútri úsečky $AP$ leží tiež a~$|MA| = |MP|$, čiže
    $$
    |MA| \cdot |MP| = |MB| \cdot |MQ'|.
    $$
    Podľa kritéria tetivovosti ležia $A$, $B$, $P$, $Q'$ na jednej kružnici. (Sú to štyri rôzne body: z~$Q' = A$ by vyšlo $Q = P$, ktorý na $k$ neleží, z~$Q' = P$ zase $Q = A$, a~$Q' \neq B$, lebo $M$ leží medzi nimi. Body $A$, $B$, $P$ nie sú kolineárne, lebo $B$ neleží na $t$.) Uhlopriečky štvoruholníka $ABPQ'$ sa pretínajú v~jeho vnútornom bode $M$, takže je konvexný a~jeho protiľahlé uhly pri vrcholoch $B$ a~$Q'$ dávajú
    $$
    |\angle AQ'P| = 180^\circ - |\angle ABP|.
    $$
    Zostáva nám teda dokázať
    $$
    |\angle AQC| + |\angle ABP| = 180^\circ. \eqno(*)
    $$

    Bod $P$ leží zvonka $k$, čiže mimo tetivy $BC$: buď leží $B$ medzi $P$ a~$C$, alebo $C$ medzi $P$ a~$B$. Na priamke $AC$ bod $P$ neleží, lebo tá pretína priamku $BC$ jedine v~$C$ a~$P \neq C$; body $A$, $P$, $C$ teda tvoria trojuholník, $M$ je vnútorným bodom jeho strany $AP$ a~priamka $BQ$ neprechádza žiadnym jeho vrcholom (cez $A$ alebo $P$ by to bola priamka $t$, na ktorej $B$ neleží, cez $C$ zase priamka $BC$, na ktorej neleží $M$). Priamka, ktorá vojde do trojuholníka vnútorným bodom jednej strany a~vyhne sa vrcholom, pretne vnútro práve jednej zo zvyšných dvoch strán. S~priamkou $PC$ má priamka $BQ$ spoločný jediný bod $B$: ak teda $B$ leží vnútri úsečky $PC$, je to práve tá pretnutá strana a~priamka $BQ$ nechá $A$ i~$C$ v~tej istej polrovine; ak $B$ vnútri úsečky $PC$ neleží, teda ak $C$ leží medzi $P$ a~$B$, pretne priamka $BQ$ vnútro úsečky $AC$ a~body $A$, $C$ oddelí.

    Dve tetivy kružnice sa pretínajú v~jej vnútornom bode práve vtedy, keď každá z~nich oddeľuje krajné body tej druhej. V~prvom prípade preto ležia $Q$ a~$B$ v~tej istej polrovine s~hraničnou priamkou $AC$, čiže $|\angle AQC| = |\angle ABC|$, a~polpriamky $BP$, $BC$ sú opačné, čiže $|\angle ABP| = 180^\circ - |\angle ABC|$. V~druhom prípade priamka $AC$ body $Q$ a~$B$ oddeľuje, čiže $|\angle AQC| = 180^\circ - |\angle ABC|$, a~polpriamky $BP$, $BC$ splývajú, čiže $|\angle ABP| = |\angle ABC|$. Tak či onak platí $(*)$, a~teda
    $$
    |\angle PQA| = |\angle AQ'P| = 180^\circ - |\angle ABP| = |\angle AQC|.
    $$
}

\EnvId{JDPS5QjvQROqZ2N6JEwq6-lichobeznik-a-body-na-kl}
\Problem{0}{}{
    Nech $ABCD$ je lichobežník s~$AB \parallel CD$. Body $K$ a~$L$ ležia postupne na~$AB$ a~$CD$ tak, že $|AK| \cdot |CL| = |BK| \cdot |DL|$. Body $P$ a~$Q$ ležia na úsečke~$KL$ tak, že $|\angle APB| = |\angle DCB|$ a~$|\angle CQD|=|\angle CBA|$. Dokážte, že $B$, $C$, $P$, $Q$ ležia na jednej kružnici.
}{
    Podmienka na polohy~$K$ a~$L$ je zaujímavá. Jej prepis do tvaru pomeru naznačuje, že v~hre je nejaká rovnoľahlosť. Čo tak zrkadliť pomocou nej niečo ďalšie?
}{
    Uvažujte rovnoľahlosť zobrazujúcu~$AB$ na $DC$; pri nej sa~$K$ zobrazí na~$L$. Ak pri nej zobrazíme~$P$ na~$P'$, tak body $K$, $L$, $P$, $Q$, $P'$ všetky ležia na jednej priamke. Úloha sa zmení na pekné cvičenie z~počítania uhlov.
}{
    Body $K$ a~$L$ berieme ako vnútorné body strán $AB$ a~$CD$. Keby bol $K$ krajným bodom strany $AB$, zadaná rovnosť by vynútila, že $L$ je príslušný krajný bod strany $CD$, a~úsečka $KL$ by splynula s~ramenom lichobežníka. Na ramene $AD$ však pre každý jeho bod $P$ platí $|\angle APB| < 180^\circ - |\angle DAB| = |\angle ADC|$ a~pre každý jeho bod $Q$ platí $|\angle CQD| < |\angle DAB|$, takže predpísané rovnosti by naraz dávali $|\angle DCB| < |\angle ADC|$ aj $|\angle ADC| < |\angle DCB|$. Na ramene $BC$ nevyhovuje dokonca ani jeden z~tých dvoch bodov, lebo tam je $|\angle APB| < |\angle DCB|$ a~$|\angle CQD| < |\angle CBA|$. Taká poloha teda do zadania nespadá. Zadanú podmienku potom vieme prepísať na pomer
    $$
    \frac{|AK|}{|KB|} = \frac{|DL|}{|LC|},
    $$
    ktorý hovorí, že $K$ a~$L$ delia rovnobežné strany $AB$ a~$DC$ v~rovnakom pomere. Presne tak si zodpovedajú body v~rovnoľahlosti zobrazujúcej $AB$ na $DC$, a~keď nám rovnobežné úsečky takú rovnoľahlosť ponúkajú, oplatí sa ňou zobraziť ďalší bod. Zobrazíme ňou bod $P$.

    Keďže $ABCD$ je lichobežník, ramená $AD$ a~$BC$ nie sú rovnobežné a~ich priamky sa pretínajú v~bode $S$. Keby $S$ ležal vnútri úsečky $AD$, ležal by striktne medzi rovnobežkami $AB$ a~$CD$, a~teda aj vnútri úsečky $BC$, čiže protiľahlé strany lichobežníka by sa pretínali. Bod $S$ preto na úsečke $AD$ neleží: body $A$, $D$ ležia na tej istej polpriamke vychádzajúcej z~bodu $S$ a~úsečka $AD$ nepretína priamku $BC$, takže $A$ a~$D$ ležia v~tej istej polrovine určenej priamkou $BC$.

    Nech $h$ je rovnoľahlosť so stredom $S$ a~koeficientom $k = |SD| \, / \, |SA| > 0$, takže $h(A) = D$. Priamka $AB$ sa pri nej zobrazí na rovnobežku s~$AB$ prechádzajúcu bodom $D$, čiže na priamku $CD$, a~priamka $SB = BC$ je samodružná; obraz $h(B)$ je teda priesečník priamok $CD$ a~$BC$, čiže $h(B) = C$. Lichobežník nie je rovnobežník, takže $|DC| \neq |AB|$ a~$k \neq 1$. Bod $S$ navyše neleží ani na jednej z~rovnobežiek $AB$, $CD$: je samodružný a~$h$ zobrazuje $AB$ na $CD$, takže by ležal na oboch.

    Rovnoľahlosť zachováva deliace pomery, takže $h(K)$ je ten bod úsečky $DC$, ktorý ju delí v~pomere $|AK| : |KB|$. Podľa prepísanej podmienky je to práve $L$, čiže
    $$
    h(K) = L.
    $$
    Bod $L$ tak leží na priamke $SK$: priamka $\ell = KL$ prechádza stredom $S$, je pri $h$ samodružná a~leží na nej aj obraz $P' = h(P)$.

    Koeficient $k$ je kladný, takže $h$ zobrazuje každú polpriamku z~bodu $S$ samu na seba. Body $K$, $L = h(K)$ a~$h(L)$ preto ležia na tej istej polpriamke z~$S$ a~ich vzdialenosti od $S$ sú postupne $|SK|$, $k|SK|$ a~$k^2|SK|$. Pre $k > 1$ táto trojica rastie, pre $k < 1$ klesá; v~oboch prípadoch je prostredná hodnota striktne medzi krajnými, takže $L$ leží medzi $K$ a~$h(L)$. Keďže $K$ aj $L$ ležia na tej istej polpriamke z~$S$ a~sú od $S$ rôzne, bod $S$ na úsečke $KL$ neleží.

    Keďže $0 < |\angle APB| = |\angle DCB| < 180^\circ$, bod $P$ neleží na priamke $AB$, špeciálne $P \neq K$. Obrazom úsečky $KL$ je úsečka $L\,h(L)$, takže $P'$ leží na nej a~$P' \neq L$. Bod $L$ teda oddeľuje $P'$ od $K$, a~keďže $\ell$ pretína priamku $CD$ jedine v~$L$, ležia $P'$ a~$K$ (a~s~ním celá priamka $AB$) v~opačných polrovinách určených priamkou $CD$.

    Rovnoľahlosť $h$ zobrazuje trojuholník $APB$ na trojuholník $DP'C$, takže
    $$
    |\angle DP'C| = |\angle APB| = |\angle DCB|.
    $$
    Bod $P'$ leží na $\ell$ a~$P' \neq L$, takže neleží na priamke $CD$; nech $\gamma$ je kružnica opísaná trojuholníku $DP'C$. Podľa vety o~úsekovom uhle zviera dotyčnica ku $\gamma$ v~bode $C$ s~polpriamkou $CD$ v~polrovine neobsahujúcej $P'$ uhol veľkosti $|\angle DP'C|$. Polpriamka $CB$ leží v~tej istej polrovine a~s~polpriamkou $CD$ zviera uhol $|\angle DCB| = |\angle DP'C|$; polpriamka z~bodu $C$ je však v~danej polrovine veľkosťou uhla s~$CD$ určená jednoznačne, takže priamka $BC$ sa dotýka $\gamma$ v~bode $C$.

    Priamka $BC$ je priečkou rovnobežiek $AB$ a~$CD$ a~body $A$, $D$ ležia pri nej v~tej istej polrovine, takže $|\angle CBA| + |\angle BCD| = 180^\circ$, čiže
    $$
    |\angle DQC| = |\angle CBA| = 180^\circ - |\angle DCB| = 180^\circ - |\angle DP'C|.
    $$
    Keďže $0 < |\angle DQC| < 180^\circ$, bod $Q$ neleží na priamke $CD$, čiže $Q \neq L$, a~keďže úsečka $KL$ pretína priamku $CD$ jedine v~$L$, leží $Q$ v~tej istej polrovine ako $K$, teda na opačnej strane priamky $CD$ než $P'$. Oblúk kružnice $\gamma$ ležiaci v~tejto polrovine je pritom množinou práve tých jej bodov, z~ktorých vidno $DC$ pod uhlom $180^\circ - |\angle DP'C|$ (obvodové uhly z~opačných oblúkov sa dopĺňajú do $180^\circ$), a~taká množina je jediná. Preto $Q$ leží na $\gamma$.

    Priamky $\ell$ a~$BC$ sú rôzne, lebo inak by $K$ bol priesečníkom priamok $AB$ a~$BC$, čiže vrchol $B$; obe prechádzajú bodom $S$, takže $\ell \cap BC = \{S\}$, a~$S$ na úsečke $KL$ neleží. Body $P$, $Q$ preto neležia na priamke $BC$ a~body $B$, $C$, ktoré sú od $S$ rôzne, neležia na $\ell$. Na $\ell$ neleží ani $D$, keďže $\ell$ pretína priamku $CD$ jedine vo vnútornom bode $L$ strany $CD$. Ak $P = Q$, tvrdenie platí (body $B$, $C$, $P$ nie sú kolineárne, stačí vziať kružnicu opísanú trojuholníku $BCP$); nech teda $P \neq Q$, takže $B$, $C$, $P$, $Q$ sú navzájom rôzne.

    Ďalej počítame s~orientovanými uhlami priamok modulo $180^\circ$, čím sa vyhneme rozboru vzájomných polôh bodov $P$, $Q$, $P'$ na priamke $\ell$; konfiguračné fakty, na ktorých úloha naozaj stojí, sme už použili. Rovnoľahlosť $h$ zobrazuje priamku $PB$ na priamku $P'C$, čiže $PB \parallel P'C$, a~body $P$, $Q$, $P'$ ležia na $\ell$. Preto
    $$
    \angle(PB, PQ) = \angle(P'C, P'Q) = \angle(DC, DQ) = \angle(CB, CQ),
    $$
    kde druhá rovnosť sú obvodové uhly nad tetivou $CQ$ kružnice $\gamma$ z~bodov $P'$ a~$D$ a~tretia je úsekový uhol pre dotyčnicu $CB$ a~tú istú tetivu. Rovnosť $\angle(PB,PQ) = \angle(CB,CQ)$ znamená, že body $B$, $C$, $P$, $Q$ ležia na jednej kružnici alebo na jednej priamke; priamka je vylúčená, lebo $B$ neleží na $\ell$.
}

\EnvId{Kr9IKyGuib12QPGum8wz1-styri-body-na-stranach-stvoruholnika}
\Problem{0}{CPSJ 2019, súťaž družstiev}{
    Nech $ABCD$ je tetivový štvoruholník. Body $K$, $L$, $M$, $N$ ležia postupne na stranách~$AB$, $BC$, $CD$, $DA$ tak, že $|\angle ADK| = |\angle BCK|$, $|\angle BAL| = |\angle CDL|$, $|\angle CBM| = |\angle DAM|$ a~$|\angle DCN|=|\angle ABN|$. Dokážte, že $KM \perp LN$.
}{
    Nakreslite len jeden bod, povedzme~$K$, a~skúste lepšie pochopiť zvláštnu uhlovú podmienku, ktorou je definovaný. Pridávanie ďalších bodov začne dávať zmysel až vtedy, keď niečo objavíte.
}{
    V~obrázku len s~bodom~$K$ vieme počítaním uhlov dostať $|\angle BKC|=|\angle KDC|$. Priamka~$AB$ je preto dotyčnicou ku kružnici opísanej~$KCD$. To pripomína mocnosť bodu: treba len pridať bod, z~ktorého ju počítame. Keď je toto jasné, vieme pospájať zvyšné body~$L$, $M$, $N$.
}{
    Body $A$, $B$, $C$, $D$ ležia v~tomto poradí na kružnici $\omega$, takže štvoruholník $ABCD$ je konvexný. Všetky štyri podmienky zo zadania sú pritom jedna a~tá istá: bod $X$ leží na strane $A_1A_2$ štvoruholníka $A_1A_2A_3A_4$ a~platí $|\angle A_1A_4X| = |\angle A_2A_3X|$, postupne pre
    $$
    (A_1,A_2,A_3,A_4) = (A,B,C,D),\ (B,C,D,A),\ (C,D,A,B),\ (D,A,B,C).
    $$
    Cyklický posun označenia $(A,B,C,D) \to (B,C,D,A)$ teda zachováva zadanie a~štvoricu $(K,L,M,N)$ prevedie na $(L,M,N,K)$. Každé tvrdenie o~bode $K$ tak stačí dokázať raz.

    Keď sa $X$ pohybuje po strane $AB$ z~$A$ do $B$, uhol $\angle ADX$ rastie od nuly po $|\angle ADB|$ a~uhol $\angle BCX$ klesá od $|\angle BCA|$ po nulu. Bod $K$ preto existuje, je jediný a~je vnútorným bodom strany $AB$; to isté platí pre $L$, $M$, $N$.

    Vybavme si najprv prípad, keď sú niektoré dve protiľahlé strany rovnobežné; vďaka posunu môžeme predpokladať $AB \parallel CD$. Potom je $ABCD$ rovnoramenný lichobežník a~os $s$ úsečky $AB$ je zároveň osou úsečky $CD$. Osová súmernosť podľa $s$ zamení $A \leftrightarrow B$ a~$C \leftrightarrow D$, čiže zachová konfiguráciu s~preznačením $(A,B,C,D) \to (B,A,D,C)$, pri ktorom prejde $(K,L,M,N)$ na $(K,N,M,L)$. Z~jednoznačnosti teda táto súmernosť zobrazí $K$ na $K$, $M$ na $M$ a~$L$ na $N$. Body $K$, $M$ preto ležia na $s$ a~sú rôzne, takže $KM = s$, kým $LN$ je na $s$ kolmá. Ďalej už predpokladajme, že žiadne dve protiľahlé strany rovnobežné nie sú.

    Kľúčové pozorovanie sa dá spraviť v~obrázku s~jediným novým bodom. Označme $\varphi = |\angle ADK| = |\angle BCK|$, $\beta = |\angle ABC|$ a~$\delta = |\angle CDA|$. Polpriamka $BK$ splýva s~polpriamkou $BA$, takže zo súčtu uhlov v~trojuholníku $BCK$ a~z~$\beta + \delta = 180^\circ$ dostaneme
    $$
    |\angle BKC| = 180^\circ - \beta - \varphi = \delta - \varphi.
    $$
    Bod $K$ je vnútorným bodom strany $AB$ konvexného štvoruholníka, takže polpriamka $DK$ leží vnútri uhla $ADC$, takže
    $$
    |\angle KDC| = \delta - \varphi = |\angle BKC|.
    $$
    Úsečka $KC$ delí štvoruholník na trojuholník $KBC$ a~štvoruholník $AKCD$, čiže body $B$ a~$D$ ležia v~opačných polrovinách priamky $KC$. Podľa obrátenej vety o~úsekovom uhle je preto priamka $AB$ dotyčnicou kružnice opísanej trojuholníku $KCD$ v~bode $K$.

    Dotyčnica je nám na nič, kým nemáme bod, z~ktorého mocnosť počítame. Stačí pretnúť dve vhodné priamky: nech $P$ je priesečník priamok $AB$ a~$CD$. Keďže $ABCD$ je konvexný, leží $P$ mimo úsečiek $AB$ aj $CD$, teda vonku $\omega$, a~mocnosť bodu $P$ ku kružnici $(KCD)$ raz cez dotyčnicu $PK$ a~raz cez sečnicu $CD$ dáva
    $$
    |PK|^2 = |PC| \cdot |PD| = |PA| \cdot |PB|,
    $$
    kde druhá rovnosť je mocnosť bodu $P$ k~$\omega$. Dvojnásobný posun $(A,B,C,D) \to (C,D,A,B)$ prevedie $K$ na $M$ a~bod $P$ na priesečník priamok $CD$, $AB$, čiže zase na $P$. Preto $|PM|^2 = |PA| \cdot |PB|$, a~teda
    $$
    |PK| = |PM|.
    $$
    Po jednom, resp. troch posunoch dostaneme rovnako $|QL| = |QN|$, kde $Q$ je priesečník priamok $BC$ a~$DA$.

    Body $K$, $M$ sú rôzne, lebo ležia na rôznych priamkach $AB$, $CD$, ktorých jediný spoločný bod $P$ nie je vnútorným bodom žiadnej strany; rovnako $L \neq N$. Označme $p$ os úsečky $KM$ a~$q$ os úsečky $LN$. Z~$|PK| = |PM|$ leží $P$ na $p$, takže osová súmernosť podľa $p$ nechá $P$ na mieste a~zamení $K$ s~$M$, čím zobrazí priamku $AB$ na priamku $CD$. Rovnako súmernosť podľa $q$ zobrazí $BC$ na $DA$.

    Ostáva porovnať smery osí $p$ a~$q$. Tu prejdeme na orientované uhly priamok modulo $180^\circ$, zapisované $\angle(g,h)$; výsledok vyjde nulový, takže rozbor vzájomných polôh úplne odpadne. Pre osovú súmernosť podľa priamky $p$ a~ľubovoľnú priamku $g$ s~obrazom $g'$ platí $\angle(g,p) = \angle(p,g')$, čo pre $g = AB$, $g' = CD$ dáva $\angle(AB,CD) = 2\angle(AB,p)$, a~rovnako $\angle(BC,DA) = 2\angle(BC,q)$. Preto
    $$
    \eqalign{
        2\angle(p,q) &= 2\angle(p,AB) + 2\angle(AB,BC) + 2\angle(BC,q) \cr
        &= -\angle(AB,CD) + 2\angle(AB,BC) + \angle(BC,DA) \cr
        &= \angle(AB,DA) - \angle(BC,CD). \cr
    }
    $$
    Tetivovosť štvoruholníka $ABCD$ znamená práve $\angle(AB,DA) = \angle(CB,CD)$, takže pravá strana je nula a~$\angle(p,q)$ je $0^\circ$ alebo $90^\circ$. Keďže $KM \perp p$ a~$LN \perp q$, priamky $KM$ a~$LN$ sú rovnobežné alebo kolmé.

    Rovnobežnosť vylúčime polohou bodov. Prienikom priamky $KM$ s~konvexným štvoruholníkom $ABCD$ je úsečka $KM$ a~jej vnútorné body ležia vnútri štvoruholníka, lebo $K$, $M$ sú vnútornými bodmi protiľahlých strán. Bod $L$ leží na hranici štvoruholníka a~je rôzny od $K$ i~$M$, takže na priamke $KM$ neleží; rovnako $N$. Úsečka $KM$ pritom delí štvoruholník na $KBCM$ a~$AKMD$, pričom $L$ je vnútorným bodom strany $BC$ prvého z~nich a~$N$ vnútorným bodom strany $DA$ druhého. Body $L$ a~$N$ tak ležia v~opačných otvorených polrovinách priamky $KM$, čiže priamky $KM$ a~$LN$ rovnobežné nie sú. Sú teda kolmé.
}

\EnvId{t8rjDz8xH7IC_FmnHz15z-stvoruholnik-s-tromi-zhodnymi-stranami}
\Problem{0}{CAPS 2024, P4}{
    Nech $ABCD$ je štvoruholník taký, že $|AB| = |BC| = |CD|$. Na polpriamkach $CA$, $BD$ sú postupne body $X$, $Y$ také, že $|BX| = |CY|$. Nech $P$, $Q$, $R$, $S$ sú stredy úsečiek $BX$, $CY$, $XD$, $YA$. Dokážte, že body $P$, $Q$, $R$, $S$ ležia na jednej kružnici.
}{
    Máme veľa stredov. Zrkadlíme, alebo zavedieme ešte viac? Tu dáva viac zmyslu to druhé.
}{
    Kľúčovým bodom na zavedenie je stred~$XY$, keďže vytvorí veľa stredných priečok. A~tie dajú~jasnú cestu k~tomu, čo použiť na dôkaz toho, že body ležia na jednej kružnici.
}{
    Body $P$, $R$ sú stredy strán $XB$, $XD$ trojuholníka $XBD$, takže $PR \parallel BD$; rovnako $Q$, $S$ sú stredy strán $YC$, $YA$ trojuholníka $YCA$, takže $QS \parallel CA$ (z~$B \neq D$ a~$A \neq C$ je $P \neq R$ a~$Q \neq S$, obe priamky sú teda dobre definované). Samotné rovnobežnosti o~kružnici nehovoria nič, chýba nám priesečník priamok $PR$, $QS$ a~dĺžky odmerané z~neho. Keďže v~úlohe vystupujú samé stredy úsečiek, pridajme ešte jeden, ktorý s~každým z~nich vytvorí strednú priečku: stred $M$ úsečky $XY$.

    Potom je $MP$ stredná priečka trojuholníka $XBY$, $MR$ trojuholníka $XDY$, $MQ$ trojuholníka $YCX$ a~$MS$ trojuholníka $YAX$. Všetky štyri sa dajú zapísať naraz: ak body chápeme ako polohové vektory, tak z~$M = \tfrac12(X+Y)$ a~definícií bodov $P$, $Q$, $R$, $S$ dostaneme
    $$
    \gather
    P - M = \tfrac12 (B - Y), \qquad R - M = \tfrac12 (D - Y), \\
    Q - M = \tfrac12 (C - X), \qquad S - M = \tfrac12 (A - X).
    \endgather
    $$

    Bod $Y$ leží na priamke $BD$, takže $P - M$ aj $R - M$ sú rovnobežné s~$BD$ a~$M$ leží na priamke $PR$; rovnako $X$ leží na priamke $CA$, takže $M$ leží na priamke $QS$. Uhlopriečky $AC$ a~$BD$ štvoruholníka $ABCD$ nie sú rovnobežné, preto sú $PR$ a~$QS$ rôzne priamky a~ich jediným spoločným bodom je $M$.

    Ďalej počítajme s~orientovanými dĺžkami na priamkach: pre kolineárne body súčin $\overline{UV} \cdot \overline{UW}$ od voľby orientácie nezávisí a~ušetrí nám rozbor toho, či $X$ leží vnútri úsečky $AC$ alebo mimo nej. Ak na priamke $PR$ zvolíme tú istú orientáciu ako na $BD$ a~na $QS$ tú istú ako na $CA$, predošlé rovnosti dávajú
    $$
    \gather
    \overline{MP} \cdot \overline{MR} = \tfrac14\, \overline{YB} \cdot \overline{YD}, \\
    \overline{MQ} \cdot \overline{MS} = \tfrac14\, \overline{XC} \cdot \overline{XA}.
    \endgather
    $$
    Podľa kritéria tetivovosti tak stačí dokázať rovnosť
    $$
    \overline{YB} \cdot \overline{YD} = \overline{XC} \cdot \overline{XA};
    $$
    ak je spoločná hodnota nenulová, sú body $P$, $Q$, $R$, $S$ všetky rôzne od $M$ a~kritérium dáva hľadanú kružnicu.

    Na oboch stranách stojí mocnosť bodu, treba len nájsť tie správne kružnice. Označme $s = |AB| = |BC| = |CD|$. Rovnosť $|BA| = |BC| = s$ hovorí presne to, že kružnica $\omega_B$ so stredom $B$ a~polomerom $s$ prechádza bodmi $A$, $C$; rovnako $|CB| = |CD| = s$ hovorí, že kružnica $\omega_C$ so stredom $C$ a~polomerom $s$ prechádza bodmi $B$, $D$. Priamka $CA$ pretína $\omega_B$ práve v~bodoch $A$ a~$C$ a~bod $X$ na nej leží, takže
    $$
    \overline{XC} \cdot \overline{XA} = |BX|^2 - s^2,
    $$
    a~rovnako $\overline{YB} \cdot \overline{YD} = |CY|^2 - s^2$. Podmienka $|BX| = |CY|$ tieto dve mocnosti stotožňuje, čím je požadovaná rovnosť dokázaná.

    Zostáva prípad, keď je spoločná hodnota $|BX|^2 - s^2$ nulová, teda $|BX| = |CY| = s$. Vtedy $X$ leží na $\omega_B$, čiže $X \in \{A, C\}$, a~$Y$ leží na $\omega_C$, čiže $Y \in \{B, D\}$. V~každej zo štyroch takých možností dva zo štyroch bodov splynú a~zvyšné tri sú stredy troch po sebe idúcich strán štvoruholníka $ABCD$: napríklad pre $X = A$, $Y = B$ je $S = P$ stred strany $AB$ a~zvyšné dva body sú stredy strán $DA$ a~$BC$. Sú to tri po sebe idúce vrcholy Varignonovho rovnobežníka, ktorého susedné strany sú rovnobežné s~$AC$ a~$BD$; tie rovnobežné nie sú, takže naše tri body neležia na jednej priamke a~kružnica cez ne existuje.
}

\EnvId{unjF6AEffP5thl4GWrGD5-uhlopriecky-lichobeznika-a-vnutorny-bod}
\Problem{0}{}{
    Uhlopriečky lichobežníka~$ABCD$ s~$AB \parallel CD$ sa pretínajú v~$P$. Bod~$Q$ leží vo vnútri trojuholníka~$BCP$ tak, že $|\angle AQB|= |\angle CQD|$. Dokážte, že $|\angle DQP| = |\angle BAQ|$.
}{
    Bod~$Q$ je definovaný pomerne ľubovoľne. Nezostáva než pridať nejaký ďalší bod, ktorý prinesie zmysluplnú štruktúru. Samotný lichobežník dáva celkom dobrý náznak.
}{
    Zrkadlite~$Q$ v~rovnoľahlosti so stredom~$P$, ktorá zobrazuje~$AB$ na~$CD$. Zrazu sa definícia~$Q$ zmení na niečo zmysluplné a~z~ničoho nič máme všetko, čo potrebujeme; už to len pospájame.
}{
    Rovnobežné strany $AB$, $CD$ ponúkajú rovnoľahlosť so stredom $P$, ktorá jednu z~nich zobrazí na druhú. Podmienka $|\angle AQB| = |\angle CQD|$ zväzuje dvojicu $A$, $B$ s~dvojicou $C$, $D$ presne tak, ako to robí táto rovnoľahlosť, tak ňou zobrazme aj bod $Q$.

    Lichobežník $ABCD$ je konvexný, takže $P$ je vnútorným bodom oboch úsečiek $AC$, $BD$. Z~rovnobežnosti $AB \parallel CD$ sú trojuholníky $PAB$ a~$PCD$ podobné, čiže
    $$
    \frac{|PC|}{|PA|} = \frac{|PD|}{|PB|}.
    $$
    Rovnoľahlosť $h$ so stredom $P$ a~koeficientom $k = -|PC| \, / \, |PA| < 0$ preto zobrazuje $A \mapsto C$ a~$B \mapsto D$ (záporné znamienko zodpovedá tomu, že $P$ leží vnútri oboch uhlopriečok). Označme $Q_1 = h(Q)$.

    Rovnoľahlosť zachováva veľkosti uhlov, takže
    $$
    |\angle CQ_1D| = |\angle AQB| = |\angle CQD|.
    $$

    Aby sme z~toho dostali tetivový štvoruholník, potrebujeme ešte, že $Q$ a~$Q_1$ ležia na tej istej strane priamky $CD$. Nech $\pi$ je otvorená polrovina s~hraničnou priamkou $CD$ obsahujúca bod $A$. Bod $B$ leží na priamke rovnobežnej s~$CD$ a~od nej rôznej, takže $B \in \pi$, a~$P$ je vnútorným bodom úsečky $AC$ s~$A \in \pi$, takže tiež $P \in \pi$. Vnútro trojuholníka $BCP$, ktorého jediný vrchol na priamke $CD$ je $C$, teda celé leží v~$\pi$; špeciálne $Q \in \pi$. Keďže $h(B) = D$, $h(C) = A$ a~$h(P) = P$, zobrazí sa trojuholník $BCP$ na trojuholník $DAP$ a~$Q_1$ leží v~jeho vnútri. Rovnakou úvahou (na priamke $CD$ leží z~jeho vrcholov jedine $D$) je aj $Q_1 \in \pi$.

    Body $Q$, $Q_1$ ležia teda na tej istej strane priamky $CD$ a~vidia úsečku $CD$ pod rovnakým uhlom, takže podľa obrátenej vety o~obvodovom uhle ležia $C$, $D$, $Q$, $Q_1$ na jednej kružnici; označme ju $\omega$.

    Koeficient $k$ je záporný a~$Q \neq P$, takže $P$ je vnútorným bodom úsečky $QQ_1$. Polpriamky $QP$ a~$QQ_1$ preto splývajú a~platí $|\angle DQP| = |\angle DQQ_1|$.

    Bod $P$ je vnútorným bodom tetivy $QQ_1$, čiže leží vnútri $\omega$. Pre bod vnútri kružnice sa cyklické poradie bodov na kružnici zhoduje s~cyklickým poradím polpriamok, ktoré k~nim z~neho vedú. Polpriamka $PC$ je opačná k~$PA$, polpriamka $PD$ je opačná k~$PB$ a~priamka $AC$ oddeľuje $B$ od $D$, takže okolo $P$ idú tieto štyri polpriamky v~poradí $PA$, $PB$, $PC$, $PD$. Bod $Q$ leží vo vnútri trojuholníka $BCP$, takže polpriamka $PQ$ leží vnútri uhla $\angle BPC$ a~opačná polpriamka $PQ_1$ vnútri uhla $\angle DPA$. Okolo $P$ tak polpriamky idú v~poradí
    $$
    PA, \quad PB, \quad PQ, \quad PC, \quad PD, \quad PQ_1,
    $$
    a~cyklické poradie bodov na $\omega$ je $Q$, $C$, $D$, $Q_1$.

    Tetiva $DQ_1$ preto necháva body $C$ a~$Q$ na tom istom oblúku a~z~vety o~obvodovom uhle dostávame $|\angle DQQ_1| = |\angle DCQ_1|$. Napokon $h$ zobrazuje $A \mapsto C$, $B \mapsto D$, $Q \mapsto Q_1$, takže $|\angle BAQ| = |\angle DCQ_1|$. Spojením troch rovností
    $$
    |\angle DQP| = |\angle DQQ_1| = |\angle DCQ_1| = |\angle BAQ|.
    $$
}

\EnvId{cvVqE4nT1JdVqFMok0P_I-osi-useciek-bd-a-ce}
\Problem{0}{IMO 2018, P1}{
    Nech $ABC$ je ostrouhlý trojuholník s~opísanou kružnicou $\Gamma$. Body $D$ a $E$ ležia vo vnútri strán $AB$ a $AC$ tak, že $|AD| = |AE|$. Osi úsečiek $BD$ a $CE$ pretínajú kratšie oblúky $AB$ a $AC$ kružnice $\Gamma$ postupne v~bodoch $F$ a $G$. Dokážte, že priamky $DE$ a $FG$ sú rovnobežné (alebo splývajú).
}{
    Máme $|FB| = |FD|$, ale nie je jasné, ako to využiť. Dáva nám to rovnaké uhly $|\angle FBD|$ a $|\angle FDB|$. Najmä ten druhý vyzerá osamotene. Skúste pretnúť niečo s~niečím, aby ste tieto uhly niekam preniesli.
}{
    Kľúčom je predĺžiť úsečku $FD$ na tetivu $\Gamma$ a~podobne predĺžiť $GE$ na druhej strane. Zrazu dostaneme rovnoramenné trojuholníky, ktoré dobre interagujú s~podmienkou $|AD| = |AE|$. Ďalšie body netreba; stane sa z~toho cvičenie z~počítania uhlov.
}{
    Označme $\alpha$, $\beta$, $\gamma$ vnútorné uhly trojuholníka $ABC$. Trojuholník je ostrouhlý, takže oblúk $AB$ neobsahujúci $C$ má veľkosť $2\gamma < 180^\circ$ a~je teda ten kratší; bod $F$ preto leží na oblúku $AB$ neobsahujúcom $C$ a~rovnako $G$ leží na oblúku $AC$ neobsahujúcom $B$.

    Z~$|FB| = |FD|$ máme $|\angle FDB| = |\angle FBD|$, lenže uhol pri $D$ nemá na priamke $AB$ s~čím interagovať. Prenesieme ho na kružnicu presne v~duchu poslednej zo situácií zo zoznamu vyššie: predĺžime $FD$ na tetivu $\Gamma$. Nech teda $T$ je druhý priesečník priamky $FD$ s~$\Gamma$ a~nech $S$ je druhý priesečník priamky $GE$ s~$\Gamma$.

    Bod $D$ je vnútorným bodom tetivy $AB$, leží teda vnútri $\Gamma$, a~preto na priamke $FT$ leží medzi $F$ a~$T$. Body $F$ a~$C$ ležia v~opačných polrovinách určených priamkou $AB$; keďže $D$ leží na priamke $AB$ medzi $F$ a~$T$, leží bod $T$ v~tej istej polrovine ako $C$, čiže na oblúku $AB$ obsahujúcom $C$. Rovnako $E$ leží medzi $G$ a~$S$ a~bod $S$ leží na oblúku $AC$ obsahujúcom $B$.

    Označme $\varphi = |\angle ABF|$. Bod $D$ leží medzi $A$ a~$B$ aj medzi $F$ a~$T$, takže z~vrcholových uhlov
    $$
    |\angle ADT| = |\angle BDF| = |\angle DBF| = \varphi.
    $$
    Oblúk $AF$ neobsahujúci $B$ je časťou oblúka $AB$ neobsahujúceho $C$, takže body $B$ a~$T$ ležia na tom istom z~oblúkov určených tetivou $AF$ a~obvodové uhly nad ňou dávajú
    $$
    |\angle ATD| = |\angle ATF| = |\angle ABF| = \varphi.
    $$
    Trojuholník $ADT$ je preto rovnoramenný a~$|AT| = |AD|$. Zámenou $B \leftrightarrow C$, $D \leftrightarrow E$, $F \leftrightarrow G$, $T \leftrightarrow S$ rovnako $|AS| = |AE|$, a~keďže $|AD| = |AE|$, dostávame
    $$
    |AT| = |AS|.
    $$

    Polpriamky $AD$ a~$AB$ splývajú, takže zo súčtu uhlov v~trojuholníku $ADT$ je $|\angle BAT| = 180^\circ - 2\varphi$. Bod $T$ leží na oblúku $AB$ obsahujúcom $C$, čiže $|\angle ATB| = \gamma$, a~súčet uhlov v~trojuholníku $ABT$ dáva
    $$
    |\angle ABT| = 180^\circ - (180^\circ - 2\varphi) - \gamma = 2\varphi - \gamma.
    $$
    Analogicky pre $\psi = |\angle ACG|$ vyjde $|\angle ACS| = 2\psi - \beta$.

    Uhly $\angle ABT$ a~$\angle ACS$ sú obvodové uhly kružnice $\Gamma$ nad zhodnými tetivami $AT$ a~$AS$. Zhodné tetivy vytínajú z~kružnice zhodnú dvojicu oblúkov a~obvodový uhol nad tetivou je polovicou toho z~jej dvoch oblúkov, ktorý neobsahuje vrchol uhla; uhly $\angle ABT$ a~$\angle ACS$ sú preto zhodné alebo majú súčet $180^\circ$.

    Ukážeme, že oba sú ostré, čím druhú možnosť vylúčime. Body $B$ a~$C$ ležia na tom istom z~oblúkov určených tetivou $AF$, takže $\varphi = |\angle ACF|$, a~keďže úsečka $CF$ pretína tetivu $AB$, leží polpriamka $CF$ vnútri uhla $ACB$ a~$\varphi < \gamma$. Číslo $2\varphi - \gamma$ je uhol trojuholníka $ABT$, čiže kladné, dokopy $0 < 2\varphi - \gamma < \gamma < 90^\circ$; rovnako $0 < 2\psi - \beta < \beta < 90^\circ$. Súčet dvoch ostrých uhlov je menší ako $180^\circ$, takže nastáva prvá možnosť:
    $$
    2\varphi - \gamma = 2\psi - \beta. \eqno(*)
    $$

    Nech $M$ je stred oblúka $BC$ neobsahujúceho $A$, teda druhý priesečník osi uhla $BAC$ s~$\Gamma$. Trojuholník $ADE$ je rovnoramenný, os uhla $BAC$ je preto osou úsečky $DE$ a~platí $AM \perp DE$. Stačí teda dokázať, že aj $AM \perp FG$.

    Na tom z~oblúkov určených tetivou $FG$, ktorý obsahuje $A$, ležia iba body oblúkov $FA$ a~$AG$, kým $B$, $M$, $C$ ležia na druhom. Tetiva $FG$ tak oddeľuje $A$ od $M$ a~úsečka $AM$ pretína $FG$ vo vnútornom bode; označme ho $X$. Body $C$ a~$F$ ležia na tom istom z~oblúkov určených tetivou $AG$ (na tom, ktorý obsahuje $B$), takže $|\angle AFG| = |\angle ACG| = \psi$. Ďalej $|\angle AFB| = 180^\circ - \gamma$, lebo $F$ a~$C$ ležia na opačných oblúkoch určených tetivou $AB$, a~zo súčtu uhlov v~trojuholníku $ABF$ je $|\angle FAB| = \gamma - \varphi$. Polpriamka $AM$ ide vnútrom uhla $BAC$, kým $F$ leží v~opačnej polrovine od $C$ vzhľadom na priamku $AB$; polpriamka $AB$ preto leží vnútri uhla $FAM$ a
    $$
    |\angle FAX| = |\angle FAB| + |\angle BAM| = \gamma - \varphi + \tfrac\alpha2.
    $$
    Zo vzťahu $(*)$ je $\varphi - \psi = \frac{\gamma - \beta}2$, čiže $\gamma - \varphi + \psi = \frac{\beta + \gamma}2$, a~súčet uhlov v~trojuholníku $AFX$ dáva
    $$
    \gather
    |\angle AXF| = 180^\circ - (\gamma - \varphi + \tfrac\alpha2) - \psi \\
    = 180^\circ - \tfrac{\beta + \gamma}2 - \tfrac\alpha2 = 90^\circ.
    \endgather
    $$
    Priamky $DE$ a~$FG$ sú tak obe kolmé na priamku $AM$, čiže sú rovnobežné alebo splývajú.
}

\EnvId{Dk4xgZzPmhroFmRx3omVb-rovnobezne-dotycnice-k-vpisanej-kruznici}
\Problem{0}{IMO 2024, P4}{
    Nech $ABC$ je trojuholník taký, že $|AB| < |AC| < |BC|$. Nech $\omega$ je kružnica vpísaná trojuholníku $ABC$ a $I$ jej stred. Nech $X$ je bod na priamke $BC$ rôzny od $C$ taký, že priamka prechádzajúca bodom $X$ rovnobežne s~$AC$ sa dotýka $\omega$. Podobne nech $Y$ je bod na priamke $BC$ rôzny od $B$ taký, že priamka prechádzajúca bodom $Y$ rovnobežne s~$AB$ sa dotýka $\omega$. Priamka $AI$ pretína kružnicu opísanú trojuholníku $ABC$ druhýkrát v~bode $P \neq A$. Nech $K$ a $L$ sú stredy strán $AC$ a $AB$. Dokážte, že $|\angle KIL| + |\angle YPX| = 180^\circ$.
}{
    Tými rovnobežnými priamkami máme skoro daný rovnobežník.
}{
    Nech~$Z$ je priesečník rovnobežných priamok zo zadania; potom pomocou rovnoľahlosti so stredom~$A$ a~koeficientom~$2$ máme, že $|\angle KIL| = |\angle BZC|$. Tie škaredé stredy už nemáme, lebo dokazujeme len $|\angle BZC| + |\angle YPX| = 180^\circ$.
}{
    Kľúčom k~dotiahnutiu úlohy je nájsť tetivové štvoruholníky. Ale aspoň netreba zavádzať viac bodov.
}{
    Dotyčnice zo zadania sú rovnobežné so stranami $AC$ a~$AB$, a~pri rovnobežkách sa oplatí uvažovať o~zobrazení, ktoré jednu z~nich pošle na druhú. Tu je také poruke: stredová súmernosť podľa $I$ zobrazuje $\omega$ samu na seba, teda dotyčnicu na dotyčnicu, a~každú priamku na rovnobežku s~ňou. Nech $Z$ je obraz bodu $A$ v~tejto súmernosti, teda bod, pre ktorý je $I$ stredom úsečky $AZ$.

    Priamka vedená bodom $X$ rovnobežne s~$AC$ je dotyčnicou $\omega$ rôznou od priamky $AC$: keby s~$AC$ splývala, ležal by $X$ na priamkach $AC$ aj $BC$, čiže by bolo $X = C$, čo zadanie vylučuje. Dotyčnice $\omega$ daného smeru sú však len dve, takže táto priamka je obrazom priamky $AC$ v~súmernosti podľa $I$, a~preto prechádza obrazom bodu $A$, čiže bodom $Z$. Rovnako priamka vedená bodom $Y$ rovnobežne s~$AB$ je obrazom priamky $AB$ a~tiež prechádza bodom $Z$. Bod $Z$ je teda priesečníkom oboch dotyčníc zo zadania.

    Rovnoľahlosť so stredom $A$ a~koeficientom $2$ zobrazí $K$ na $C$, $L$ na $B$ a~$I$ na $Z$, pričom zachováva veľkosti uhlov. Preto $|\angle KIL| = |\angle CZB|$ a~zostáva dokázať
    $$
    |\angle BZC| + |\angle YPX| = 180^\circ.
    $$
    Nepríjemné stredy $K$, $L$ sa tým z~úlohy stratili.

    Označme $a = |BC|$, $b = |CA|$, $c = |AB|$, $\alpha = |\angle BAC|$, ďalej $s$ polobvod, $S$ obsah trojuholníka $ABC$, $r$ polomer $\omega$ a~$v_a$, $v_b$, $v_c$ jeho výšky. Zo vzťahov $S = rs$ a~$2S = a v_a = b v_b = c v_c$ dostaneme
    $$
    \frac{2r}{v_a} = \frac{a}{s}, \qquad \frac{2r}{v_b} = \frac{b}{s}, \qquad \frac{2r}{v_c} = \frac{c}{s},
    $$
    pričom všetky tri podiely sú menšie ako $1$, lebo trojuholníkové nerovnosti dávajú $a, b, c < s$.

    Merajme vzdialenosti od priamky $BC$ so znamienkom, kladne na strane bodu $A$. Bod $I$ má vzdialenosť $r$ a~je stredom úsečky $AZ$, takže vzdialenosť bodu $Z$ je $2r - v_a < 0$. Bod $Z$ teda leží v~opačnej polrovine ako $A$ a~úsečka $AZ$ pretína priamku $BC$; ten priesečník označme $T$. Leží na osi uhla $\angle BAC$, čiže vnútri úsečky $BC$.

    Druhá dotyčnica $\omega$ rovnobežná s~$AC$ má od priamky $AC$ vzdialenosť $2r$, a~to na tej strane, kde leží $I$, čiže kde leží aj $B$. Z~nerovnosti $2r < v_b$ preto oddeľuje $B$ od priamky $AC$, špeciálne od bodu $C$, a~tak pretína vnútro úsečky $BC$: bod $X$ leží vnútri $BC$. Rovnako z~$2r < v_c$ leží vnútri $BC$ aj bod $Y$.

    Nech $h$ je rovnoľahlosť so stredom $T$, ktorá zobrazuje $A$ na $Z$. Jej koeficient je záporný, lebo $T$ leží vnútri úsečky $AZ$. Priamku $BC$ zobrazuje $h$ na seba a~priamku $AC$ na rovnobežku s~$AC$ prechádzajúcu bodom $Z$, čo je priamka $ZX$; preto $h(C) = X$ a~rovnako $h(B) = Y$. Zo záporného koeficientu leží $X$ na polpriamke opačnej k~$TC$ a~$Y$ na polpriamke opačnej k~$TB$, takže na priamke $BC$ máme poradie $B$, $X$, $T$, $Y$, $C$. Ak označíme $k$ koeficient rovnoľahlosti $h$ v~absolútnej hodnote, tak
    $$
    |TZ| = k \, |TA|, \qquad |TX| = k \, |TC|, \qquad |TY| = k \, |TB|.
    $$

    Priamka $AI$ je osou uhla $\angle BAC$, takže $P$ je stredom oblúka $BC$ neobsahujúceho $A$; leží preto v~opačnej polrovine ako $A$ a~na priamke $BC$ neleží. Zo zadania je $|AB| \neq |AC|$, takže podľa tvrdenia o~Švrčkovom bode je $P$ stredom kružnice opísanej štvoruholníku $BICI_a$, kde $I_a$ je stred kružnice pripísanej oproti vrcholu $A$; špeciálne $|PI| = |PB|$.

    Nech $F$ je bod dotyku $\omega$ so stranou $AB$ a~nech $M$ je stred strany $BC$. Uhol pri $F$ v~trojuholníku $AFI$ je pravý a~$|AF| = s - a$. Bod $P$ je stredom oblúka $BC$, takže $|PB| = |PC|$ a~priamka $PM$ je osou úsečky $BC$; uhol pri $M$ v~trojuholníku $BMP$ je preto tiež pravý a~$|BM| = a/2$. Obvodové uhly nad tetivou $PC$ napokon dávajú $|\angle PBC| = |\angle PAC| = \alpha/2 = |\angle FAI|$. Pravouhlé trojuholníky $AFI$ a~$BMP$ sa tak zhodujú v~ostrom uhle, čiže sú podobné (s~$F \leftrightarrow M$, $I \leftrightarrow P$), a~porovnaním odvesny s~preponou dostaneme
    $$
    \frac{|AI|}{|AF|} = \frac{|BP|}{|BM|}, \qquad \text{teda} \qquad |AI| = \frac{2(s-a)}{a}\,|PB|.
    $$
    Keďže $|PI| = |PB|$, je preto
    $$
    \frac{|AI|}{|PI|} = \frac{2(s-a)}{a} = \frac{b+c-a}{a} < 1,
    $$
    lebo zo zadania je $b < a$ aj $c < a$; práve tu sa využíva, že $BC$ je najdlhšia strana.

    Bod $I$ leží vnútri kružnice opísanej, teda vnútri tetivy $AP$. Bod $Z$ leží na polpriamke $AI$ a~$|AZ| = 2|AI| < |AI| + |IP| = |AP|$, takže $Z$ leží vnútri úsečky $IP$. Na priamke $AP$ tak máme poradie $A$, $T$, $Z$, $P$.

    Tetivy $BC$ a~$AP$ kružnice opísanej sa pretínajú v~bode $T$, takže z~mocnosti bodu $T$ je $|TB| \cdot |TC| = |TA| \cdot |TP|$; túto spoločnú hodnotu označme $m$. Odtiaľ
    $$
    \gather
    |TB| \cdot |TX| = k \, m = |TZ| \cdot |TP|, \\
    |TC| \cdot |TY| = k \, m = |TZ| \cdot |TP|.
    \endgather
    $$
    Bod $T$ leží mimo úsečky $BX$ (poradie $B$, $X$, $T$), mimo úsečky $CY$ (poradie $T$, $Y$, $C$) aj mimo úsečky $ZP$ (poradie $T$, $Z$, $P$), takže podľa kritéria tetivovosti ležia na jednej kružnici body $B$, $X$, $Z$, $P$ a~na jednej kružnici aj body $C$, $Y$, $Z$, $P$.

    Body $Z$ a~$P$ ležia v~tej istej polrovine určenej priamkou $BC$, teda na tom istom oblúku prvej kružnice určenom tetivou $BX$. Obvodové uhly nad touto tetivou sa preto rovnajú, čiže $|\angle BZX| = |\angle BPX|$, a~z~druhej kružnice rovnako $|\angle CZY| = |\angle CPY|$.

    Body $B$, $X$, $Y$, $C$ ležia na priamke $BC$ v~tomto poradí a~body $Z$, $P$ na nej neležia, takže polpriamky z~$Z$ aj polpriamky z~$P$ idú k~nim v~tom istom poradí a~platí
    $$
    \gather
    |\angle BZC| = |\angle BZX| + |\angle XZY| + |\angle YZC|, \\
    |\angle BPC| = |\angle BPX| + |\angle XPY| + |\angle YPC|.
    \endgather
    $$

    Rovnoľahlosť $h$ zobrazuje trojuholník $ABC$ na trojuholník $ZYX$, takže $|\angle XZY| = \alpha$. Body $A$ a~$P$ ležia na kružnici opísanej po opačných stranách tetivy $BC$, takže $|\angle BPC| = 180^\circ - \alpha$. Po dosadení dvoch dokázaných rovností uhlov dostávame
    $$
    \gather
    |\angle BZC| = |\angle BPX| + \alpha + |\angle YPC|, \\
    |\angle YPX| = 180^\circ - \alpha - |\angle BPX| - |\angle YPC|,
    \endgather
    $$
    a~súčet týchto dvoch uhlov je $180^\circ$. Keďže $|\angle KIL| = |\angle BZC|$, je tvrdenie dokázané.
}

\EnvId{h_nSioAPSd9vwFsAInRyp-stredy-oblukov-dvoch-kruznic}
\Problem{0}{}{
    Nech $ABC$ je ostrouhlý trojuholník so stredom~$M$ strany~$BC$. Nech~$P \neq M$ je bod na strane~$BC$. Označme~$Q$ stred oblúka~$AC$ neobsahujúceho~$P$ kružnice opísanej trojuholníku~$APC$. Podobne~$R$ je stred oblúka~$AB$ neobsahujúceho~$P$ kružnice opísanej trojuholníku~$APB$. Dokážte, že $Q, M, P, R$ ležia na jednej kružnici.
}{
    Najproblematickejším bodom je~$M$: ako stred nemá okolo seba pekné uhly. Treba mu dať lepší význam. Aj tvrdenie, ktoré máme dokázať, naznačuje, čo chceme urobiť.
}{
    Kľúčom je zrkadliť napríklad~$R$ podľa~$M$. Potom sa oplatí zafarbiť rovnako dlhé úsečky; cieľom je zafarbiť ďalšiu dvojicu. Od želaného konca nie sme ďaleko.
}{
    Aby kružnice zo zadania existovali, je $P$ vnútorným bodom strany $BC$. Uhly trojuholníka $ABC$ označme ako obvykle $\alpha$, $\beta$, $\gamma$; trojuholník je ostrouhlý, takže všetky sú menšie než $90^\circ$. Položme $2\varphi = |\angle APB|$, čiže $|\angle APC| = 180^\circ - 2\varphi$ a~$0^\circ < \varphi < 90^\circ$.

    Najprv si upracme polohu bodov $Q$ a~$R$. Bod $R$ leží na oblúku $AB$ neobsahujúcom $P$, takže na kružnici $(APB)$ ležia body v~cyklickom poradí $A$, $P$, $B$, $R$. Tetiva $PB$ leží na priamke $BC$ a~body $A$, $R$ ležia na tom istom oblúku, ktorý táto tetiva vytína, preto $A$ a~$R$ ležia v~tej istej polrovine určenej priamkou $BC$. Tetiva $AB$ naopak oddeľuje $P$ od $R$, takže $R$ a~$C$ ležia v~opačných polrovinách určených priamkou $AB$ (bod $C$ je pri priamke $AB$ na tej istej strane ako $P$). Rovnaká úvaha pre kružnicu $(APC)$ s~cyklickým poradím $A$, $P$, $C$, $Q$ dáva, že $Q$ leží v~tej istej polrovine určenej priamkou $BC$ ako $A$ a~že $Q$, $B$ ležia v~opačných polrovinách určených priamkou $AC$. Špeciálne $Q$ ani $R$ neležia na priamke $BC$.

    Štvoruholník $APBR$ je konvexný, takže polpriamka $PR$ leží vnútri uhla $APB$; z~rovnosti oblúkov $AR$ a~$RB$ je $|\angle APR| = |\angle RPB| = \varphi$ a~$|RA| = |RB|$. Analogicky polpriamka $PQ$ leží vnútri uhla $APC$, $|\angle QPC| = 90^\circ - \varphi$ a~$|QA| = |QC|$. Polpriamky $PB$ a~$PC$ sú opačné a~polpriamka $PA$ delí polrovinu s~hraničnou priamkou $BC$ obsahujúcu $A$ na uhly $APB$ a~$APC$, vnútri ktorých postupne ležia polpriamky $PR$ a~$PQ$. Preto
    $$
    |\angle RPQ| = 180^\circ - \varphi - (90^\circ - \varphi) = 90^\circ.
    $$
    Bod $P$ tak leží na kružnici s~priemerom $QR$ a~zostáva dokázať, že na nej leží aj $M$, čiže že $|\angle QMR| = 90^\circ$.

    Bod $M$ je zatiaľ len stredom úsečky a~okolo stredu úsečky niet uhlov, ktoré by sa dali počítať. Dáme mu lepší význam tým, čo si stred úsečky vždy pýta: nech $R_0$ je obraz bodu $R$ v~stredovej súmernosti so stredom $M$, takže $BRCR_0$ je rovnobežník. Keďže $R$ neleží na priamke $BC$, je $R_0 \neq R$, a~keďže na nej neleží ani $Q$, je $Q \neq M$. Bod $M$ je stredom úsečky $RR_0$, takže dokazovaná rovnosť $|\angle QMR| = 90^\circ$ hovorí presne to, že $Q$ leží na osi úsečky $RR_0$, čiže
    $$
    |QR| = |QR_0|.
    $$

    Teraz sa oplatí vyznačiť rovnako dlhé úsečky. Jednu dvojicu dáva $Q$, totiž $|QA| = |QC|$. Druhú dáva $R$: stredová súmernosť podľa $M$ zobrazuje $B$ na $C$ a~$R$ na $R_0$, takže $|CR_0| = |BR| = |AR|$. Trojuholníky $QAR$ a~$QCR_0$ tak majú zhodné dve dvojice strán a~tretiu dvojicu $QR$, $QR_0$ dostaneme podľa vety sus, len čo overíme
    $$
    |\angle RAQ| = |\angle R_0CQ|.
    $$

    Pri tom budeme opakovane používať nasledujúce pozorovanie. Ak body $X$, $Z$ ležia v~opačných polrovinách určených priamkou $VY$, leží vnútri uhla $XVZ$ práve jedna z~dvoch polpriamok tejto priamky s~počiatkom $V$. V~prvom prípade je $|\angle XVZ| = |\angle XVY| + |\angle YVZ|$, v~druhom $|\angle XVZ| = 360^\circ - |\angle XVY| - |\angle YVZ|$, a~keďže $|\angle XVZ| \leq 180^\circ$, rozhoduje o~tom, ktorý prípad nastal, práve súčet $|\angle XVY| + |\angle YVZ|$.

    Obvodové uhly nad tetivou $RB$ z~toho istého oblúka dávajú $|\angle RAB| = |\angle RPB| = \varphi$, analogicky $|\angle QAC| = |\angle QPC| = 90^\circ - \varphi$. Body $R$, $C$ ležia v~opačných polrovinách určených priamkou $AB$ a~$\varphi + \alpha < 180^\circ$, takže
    $$
    |\angle RAC| = |\angle RAB| + |\angle BAC| = \varphi + \alpha
    $$
    a~polpriamka $AB$ leží vnútri uhla $RAC$. Body $R$, $B$ preto ležia v~tej istej polrovine určenej priamkou $AC$, a~keďže $Q$, $B$ ležia v~opačných, oddeľuje priamka $AC$ body $R$ a~$Q$. Súčet $(\varphi + \alpha) + (90^\circ - \varphi) = 90^\circ + \alpha$ je menší než $180^\circ$, čiže
    $$
    |\angle RAQ| = |\angle RAC| + |\angle CAQ| = 90^\circ + \alpha.
    $$

    Pri vrchole $C$ počítajme rovnako. Z~$|QA| = |QC|$ je $|\angle QCA| = |\angle QAC| = 90^\circ - \varphi$, priamka $AC$ oddeľuje $Q$ od $B$ a~$(90^\circ - \varphi) + \gamma < 180^\circ$, takže $|\angle QCB| = 90^\circ - \varphi + \gamma$. Z~$|RA| = |RB|$ je $|\angle RBA| = |\angle RAB| = \varphi$, priamka $AB$ oddeľuje $R$ od $C$ a~$\varphi + \beta < 180^\circ$, takže $|\angle RBC| = \varphi + \beta$. Stredová súmernosť podľa $M$ zobrazuje polpriamku $BC$ na polpriamku $CB$ a~polpriamku $BR$ na polpriamku $CR_0$, preto $|\angle R_0CB| = |\angle RBC| = \varphi + \beta$.

    Tá istá súmernosť zobrazuje polrovinu určenú priamkou $BC$ a~obsahujúcu $A$ na tú druhú, takže $R_0$ a~$Q$ ležia v~opačných polrovinách určených priamkou $BC$. Tentoraz je
    $$
    |\angle QCB| + |\angle BCR_0| = 90^\circ + \beta + \gamma = 270^\circ - \alpha,
    $$
    čo je vďaka $\alpha < 90^\circ$ viac než $180^\circ$. Vnútri uhla $QCR_0$ teda leží polpriamka opačná k~polpriamke $CB$ a~
    $$
    |\angle QCR_0| = 360^\circ - (270^\circ - \alpha) = 90^\circ + \alpha.
    $$

    Trojuholníky $QAR$ a~$QCR_0$ sa preto zhodujú: $|QA| = |QC|$, $|AR| = |CR_0|$ a~uhly medzi týmito stranami sú oba $90^\circ + \alpha$. Odtiaľ $|QR| = |QR_0|$, čiže $|\angle QMR| = 90^\circ$ a~aj bod $M$ leží na kružnici s~priemerom $QR$.

    Ostáva overiť, že ide o~štyri rôzne body. Z~$|\angle RPQ| = 90^\circ$ sú $Q$, $P$, $R$ navzájom rôzne, bod $M$ je rôzny od $P$ podľa zadania a~od $Q$ i~$R$ preto, že tie neležia na priamke $BC$. Body $Q$, $M$, $P$, $R$ tak ležia na jednej kružnici, a~to na kružnici s~priemerom $QR$.
}

\EnvId{aq53bc6NlKQWV07NwXRnF-dotycnica-a-obraz-bodu-r}
\Problem{0}{IMO 2017, P4}{
    Nech~$R$ a~$S$ sú rôzne body na kružnici~$k$ a~nech~$t$ označuje dotyčnicu ku~$k$ v~bode~$R$. Nech~$R_0$ je zrkadlový obraz~$R$ podľa~$S$. Bod~$I$ je zvolený na kratšom oblúku~$RS$ kružnice~$k$ tak, že kružnica~$k_0$ opísaná trojuholníku~$ISR_0$ pretína~$t$ v~dvoch rôznych bodoch. Nech~$A$ je spoločný bod~$k_0$ a~$t$ bližší k~$R$. Priamka~$AI$ pretína~$k$ znova v~$J$. Dokážte, že $JR_0$ je dotyčnicou ku~$k_0$.
}{
    Bod~$S$ je stred nejakej úsečky a~to robí počítanie uhlov neuskutočniteľným. Napriek tomu sa niečo počítaním uhlov dá, čo môže naznačiť, ako pomocou stredu~$S$ definovať nový bod.
}{
    Kľúčom je objaviť $AR_0 \parallel RJ$. To dáva celkom jasný náznak: obraz~$J$ v~stredovej súmernosti podľa~$S$ dopadne na~$AR_0$. Odtiaľto to dorazíme, na počítanie uhlov máme všetko potrebné.
}{
    V~celom riešení počítame s~orientovanými uhlami priamok modulo $180^\circ$: symbol $\angle(p, q)$ označuje uhol, o~ktorý treba otočiť priamku $p$, aby bola rovnobežná s~priamkou $q$. Konfigurácia sa totiž mení podľa polohy $I$ na oblúku aj podľa toho, ktorý z~priesečníkov priamky $t$ a~kružnice $k_0$ je $A$, a~orientované uhly nás týchto rozborov zbavia. Použijeme dve ich štandardné vlastnosti: štyri navzájom rôzne body $P$, $Q$, $X$, $Y$ ležia na jednej kružnici alebo priamke práve vtedy, keď $\angle(XP, XQ) = \angle(YP, YQ)$, a~pre dotyčnicu $\tau$ kružnice v~jej bode $P$ a~ďalšie jej body $Q$, $X$ platí $\angle(\tau, PQ) = \angle(XP, XQ)$.

    Odbavme si najprv, že všetky body, s~ktorými pracujeme, sú navzájom rôzne. Kružnica $k_0$ existuje, takže $I$ neleží na priamke $SR_0$, čiže ani na priamke $RS$; zvlášť $I \neq R$. Priamka $t$ má s~$k$ jediný spoločný bod $R$, takže na nej neležia $I$ ani $S$; priamka $RS$ je preto od $t$ rôzna a~pretínajú sa jedine v~$R$, čiže na $t$ neleží ani $R_0$. Bod $A$ je tak rôzny od $I$, $S$ i~$R_0$. Ďalej $R$ neleží na $k_0$, lebo inak by $k_0$ mala s~priamkou $RS$ tri rôzne spoločné body $R$, $S$, $R_0$; preto $A \neq R$. Bod $R_0$ neleží na $k$, lebo stred tetivy $RR_0$ by ležal vnútri $k$, kým $S$ leží na $k$; preto $J \neq R_0$. Priamka $AI$ pretína $k_0$ práve v~bodoch $A$ a~$I$, takže na nej neleží $S$ a~máme $J \neq S$; keby na nej ležal $R$, splynula by s~$t$ (obe by prechádzali rôznymi bodmi $A$ a~$R$) a~bod $I$ by ležal na $t$, takže aj $J \neq R$. Napokon zo zadania má priamka $AI$ s~$k$ druhý priesečník, čiže $J \neq I$.

    Bod $S$ je stredom úsečky $RR_0$ a~s~takouto podmienkou uhly nič nezmôžu. Priamka $R_0S$ je však tá istá ako $RS$, a~na jedno pozorovanie to stačí. Body $A$, $I$, $S$, $R_0$ ležia na $k_0$, body $R$, $S$, $I$, $J$ ležia na $k$ a~body $A$, $I$, $J$ ležia na jednej priamke, takže
    $$
    \gather
    \angle(R_0A, RS) = \angle(R_0A, R_0S) = \angle(IA, IS) \\
    = \angle(IJ, IS) = \angle(RJ, RS),
    \endgather
    $$
    a~teda
    $$
    AR_0 \parallel RJ.
    $$

    Rovnobežnosť rovno hovorí, ako stred $S$ využiť: stredová súmernosť $\sigma$ so stredom $S$ zobrazí $R$ na $R_0$, takže obraz $J' = \sigma(J)$ leží na priamke vedenej bodom $R_0$ rovnobežne s~$RJ$, čo je práve priamka $AR_0$. Ide o~známu situáciu, keď stred úsečky v~zadaní volá po stredovej súmernosti, ktorá vyrobí rovnobežník $RJR_0J'$. Body $A$, $R_0$, $J'$ teda ležia na jednej priamke.

    Ukážeme, že $A$, $R$, $S$, $J'$ ležia na jednej kružnici. Priamka $J'S$ je tá istá ako $JS$, lebo $S$ je stredom úsečky $JJ'$ a~$J \neq S$. Spolu s~$AR_0 \parallel RJ$ a~s~vetou o~úsekovom uhle pre dotyčnicu $t$ a~tetivu $RS$ kružnice $k$ dostaneme
    $$
    \angle(J'A, J'S) = \angle(JR, JS) = \angle(t, RS) = \angle(RA, RS).
    $$
    Naše štyri body sú navzájom rôzne, netriviálne sú len $J' \neq R$ (lebo $J \neq R_0$), $J' \neq S$ (lebo $J \neq S$) a~$J' \neq A$ (inak by $S$ ako stred úsečky $JJ'$ ležal na priamke $AJ$, čiže na priamke $AI$). Na jednej priamke neležia, lebo inak by $A$ ležal na priamke $RS$, a~teda by bol $A = R$. Ležia teda na jednej kružnici.

    Zostáva poznatok preniesť späť k~$R_0$. Súmernosť $\sigma$ zachováva orientované uhly, zobrazuje $J$ na $J'$, $R_0$ na $R$ a~priamku $R_0S$, čiže priamku $RR_0$ prechádzajúcu jej stredom, samu na seba. Preto
    $$
    \gather
    \angle(JR_0, R_0S) = \angle(J'R, RS) \\
    = \angle(AJ', AS) = \angle(AR_0, AS),
    \endgather
    $$
    kde druhá rovnosť plynie z~tetivovosti bodov $A$, $R$, $S$, $J'$ a~tretia z~kolinearity bodov $A$, $R_0$, $J'$. Body $A$, $S$, $R_0$ sú tri rôzne body kružnice $k_0$, takže pravá strana je uhol, ktorý s~tetivou $R_0S$ zviera dotyčnica $k_0$ v~bode $R_0$. Priamka prechádzajúca bodom $R_0$ je pritom svojím orientovaným uhlom s~$R_0S$ určená jednoznačne, a~preto je $JR_0$ touto dotyčnicou.
}

\EnvId{ZuNLEoEwr-J_fVOQWtSbR-dva-priemery-a-ortocentrum}
\Problem{0}{IMO 2013, P4}{
    Nech~$ABC$ je ostrouhlý trojuholník s~ortocentrom~$H$ a~nech~$W$ je bod na strane~$BC$. Označme~$M$ a~$N$ päty výšok z~$B$ a~$C$ v~trojuholníku~$ABC$. Nech~$k_1$ je kružnica opísaná trojuholníku~$BWN$ a~nech~$X$ je bod na~$k_1$ taký, že~$XW$ je priemer~$k_1$. Analogicky nech~$k_2$ je kružnica opísaná trojuholníku~$CWM$ a~nech~$Y$ je bod na~$k_2$ taký, že~$YW$ je priemer~$k_2$. Dokážte, že $X$, $Y$, $H$ ležia na jednej priamke.
}{
    Úloha nás takmer pozýva definovať druhý priesečník kružníc~$k_1$ a~$k_2$. To ich prirodzene poprepletá a~vieme tiež dobre využiť ich priemery.
}{
    Ak~$T$ je druhý priesečník~$k_1$ a~$k_2$, tak priemery nám dajú pravé uhly, z~ktorých~$T$, $X$, $Y$ ležia na jednej priamke. Dôkaz, že~$H$ tiež leží na tejto priamke, sa dá preformulovať ako kolmosť. Neeliminovali sme náhodou~$X$ a~$Y$?
}{
    Keďže je trojuholník $ABC$ ostrouhlý, päty $N$, $M$ ležia vnútri strán $AB$, $AC$ a~ortocentrum $H$ vo vnútri trojuholníka, konkrétne vnútri úsečky $AA'$, kde $A'$ je päta výšky z~$A$. Aby kružnice $k_1$, $k_2$ existovali, musí byť bod $W$ rôzny od $B$ i~od $C$; sú to potom dve rôzne kružnice, lebo tri rôzne kolineárne body $B$, $W$, $C$ na jednej kružnici ležať nemôžu.

    Vzájomné polohy bodov v~obrázku závisia od toho, kde na strane $BC$ bod $W$ leží, no na podstate argumentu to nič nemení. Počítame preto v~orientovaných uhloch modulo $180^\circ$: $\angle(p,q)$ značí uhol, o~ktorý treba otočiť priamku $p$, aby bola rovnobežná s~$q$.

    Kružnice $k_1$ a~$k_2$ prechádzajú spoločným bodom $W$, a~to je presne situácia, keď sa oplatí pridať ich druhý priesečník; označme ho $T$ a~jeho existenciu dokážme o~chvíľu. Vďaka nemu sa dajú využiť priemery zo zadania: z~priemeru $WX$ a~Tálesovej vety je $|\angle WTX| = 90^\circ$ (pre $X = T$ triviálne), takže $X$ leží na priamke $\ell$ vedenej bodom $T$ kolmo na $TW$, a~rovnako z~priemeru $WY$ leží na $\ell$ aj $Y$. Body $X$ a~$Y$ sme tým z~úlohy odstránili a~zostáva dokázať, že na $\ell$ leží aj $H$, čo pre $T \neq H$ znamená
    $$
    TH \perp TW.
    $$
    (Ak $T = H$, sme hotoví, lebo $T$ na $\ell$ leží.)

    Kľúčovou kružnicou je $\omega$ nad priemerom $AH$. Bod $H$ leží na výške $CN$ kolmej na $AB$ a~na výške $BM$ kolmej na $AC$, takže $|\angle ANH| = |\angle AMH| = 90^\circ$ a~okrem $A$ a~$H$ ležia na $\omega$ aj päty $N$ a~$M$. Priamku $BC$ pritom $\omega$ nepretína: body $A$ i~$H$ sa premietajú do $A'$, takže stred $\omega$, čiže stred úsečky $AH$, má od $BC$ vzdialenosť $(|AA'| + |HA'|)/2$, kým polomer $\omega$ je $|AH|/2 = (|AA'| - |HA'|)/2$, čo je menej.

    Teraz existencia bodu $T$. Keby sa $k_1$ a~$k_2$ v~bode $W$ dotýkali, mali by tam spoločnú dotyčnicu $t$ a~veta o~úsekovom uhle v~oboch kružniciach by dala
    $$
    \gather
    \angle(WN, WM) = \angle(WN, t) + \angle(t, WM) \\
    = \angle(BN, BW) + \angle(CW, CM) \\
    = \angle(AB, BC) + \angle(BC, AC) = \angle(AN, AM).
    \endgather
    $$
    Bod $W$ by teda ležal na kružnici opísanej trojuholníku $ANM$, čiže na $\omega$, ktorá priamku $BC$ nepretína. To je spor, kružnice $k_1$ a~$k_2$ majú okrem $W$ ešte druhý spoločný bod $T$.

    Ten istý výpočet, len s~obvodovými uhlami namiesto úsekového, ukáže, že $T$ leží na $\omega$. Najprv $T \neq N$: inak by na $k_2$ ležali body $C$, $M$, $N$, ktoré ležia aj na kružnici $\gamma$ nad priemerom $BC$ (veď $|\angle BNC| = |\angle BMC| = 90^\circ$), takže by bolo $k_2 = \gamma$ a~bod $W$ by ako spoločný bod $\gamma$ a~priamky $BC$ splynul s~$B$ alebo $C$. Rovnako $T \neq M$. Obvodové uhly v~$k_1$ a~v~$k_2$ tak dávajú
    $$
    \gather
    \angle(TN, TM) = \angle(TN, TW) + \angle(TW, TM) \\
    = \angle(BN, BW) + \angle(CW, CM) = \angle(AN, AM),
    \endgather
    $$
    čiže body $A$, $N$, $M$, $T$ ležia na jednej kružnici a~$T$ leží na $\omega$.

    Zostáva kolmosť; nech teda $T \neq H$. Body $T$, $B$, $W$, $N$ ležia na $k_1$ a~body $T$, $A$, $N$, $H$ na $\omega$, takže
    $$
    \gather
    \angle(TW, TH) = \angle(TW, TN) + \angle(TN, TH) \\
    = \angle(BW, BN) + \angle(AN, AH).
    \endgather
    $$
    Prvý sčítanec je $\angle(BC, AB)$ a~druhý $\angle(AB, AH)$, čo je kvôli $AH \perp BC$ presne $\angle(AB, BC) + 90^\circ$. Súčet je teda $90^\circ$, čiže $TH \perp TW$. Bod $H$ preto leží na priamke $\ell$, na ktorej ležia aj $X$ a~$Y$.
}

\EnvId{FhO0eO_oi-IndJmQ9N6Ql-dva-body-na-strane-bc}
\Problem{0}{IMO 2014, P4}{
    Body~$P$ a~$Q$ sú zvolené na strane~$BC$ ostrouhlého trojuholníka $ABC$ tak, že $|\angle PAB| = |\angle ACB|$ a~$|\angle QAC| = |\angle CBA|$. Body~$M$ a~$N$ ležia postupne na polpriamkach~$AP$ a~$AQ$ tak, že $|AP| = |PM|$ a~$|AQ| = |QN|$. Dokážte, že priamky~$BM$ a~$CN$ sa pretínajú na kružnici opísanej trojuholníku~$ABC$.
}{
    Body~$Q$ a~$P$ sú stredy~$AN$ a~$AM$. Tento fakt sa sám o~sebe využíva ťažko. Kľúčom je definovať nejaké ďalšie stredy. V~najlepšom prípade nám tieto stredy pomôžu eliminovať nie práve prirodzene vyzerajúce body~$M$ a~$N$.
}{
    Definujte stredy~$J$ a~$K$ strán~$AC$ a~$AB$. Stredné priečky nám dajú $BM \parallel KP$ a~$CN \parallel QJ$. Takže treba ukázať, že uhol medzi úsečkami~$QJ$ a~$KP$ je $180^\circ - \alpha$. Zdá sa, že sme nepokročili, keďže uhly pri ťažniciach sú vo všeobecnosti škaredé. V~konfigurácii je však množstvo podobných trojuholníkov.
}{
    Označme $\alpha = |\angle BAC|$, $\beta = |\angle CBA|$, $\gamma = |\angle ACB|$. Rovnosti $|AP| = |PM|$ a~$|AQ| = |QN|$ hovoria, že $P$ je stred úsečky $AM$ a~$Q$ je stred úsečky $AN$.

    Najprv preložme uhlové podmienky do dĺžok. Bod $P$ leží na strane $BC$, takže $|\angle ABP| = \beta = |\angle CBA|$, a~zo zadania $|\angle BAP| = \gamma = |\angle BCA|$; podľa vety $uu$ je $\triangle BAP \sim \triangle BCA$ (v~poradí vrcholov $B \leftrightarrow B$, $A \leftrightarrow C$, $P \leftrightarrow A$). Analogicky $\triangle CAQ \sim \triangle CBA$, a~z~pomerov strán
    $$
    \frac{|BP|}{|BA|} = \frac{|BA|}{|BC|}, \qquad \frac{|CQ|}{|CA|} = \frac{|CA|}{|CB|}. \eqno(*)
    $$

    Body $M$ a~$N$ vstupujú do zadania jedine cez stredy $P$ a~$Q$ a~samy o~sebe sa uchopiť nedajú. Keď sú v~úlohe stredy úsečiek, oplatí sa doplniť ďalšie stredy, aby vznikli stredné priečky: nech $K$, $J$, $L$ sú postupne stredy strán $AB$, $AC$, $BC$. V~trojuholníku $ABM$ sú $K$ a~$P$ stredy strán $AB$ a~$AM$, takže $KP \parallel BM$; v~trojuholníku $ACN$ sú $J$ a~$Q$ stredy strán $AC$ a~$AN$, takže $JQ \parallel CN$. Body $M$, $N$ tým z~obrázka zmizli: uhol priamok $BM$ a~$CN$ je ten istý ako uhol priamok $KP$ a~$JQ$ a~stačí ukázať, že je rovný $180^\circ - \alpha$.

    Obe priamky pretínajú $BC$, zmerajme teda ich uhly voči nej. Keďže $|BK| = \tfrac12|AB|$ a~$|BL| = \tfrac12|BC|$, dáva $(*)$
    $$
    \frac{|BK|}{|BL|} = \frac{|AB|}{|BC|} = \frac{|BP|}{|BA|},
    $$
    pričom $\angle KBP$ a~$\angle LBA$ je ten istý uhol $\beta$ (bod $K$ leží na úsečke $AB$ a~body $P$, $L$ na úsečke $BC$). Podľa vety $sus$ je $\triangle BKP \sim \triangle BLA$ (v~poradí vrcholov $B \leftrightarrow B$, $K \leftrightarrow L$, $P \leftrightarrow A$), a~teda $|\angle BPK| = |\angle BAL|$. Rovnako z~$|CJ| : |CL| = |AC| : |BC| = |CQ| : |CA|$ a~zo spoločného uhla pri $C$ je $\triangle CJQ \sim \triangle CLA$ a~$|\angle CQJ| = |\angle CAL|$. Uhly pri ťažnici $AL$ sú síce škaredé, no potrebujeme len ich súčet: bod $L$ leží vnútri strany $BC$, takže polpriamka $AL$ leží vnútri uhla $BAC$ a
    $$
    |\angle BPK| + |\angle CQJ| = |\angle BAL| + |\angle LAC| = \alpha.
    $$

    Zostáva preniesť tieto dva uhly k~vrcholom $B$ a~$C$. Bod $M$ je obrazom bodu $A$ v~stredovej súmernosti podľa bodu $P$ ležiaceho na $BC$, leží teda v~polrovine $\rho$ s~hraničnou priamkou $BC$ neobsahujúcej $A$; to isté platí pre $N$. Naopak $K$ a~$J$ ležia vnútri úsečiek $AB$ a~$AC$, čiže mimo $\rho$. Priamka $BC$ je priečkou rovnobežiek $KP$, $BM$, body $K$, $M$ ležia pri nej na opačných stranách a~polpriamky $PB$, $BC$ mieria opačne (bod $P$ leží na $BC$ a~$P \neq B$), takže ide o~striedavé uhly; rovnako pre rovnobežky $JQ$, $CN$. Dostávame
    $$
    |\angle MBC| = |\angle BPK|, \qquad |\angle NCB| = |\angle CQJ|.
    $$

    Polpriamky $BM$ a~$CN$ tak vychádzajú z~koncov úsečky $BC$ do tej istej polroviny $\rho$ a~zvierajú s~$BC$ uhly, ktorých súčet je $\alpha < 180^\circ$. Pretnú sa preto v~bode $X$ ležiacom v~$\rho$ a~zo súčtu uhlov v~trojuholníku $BXC$ je
    $$
    |\angle BXC| = 180^\circ - |\angle MBC| - |\angle NCB| = 180^\circ - \alpha.
    $$
    Body $A$ a~$X$ ležia v~opačných polrovinách určených priamkou $BC$ a~$|\angle BXC| + |\angle BAC| = 180^\circ$, takže $ABXC$ je tetivový štvoruholník. Priesečník $X$ priamok $BM$ a~$CN$ preto leží na kružnici opísanej trojuholníku $ABC$.
}

\EnvId{3AK1oHdV1pF38vTXfQc41-stredy-kruznic-pri-patach-vysok}
\Problem{0}{ISL 2012, G3}{
    Nech $ABC$ je ostrouhlý trojuholník. Označme $D$, $E$, $F$ päty výšok z~$A$, $B$, $C$. Nech~$I_1$ a~$I_2$ sú stredy kružníc vpísaných trojuholníkom~$BFD$ a~$CED$ a~nech~$O_1$ a~$O_2$ sú stredy kružníc opísaných trojuholníkom~$AI_1B$ a~$AI_2C$. Dokážte, že $I_1I_2 \parallel O_1O_2$.
}{
    Body~$O_1$ a~$O_2$ vyzerajú naozaj umelo. Naozaj také aj sú. Oba sa dajú eliminovať definovaním nového bodu a~preformulovaním rovnobežnosti na niečo iné. Aj bez týchto zvláštnych bodov v~konfigurácii platí viac, než úloha tvrdí, a~nie je to celkom ľahko vidieť. Oplatí sa to uhádnuť.
}{
    Nech~$T$ je druhý priesečník našich dvoch kružníc so stredmi~$O_1$ a~$O_2$. Chceme ukázať, že~$AT$ je kolmé na~$I_1I_2$. To stále nie je zadarmo. Pomáha uhádnuť a~dokázať, že $BCI_2I_1$ je tetivový štvoruholník.
}{
    Označme $\alpha$, $\beta$, $\gamma$ vnútorné uhly trojuholníka $ABC$, ďalej $I$ stred jeho vpísanej kružnice a~$r$ jej polomer. Kružnice opísané trojuholníkom $AI_1B$ a~$AI_2C$, teda kružnice so stredmi $O_1$ a~$O_2$, označme $\omega_1$ a~$\omega_2$.

    Body $O_1$, $O_2$ vstupujú do tvrdenia jedine cez priamku $O_1O_2$, a~tá je osou spoločnej tetivy $AT$, kde $T$ je druhý spoločný bod kružníc $\omega_1$, $\omega_2$. Stačí teda dokázať, že $AT \perp I_1I_2$; existenciu bodu $T$ dokážeme o~chvíľu a~ukáže sa pritom, že priamka $AT$ nie je nič iné ako os uhla $\angle BAC$. Z~presného obrázka sa dá uhádnuť, že v~konfigurácii platí viac, než úloha tvrdí: $BCI_2I_1$ je tetivový štvoruholník. To je jadro riešenia a~začneme ním.

    Trojuholník $ABC$ je ostrouhlý, takže $D$, $E$, $F$ ležia postupne vnútri strán $BC$, $CA$, $AB$. Uhly $\angle AFC$ a~$\angle ADC$ sú pravé, takže $A$, $F$, $D$, $C$ ležia na kružnici s~priemerom $AC$ a~mocnosť bodu $B$ vzhľadom na ňu dáva $|BF| \cdot |BA| = |BD| \cdot |BC|$, čiže
    $$
    \frac{|BF|}{|BC|} = \frac{|BD|}{|BA|}.
    $$
    Tento spoločný pomer označme $k$; keďže $BD$ je odvesna a~$BA$ prepona pravouhlého trojuholníka $ABD$, je $0 < k < 1$. Uhol $\angle FBD$ je pritom ten istý ako $\angle CBA$, takže trojuholníky $BFD$ a~$BCA$ sa zhodujú v~uhle pri $B$ aj v~pomere jeho ramien, čiže sú podobné s~koeficientom $k$, pričom $F \leftrightarrow C$ a~$D \leftrightarrow A$.

    Bod $I_1$ je stred vpísanej kružnice trojuholníka $BFD$, takže priamka $BI_1$ je osou uhla $\angle FBD$, čiže osou uhla $\angle ABC$, a~$I_1$ leží na polpriamke $BI$. Podobnosť zobrazuje stred vpísanej kružnice na stred vpísanej kružnice, takže $|BI_1| = k|BI|$, a~z~$0 < k < 1$ vyplýva, že $I_1$ leží vnútri úsečky $BI$.

    Súčin $|IB| \cdot |II_1|$ odmeriame pomocou obrazu stredu $I$ v~osovej súmernosti. Nech $X$ je dotykový bod vpísanej kružnice so stranou $BC$ a~nech $I'$ je obraz bodu $I$ v~súmernosti podľa priamky $BC$. Bod $X$ je pätou kolmice z~$I$ na $BC$, teda stredom úsečky $II'$; preto $|IX| = r$, $|II'| = 2r$ a~polpriamky $IX$, $II'$ splývajú. Ďalej $|BI'| = |BI|$ a~$|\angle IBI'| = 2|\angle IBC| = \beta$.

    Nech $Q$ je päta kolmice z~$I'$ na priamku $BI$. Uhol $\angle IBI'$ je nenulový a~ostrý, takže $I'$ neleží na priamke $BI$ a~$Q$ leží na polpriamke $BI$. Pravouhlé trojuholníky $BQI'$ a~$BDA$ majú pravé uhly pri $Q$ a~$D$ a~zhodujú sa v~uhle pri $B$, lebo $|\angle QBI'| = \beta = |\angle DBA|$. Sú teda podobné a
    $$
    |BQ| = |BI'| \cdot \frac{|BD|}{|BA|} = k|BI| = |BI_1|,
    $$
    a~keďže $Q$ i~$I_1$ ležia na polpriamke $BI$, je $Q = I_1$. Zvlášť $|\angle II_1I'| = 90^\circ$.

    Bod $I_1$ leží vnútri úsečky $BI$, takže polpriamky $II_1$ a~$IB$ splývajú, a~ako sme videli, splývajú aj polpriamky $II'$ a~$IX$. Pravouhlé trojuholníky $II_1I'$ a~$IXB$ (s~pravými uhlami pri $I_1$ a~$X$) sa tak zhodujú v~uhle pri $I$, čiže sú podobné, pričom $I_1 \leftrightarrow X$ a~$I' \leftrightarrow B$. Odtiaľ
    $$
    |IB| \cdot |II_1| = |IX| \cdot |II'| = 2r^2.
    $$

    Pri vrchole $C$ je situácia rovnaká. Uhly $\angle AEB$ a~$\angle ADB$ sú pravé, takže $A$, $E$, $D$, $B$ ležia na kružnici s~priemerom $AB$ a~mocnosť bodu $C$ dáva $|CE| \cdot |CA| = |CD| \cdot |CB|$, čiže
    $$
    \frac{|CE|}{|CB|} = \frac{|CD|}{|CA|} = k',
    $$
    kde $0 < k' < 1$, lebo $CD$ je odvesna a~$CA$ prepona pravouhlého trojuholníka $ACD$. Trojuholníky $CED$ a~$CBA$ sú preto podobné s~koeficientom $k'$ (s~$E \leftrightarrow B$ a~$D \leftrightarrow A$), lebo uhol $\angle ECD$ je ten istý ako $\angle BCA$. Priamka $CI_2$ je preto osou uhla $\angle ACB$, bod $I_2$ leží vnútri úsečky $CI$ a~$|CI_2| = k'|CI|$. Ten istý bod $I'$ spĺňa $|CI'| = |CI|$ a~$|\angle ICI'| = 2|\angle ICB| = \gamma$; päta kolmice z~$I'$ na priamku $CI$ tak leží na polpriamke $CI$ a~z~podobnosti s~pravouhlým trojuholníkom $CDA$ je jej vzdialenosť od $C$ rovná $k'|CI|$, čiže je to bod $I_2$. Tá istá dvojica podobných pravouhlých trojuholníkov ako pri $B$ potom dáva $|IC| \cdot |II_2| = |IX| \cdot |II'| = 2r^2$.

    Priamky $BI$ a~$CI$ sú rôzne a~pretínajú sa jedine v~$I$, takže body $B$, $I_1$, $I_2$, $C$ (žiadny z~nich nie je $I$) sú navzájom rôzne a~neležia na jednej priamke. Bod $I$ leží mimo úsečky $BI_1$ i~mimo úsečky $CI_2$, takže z~rovnosti $|IB| \cdot |II_1| = |IC| \cdot |II_2|$ podľa kritéria tetivovosti ležia $B$, $I_1$, $I_2$, $C$ na jednej kružnici.

    Tá istá rovnosť nám dá aj bod $T$. Body $B$, $I_1$ ležia na $\omega_1$ a~bod $I$ leží na priamke $BI_1$ mimo úsečky $BI_1$, takže $I$ leží zvonka $\omega_1$ a~jeho mocnosť vzhľadom na $\omega_1$ je $|IB| \cdot |II_1| = 2r^2$; rovnako mocnosť $I$ vzhľadom na $\omega_2$ je $|IC| \cdot |II_2| = 2r^2$. Bod $I$ teda leží na chordále kružníc $\omega_1$, $\omega_2$, na ktorej leží aj ich spoločný bod $A$.

    Nech $Z$ je dotykový bod vpísanej kružnice so stranou $AB$. Trojuholník $AZI$ má pravý uhol pri $Z$, uhol $\frac\alpha2$ pri $A$ a~preponu $AI$. Z~$\alpha < 90^\circ$ je uhol pri $A$ menší než $45^\circ$, takže uhol pri $I$ je väčší než $45^\circ$; proti väčšiemu uhlu leží väčšia strana, čiže $|AZ| > |ZI| = r$ a~podľa Pytagorovej vety
    $$
    |IA|^2 = r^2 + |AZ|^2 > 2r^2.
    $$
    Keby sa priamka $AI$ dotýkala $\omega_1$, dotýkala by sa jej v~bode $A$ a~mocnosť $I$ by bola $|IA|^2 \neq 2r^2$. Priamka $AI$ teda pretína $\omega_1$ ešte v~bode $T_1 \neq A$, pričom $I$ leží zvonka $\omega_1$, takže $A$ a~$T_1$ ležia na tej istej polpriamke z~$I$ a~$|IT_1| = 2r^2/|IA|$. To isté s~$\omega_2$ dáva bod $T_2$ na tej istej polpriamke s~$|IT_2| = 2r^2/|IA|$, čiže $T_1 = T_2$. Kružnice $\omega_1$, $\omega_2$ majú teda okrem $A$ ešte druhý spoločný bod $T = T_1$ a~ten leží na priamke $AI$.

    Bod $I_1$ leží vnútri úsečky $BI$, teda vo vnútri trojuholníka $ABC$. Keby $\omega_1 = \omega_2$, prechádzala by táto kružnica bodmi $A$, $B$, $C$, bola by to kružnica opísaná trojuholníku $ABC$ a~bod $I_1$ by na nej ležal, nie vo vnútri. Kružnice $\omega_1$, $\omega_2$ sú teda rôzne a~majú spoločný bod, takže nie sú sústredné a~$O_1 \neq O_2$. Oba stredy sú rovnako vzdialené od $A$ i~od $T$, takže priamka $O_1O_2$ je osou úsečky $AT$, a~keďže priamka $AT$ je priamka $AI$, dostávame $O_1O_2 \perp AI$.

    Zostáva dokázať, že $I_1I_2 \perp AI$. Nech $I_a$ je stred kružnice pripísanej oproti vrcholu $A$. Leží na osi uhla $\angle BAC$ na opačnej strane priamky $BC$ než $A$, takže $I$ leží vnútri úsečky $AI_a$. Priamky $BI$, $BI_a$ sú vnútorná a~vonkajšia os uhla pri $B$, takže $|\angle IBI_a| = 90^\circ$, a~rovnako $|\angle ICI_a| = 90^\circ$. Podľa vety o~základných uhloch (po premenovaní vrcholov) je
    $$
    |\angle AIB| = 90^\circ + \tfrac\gamma2, \qquad |\angle AIC| = 90^\circ + \tfrac\beta2.
    $$

    Nech $J$ je päta kolmice z~$I_1$ na priamku $AI$. Bod $I_1$ leží na priamke $BI$ a~je rôzny od $I$, takže na priamke $AI$ neleží a~uhly trojuholníka $IJI_1$ pri $I$ a~$I_1$ sú ostré (zvlášť $J \neq I$). Keďže $I_1$ leží na polpriamke $IB$, je $|\angle AII_1| = |\angle AIB| > 90^\circ$, takže $J$ neleží na polpriamke $IA$, ale na polpriamke k~nej opačnej, čiže na polpriamke $II_a$. Uhly $\angle JII_1$ a~$\angle I_aIB$ sú tak ten istý uhol a~pravouhlé trojuholníky $IJI_1$, $IBI_a$ sú podobné (s~$J \leftrightarrow B$, $I_1 \leftrightarrow I_a$), odkiaľ
    $$
    |IJ| \cdot |II_a| = |II_1| \cdot |IB| = 2r^2.
    $$
    Rovnako päta $J'$ kolmice z~$I_2$ na $AI$ leží na polpriamke $II_a$, lebo $|\angle AII_2| = |\angle AIC| > 90^\circ$, a~z~podobnosti trojuholníkov $IJ'I_2$, $ICI_a$ dostaneme
    $$
    |IJ'| \cdot |II_a| = |II_2| \cdot |IC| = 2r^2.
    $$

    Body $J$, $J'$ teda ležia na tej istej polpriamke z~$I$ a~$|IJ| = |IJ'|$, čiže $J = J'$. Priamky $I_1J$ a~$I_2J$ sú potom obe kolmicou na $AI$ v~bode $J$, a~preto splývajú. Body $I_1$, $J$, $I_2$ tak ležia na jednej priamke a~$I_1I_2 \perp AI$. Spolu s~$O_1O_2 \perp AI$ to dáva $I_1I_2 \parallel O_1O_2$.
}

\EnvId{a0vidDqjikMyjuIYb5vVc-stvoruholnik-z-osi-uhlov}
\Problem{0}{MEMO 2012, T6}{
    Nech~$ABCD$ je konvexný štvoruholník, ktorého žiadne dve strany nie sú rovnobežné, a~$|\angle ABC| = |\angle CDA|$. Predpokladajme, že vzájomné priesečníky osí uhlov pri susedných vrcholoch~$ABCD$ tvoria štvoruholník~$EFGH$. Nech~$K$ je priesečník uhlopriečok~$EFGH$. Dokážte, že priesečník priamok~$AB$ a~$CD$ leží na kružnici opísanej trojuholníku~$BKD$.
}{
    Úlohu je trochu neohrabané kresliť. Nedajte sa zastrašiť. Neštandardná uhlová podmienka vyzerá, akoby mohla znamenať niečo pekné. Aj znamená: treba len niečo pretnúť. Čo sa týka osí uhlov, nebojte sa ich; dve často implikujú tretiu a~tu ich je habadej.
}{
    Nech~$E$ leží na osiach uhlov pri~$A$ a~$B$, $F$ na osiach pri~$B$ a~$C$, $G$ na osiach pri~$C$ a~$D$ a~$H$ na osiach pri~$D$ a~$A$. Všetky body~$E$, $F$, $G$, $H$ sú stredy vpísaných alebo pripísaných kružníc nejakých trojuholníkov vďaka osiam uhlov. V~každom prípade body~$P$, $F$, $H$ a $Q$, $E$, $G$ ležia na jednej priamke a~tieto priamky sú osami uhlov~$\angle CPB$ a~$\angle CQD$. Sme už veľmi blízko.
}{
    Stačí pretnúť vhodné priamky. Okrem priesečníka $P$ priamok $AB$ a~$CD$ zo zadania pridáme priesečník $Q$ priamok $BC$ a~$AD$; obidva existujú, lebo žiadne dve strany nie sú rovnobežné. Až táto dvojica dá podmienke $|\angle ABC| = |\angle CDA|$ tvar, s~ktorým sa dá pracovať.

    Označme $\alpha = |\angle DAB|$, $\beta = |\angle ABC|$, $\gamma = |\angle BCD|$ a~$\delta = |\angle CDA|$, takže $\beta = \delta$ a~$\alpha + 2\beta + \gamma = 360^\circ$.

    Priamka každej strany konvexného štvoruholníka je opornou priamkou a~zvyšné dva vrcholy ležia striktne v~jednej polrovine. Body $A$, $B$ teda ležia striktne na tej istej strane priamky $CD$, čiže $P$ neleží na úsečke $AB$ a~polpriamky $PA$, $PB$ splývajú; symetricky splývajú $PC$ s~$PD$, $QB$ s~$QC$ a~$QA$ s~$QD$. Uhly $\angle BPC$, $\angle APD$ a~$\angle BPD$ sú preto ten istý uhol, rovnako ako $\angle AQB$, $\angle CQD$ a~$\angle BQD$.

    Priamka $BC$ je priečkou priamok $AB$ a~$CD$ a~body $A$, $D$ ležia na tej istej jej strane, takže $AB \parallel CD$ práve vtedy, keď $\beta + \gamma = 180^\circ$; to je vylúčené. Zámena označenia $(A,B,C,D) \to (C,D,A,B)$ je cyklickým posunom: zachová predpoklady aj body $P$ a~$Q$, zamení $E$ s~$G$ a~$F$ s~$H$ (uhlopriečky $EFGH$, a~teda ani bod $K$, sa nezmenia) a~zamení $B$ s~$D$, pričom $\beta + \gamma$ prejde na $\alpha + \beta = 360^\circ - \beta - \gamma$. Dokazované tvrdenie je voči nej invariantné, takže smieme predpokladať $\beta + \gamma < 180^\circ$.

    Body $B$, $C$ ležia striktne na tej istej strane priamky $AD$, zatiaľ čo $A$ a~$D$ ležia na nej. Preto $A$ leží medzi $P$ a~$B$ práve vtedy, keď sú $P$ a~$B$ v~opačných polrovinách, čo nastane práve vtedy, keď sú v~opačných polrovinách $P$ a~$C$, čiže práve vtedy, keď $D$ leží medzi $P$ a~$C$. Uhly trojuholníka $PBC$ pri vrcholoch $B$ a~$C$ sú tak buď $\beta$ a~$\gamma$, alebo $180^\circ - \beta$ a~$180^\circ - \gamma$; druhá možnosť dáva súčet $360^\circ - \beta - \gamma > 180^\circ$, takže platí prvá a~bod $A$ leží medzi $P$ a~$B$ a~bod $D$ medzi $P$ a~$C$. Tá istá úvaha s~priamkou $AB$ dáva, že $B$ leží medzi $Q$ a~$C$ práve vtedy, keď $A$ leží medzi $Q$ a~$D$, a~uhly trojuholníka $QCD$ pri $C$ a~$D$ sú $\gamma$ a~$\delta$, lebo $360^\circ - \gamma - \delta = 360^\circ - \gamma - \beta > 180^\circ$. Odtiaľ
    $$
    |\angle BPC| = 180^\circ - \beta - \gamma = 180^\circ - \gamma - \delta = |\angle CQD|.
    $$
    To je celý obsah podmienky $\beta = \delta$: uhly pri $P$ a~$Q$ sú zhodné.

    Teraz osi uhlov. V~trojuholníku $PBC$ sú osi štvoruholníka pri $B$ a~$C$ jeho vnútornými osami, takže $F$ je stred jeho vpísanej kružnice. V~trojuholníku $PAD$ je $|\angle PAD| = 180^\circ - \alpha$ a~$|\angle PDA| = 180^\circ - \delta$, takže osi štvoruholníka pri $A$ a~$D$ sú jeho vonkajšími osami a~$H$ je stred kružnice pripísanej oproti vrcholu $P$. Oba body preto ležia na vnútornej osi $p$ uhla $\angle BPC$. Analogicky je $G$ stred vpísanej kružnice trojuholníka $QCD$ a~$E$ stred kružnice pripísanej trojuholníku $QAB$ oproti vrcholu $Q$, takže oba ležia na vnútornej osi $q$ uhla $\angle CQD$. Uhlopriečkami štvoruholníka $EFGH$ sú teda priamky $EG = q$ a~$FH = p$ a~bod $K$ je priesečníkom priamok $p$ a~$q$.

    Zo zistených polpriamok: $BP$ je polpriamka $BA$ a~$BQ$ je opačná k~$BC$, takže $|\angle PBQ| = 180^\circ - \beta$; podobne $DP$ je opačná k~$DC$ a~$DQ$ je polpriamka $DA$, takže $|\angle PDQ| = 180^\circ - \delta$. Tieto dva uhly sú zhodné. Bod $A$ leží striktne medzi $P$ a~$B$ i~striktne medzi $Q$ a~$D$, pričom neleží na priamke $PQ$ (inak by priamky $AB$ a~$AD$ splynuli). Bod $B$ leží na polpriamke z~bodu $P$ cez $A$ za bodom $A$, a~keďže $P$ leží na priamke $PQ$, je $B$ striktne v~tej istej polrovine ako $A$; rovnako $D$ leží na polpriamke z~bodu $Q$ cez $A$ za bodom $A$, čiže tiež v~tejto polrovine. Body $B$ a~$D$ tak vidia úsečku $PQ$ pod rovnakým uhlom a~z~tej istej strany, takže $P$, $B$, $D$, $Q$ ležia na jednej kružnici $\omega$.

    Keďže $B$ a~$D$ ležia v~tej istej polrovine s~hraničnou priamkou $PQ$, tetivy $BD$ a~$PQ$ sa nepretínajú, a~teda $P$ i~$Q$ ležia na tom istom z~oblúkov určených bodmi $B$ a~$D$. Priamka $p$ je vnútornou osou uhla $\angle BPD$, takže jej druhý priesečník $M$ s~$\omega$ delí oblúk $BD$ neobsahujúci $P$ na dva zhodné oblúky, lebo obvodové uhly $\angle BPM$ a~$\angle MPD$ sú zhodné. Bod $Q$ leží na tom istom oblúku ako $P$, takže ani jeden z~oblúkov $BM$, $MD$ neobsahuje $Q$; obvodové uhly nad nimi dávajú $|\angle BQM| = |\angle MQD|$ a~z~polohy $M$ na oblúku $BD$ neobsahujúcom $Q$ leží polpriamka $QM$ vnútri uhla $\angle BQD$. Priamka $QM$ je teda vnútornou osou uhla $\angle BQD$, čiže $QM = q$.

    Bod $M$ leží na $p$ i~na $q$, preto $M = K$. Body $B$, $K$, $D$ sú tri rôzne body kružnice $\omega$, takže $\omega$ je kružnica opísaná trojuholníku $BKD$, a~leží na nej aj bod $P$.
}

\EnvId{DzF1H7zPMIba3UVbwJoNr-bod-na-kruznici-apob}
\Problem{0}{MEMO 2016, I3}{
    Nech~$ABC$ je ostrouhlý trojuholník so stredom opísanej kružnice~$O$ a~$|\angle BAC| > 45^\circ$. Bod~$P$ leží v~jeho vnútri tak, že $A$, $P$, $O$, $B$ ležia na jednej kružnici a~priamka~$BP$ je kolmá na~$CP$. Bod~$Q$ leží na úsečke~$BP$ tak, že~$AQ$ a~$PO$ sú rovnobežné. Dokážte, že $|\angle QCB| = |\angle PCO|$.
}{
    Je zjavné, že v~úlohe niečo chýba. Oplatí sa preto najprv získať lepší obraz situácie preformulovaním tvrdenia na rovnosť nejakých iných uhlov, ideálne takých, ktoré sa dajú vyjadriť pomocou základných uhlov~$ABC$. To môže inšpirovať pridanie. Nezabúdajme, že úloha stále obsahuje známy indikátor potenciálne dobrých pridaní: pravý uhol.
}{
    Tvrdenie je ekvivalentné s~$|\angle PCQ|=|\angle OCB|$, čo je ekvivalentné s~$|\angle BAC|=|\angle CQP|$. Takže ak zrkadlíme~$Q$ podľa~$P$ do~$Q'$, stačí dokázať, že body ležia na jednej kružnici. Ani to nie je zjavné, no v~tomto bode sa to dá dotiahnuť.
}{
    Označme $\alpha = |\angle BAC|$ a~$\gamma = |\angle BCA|$, nech $\omega$ je kružnica prechádzajúca bodmi $A$, $P$, $O$, $B$ a~nech $k$ je kružnica s~priemerom $BC$. Z~$|\angle BPC| = 90^\circ$ leží $P$ na $k$. Trojuholník $ABC$ je ostrouhlý, takže $O$ leží v~jeho vnútri a~podľa tvrdenia o~základných uhloch je $|\angle BOC| = 2\alpha$ a~$|\angle AOB| = 2\gamma$. Z~rovnoramenných trojuholníkov $OBC$ a~$OAB$ tak dostávame
    $$
    |\angle OCB| = 90^\circ - \alpha, \qquad |\angle OAB| = 90^\circ - \gamma.
    $$

    Jediný raz využijeme predpoklad $\alpha > 45^\circ$, a~to hneď teraz: dáva $|\angle BOC| = 2\alpha > 90^\circ$, čiže $O$ leží vnútri kružnice $k$, kým z~$\alpha < 90^\circ$ leží $A$ mimo $k$. Kružnica $\omega$ prechádza vnútorným bodom $O$ kružnice $k$, takže $\omega \neq k$ a~tieto dve kružnice majú spoločné práve body $B$ a~$P$ (tie sú rôzne, lebo $P$ leží vo vnútri trojuholníka). Delia $\omega$ na dva oblúky, z~ktorých jeden leží vnútri $k$ a~druhý mimo $k$, takže $A$ a~$O$ ležia na rôznych oblúkoch: na $\omega$ idú naše štyri body v~poradí $A$, $P$, $O$, $B$. Špeciálne $P \neq O$, takže priamka $PO$ je vôbec dobre definovaná; a~keďže $A$, $P$, $O$ sú tri rôzne body kružnice, nie sú kolineárne, čiže z~$AQ \parallel PO$ máme $Q \neq P$.

    Dokážeme, že $|\angle PCQ| = 90^\circ - \alpha$, čo už stačí. Bod $Q$ totiž leží na úsečke $BP$, takže polpriamka $CQ$ leží v~uhle $BCP$ a
    $$
    |\angle BCP| = |\angle BCQ| + |\angle QCP| \geq |\angle QCP|.
    $$
    Ak teda $|\angle QCP| = 90^\circ - \alpha$, je $|\angle BCO| = 90^\circ - \alpha \leq |\angle BCP|$, a~keďže $O$ aj $P$ ležia vo vnútri trojuholníka, teda v~tej istej polrovine s~hraničnou priamkou $BC$, leží v~uhle $BCP$ aj polpriamka $CO$ a
    $$
    |\angle BCP| = |\angle BCO| + |\angle OCP|.
    $$
    Porovnaním oboch rozkladov uhla $BCP$ a~vykrátením zhodných uhlov $|\angle QCP| = |\angle BCO|$ dostaneme $|\angle BCQ| = |\angle OCP|$, čo je dokazovaná rovnosť.

    Keďže $Q \neq P$ leží na úsečke $BP$, splývajú polpriamky $PQ$ a~$PB$, a~tak $|\angle QPC| = |\angle BPC| = 90^\circ$. V~pravouhlom trojuholníku $QPC$ je preto rovnosť $|\angle PCQ| = 90^\circ - \alpha$ ekvivalentná s~$|\angle CQP| = \alpha$; stačí teda dokázať $|\angle CQP| = |\angle BAC|$.

    Pravý uhol pri $P$ si pýta doplnenie pravouhlého trojuholníka $QPC$ na rovnoramenný: nech $Q'$ je obraz $Q$ v~stredovej súmernosti podľa $P$. Priamka $CP$ je osou úsečky $QQ'$, takže $|CQ| = |CQ'|$ a~$|\angle CQP| = |\angle CQ'P|$. Bod $Q$ leží medzi $B$ a~$P$ a~bod $P$ medzi $Q$ a~$Q'$, čiže na polpriamke $BP$ idú body v~poradí $Q$, $P$, $Q'$; polpriamky $Q'P$ a~$Q'B$ preto splývajú a~$|\angle CQ'P| = |\angle CQ'B|$. Priamky $AB$ a~$BC$ prechádzajú bodom $B$, takže každý bod polpriamky $BP$ okrem $B$ leží pri nich v~tej istej polrovine ako vnútorný bod $P$: bod $Q'$ tak leží v~tej istej polrovine ako $C$ vzhľadom na $AB$ a~v~tej istej ako $A$ vzhľadom na $BC$. Stačí preto dokázať, že $Q'$ leží na kružnici opísanej trojuholníku $ABC$: z~obvodových uhlov nad tetivou $BC$ potom $|\angle CQ'B| = |\angle CAB| = \alpha$, a~teda $|\angle CQP| = \alpha$.

    Z~poradia bodov na $\omega$ tetiva $AB$ neoddeľuje $P$ od $O$ a~tetiva $OB$ neoddeľuje $P$ od $A$, takže z~obvodových uhlov
    $$
    \gather
    |\angle APB| = |\angle AOB| = 2\gamma, \\
    |\angle OPB| = |\angle OAB| = 90^\circ - \gamma.
    \endgather
    $$
    Polpriamky $PQ$ a~$PB$ splývajú, čiže $|\angle APQ| = 2\gamma$. Priamka $BP$ je priečkou rovnobežiek $QA$ a~$PO$ a~body $A$, $O$ ležia pri nej v~opačných polrovinách, takže ide o~striedavé uhly a~$|\angle AQP| = |\angle OPB| = 90^\circ - \gamma$. Zo súčtu uhlov v~trojuholníku $APQ$ je
    $$
    |\angle QAP| = 180^\circ - 2\gamma - (90^\circ - \gamma) = 90^\circ - \gamma = |\angle AQP|,
    $$
    takže trojuholník $APQ$ je rovnoramenný a~$|PA| = |PQ| = |PQ'|$.

    Rovnoramenný je tým pádom aj trojuholník $APQ'$, pričom $P$ leží medzi $Q$ a~$Q'$, takže jeho uhol pri $P$ je $|\angle APQ'| = 180^\circ - |\angle APQ| = 180^\circ - 2\gamma$. Zvyšné dva uhly sú zhodné, čiže $|\angle AQ'P| = \gamma$. Polpriamky $Q'P$ a~$Q'B$ splývajú, a~tak $|\angle AQ'B| = \gamma = |\angle ACB|$. Body $Q'$ a~$C$ ležia v~tej istej polrovine s~hraničnou priamkou $AB$, preto $Q'$ leží na oblúku $AB$ kružnice opísanej trojuholníku $ABC$ obsahujúcom $C$.
}

\EnvId{GK8F64_JP-rSXryIXL_wp-patuholnik-a-dva-stredy-kruznic}
\Problem{0}{MEMO 2017, I3}{
    V~konvexnom päťuholníku $ABCDE$ označme~$P$ priesečník priamok~$CE$ a~$BD$. Dokážte, že ak $|\angle PAD| = |\angle ACB|$ a~$|\angle CAP| = |\angle EDA|$, tak~$P$ leží na priamke určenej stredmi kružníc opísaných trojuholníkom~$ABC$ a~$ADE$.
}{
    Keď nevieme, čo robiť, chvíľu počítajme uhly. To by nám malo umožniť nájsť v~obrázku jednoduchú peknú vlastnosť, ktorá môže inšpirovať, čo pridať.
}{
    Kľúčom je objaviť $BC \parallel DE$. Obrázok obsahuje priesečník~$CE$ a~$BD$. To môže inšpirovať pridanie obrazu~$A$ v~rovnoľahlosti zobrazujúcej~$BC$ na~$DE$. Zrazu sa všetko stane krásnym.
}{
    Označme $\omega_1$ a~$\omega_2$ postupne kružnice opísané trojuholníkom $ABC$ a~$ADE$ a~ich stredy $O_1$, $O_2$.

    Najprv si ujasnime polohu bodu $P$. Uhlopriečka $AC$ oddeľuje $B$ od bodov $D$, $E$ a~uhlopriečka $AD$ oddeľuje $E$ od bodov $B$, $C$. Bod $P$ je vnútorným bodom úsečky $CE$, leží teda v~otvorenej polrovine určenej priamkou $AC$ a~bodom $E$, a~súčasne je vnútorným bodom úsečky $BD$, takže leží v~otvorenej polrovine určenej priamkou $AD$ a~bodom $B$. Leží preto vnútri uhla $\angle CAD$ a~platí
    $$
    |\angle CAP| + |\angle PAD| = |\angle CAD|.
    $$

    Uhlopriečky $CA$ a~$DA$ ležia vnútri vnútorných uhlov päťuholníka pri vrcholoch $C$ a~$D$, takže
    $$
    \gather
    |\angle BCD| = |\angle BCA| + |\angle ACD|, \\
    |\angle CDE| = |\angle CDA| + |\angle ADE|.
    \endgather
    $$
    Podľa zadania je $|\angle BCA| = |\angle PAD|$ a~$|\angle ADE| = |\angle CAP|$, zo súčtu uhlov v~trojuholníku $ACD$ je $|\angle ACD| + |\angle CDA| = 180^\circ - |\angle CAD|$, a~teda
    $$
    |\angle BCD| + |\angle CDE| = |\angle CAD| + 180^\circ - |\angle CAD| = 180^\circ.
    $$
    Body $B$ a~$E$ ležia v~tej istej polrovine s~hraničnou priamkou $CD$, takže $\angle BCD$ a~$\angle CDE$ sú vnútorné uhly pri priečke $CD$ na tej istej strane. Ich súčet $180^\circ$ dáva
    $$
    BC \parallel DE.
    $$

    Dve rovnobežné úsečky sú pozvánkou uvažovať o~rovnoľahlosti, ktorá jednu zobrazí na druhú a~ktorou vieme preniesť ďalší bod. Jej stred v~obrázku už máme: priamky $BD$ a~$CE$ sa pretínajú práve v~$P$. Štvoruholník $BCDE$ je konvexný, takže $P$ leží vnútri oboch úsečiek $BD$ aj $CE$, a~rovnoľahlosť $h$ so stredom $P$ a~koeficientom $k = -|PD| \, / \, |PB| < 0$ zobrazuje $B$ na bod $D$. Priamku $BC$ zobrazí na rovnobežku s~ňou vedenú bodom $D$, čiže na priamku $DE$, a~priamku $CE$ na seba samu, lebo tá prechádza stredom $P$. Preto $h(C) = E$.

    Pridaným bodom je obraz $A' = h(A)$. Ak leží na $\omega_2$, sme hotoví: $\omega_2$ je potom kružnica opísaná trojuholníku $A'DE$, teda obrazu trojuholníka $ABC$ v~$h$, čiže $h(\omega_1) = \omega_2$, a~rovnoľahlosť zobrazuje stred na stred.

    Keďže $k < 0$, leží $P$ medzi $A$ a~$A'$ a~polpriamky $A'A$, $A'P$ splývajú. Rovnoľahlosť $h$ zobrazuje trojuholník $APC$ na trojuholník $A'PE$ (bod $P$ je jej pevným bodom), a~tak
    $$
    |\angle AA'E| = |\angle PA'E| = |\angle PAC| = |\angle ADE|,
    $$
    kde posledná rovnosť je druhá podmienka zo zadania.

    Ostáva overiť polohu bodu $A'$. Bod $P$ je vnútorným bodom uhlopriečky $BD$, leží teda vo vnútri päťuholníka, čiže v~otvorenej polrovine s~hraničnou priamkou $AE$, v~ktorej leží aj vrchol $D$. Bod $A'$ leží na polpriamke $AP$ a~je rôzny od $A$, takže leží v~tej istej polrovine. Body $D$ a~$A'$ preto vidia úsečku $AE$ pod rovnakým uhlom a~ležia pri nej z~tej istej strany, takže $A$, $E$, $D$, $A'$ ležia na jednej kružnici. Tou je $\omega_2$, a~teda $A'$ na $\omega_2$ leží.

    Trojuholník $A'DE$ je obrazom trojuholníka $ABC$, jeho vrcholy teda nie sú kolineárne a~kružnica $h(\omega_1)$ je nimi jednoznačne určená. Všetky tri ležia na $\omega_2$, preto $h(\omega_1) = \omega_2$ a~$h(O_1) = O_2$. Ak $O_1 = O_2$, je $O_1$ pevným bodom rovnoľahlosti s~koeficientom $k \neq 1$, čiže $O_1 = O_2 = P$ a~priamka zo zadania nie je určená. V~opačnom prípade je $O_1 \neq P$ a~body $P$, $O_1$, $O_2 = h(O_1)$ ležia z~definície rovnoľahlosti na jednej priamke.
}

\EnvId{NsM3t1SwEv7CD5wXQgwbU-usecka-pq-rovnobezna-s-bc}
\Problem{0}{IMO 2019, P2}{
    V~trojuholníku~$ABC$ nech~$B_1$ je bod na strane~$AC$ a~$C_1$ bod na strane~$AB$. Nech $P$ a~$Q$ sú body na úsečkách~$BB_1$ a~$CC_1$ tak, že priamka~$PQ$ je rovnobežná s~$BC$. Nech~$P_1$ je bod na priamke~$PC_1$ taký, že~$C_1$ leží striktne medzi~$P$ a~$P_1$ a~$|\angle PP_1A|=|\angle CBA|$. Nech~$Q_1$ je bod na priamke~$QB_1$ taký, že~$B_1$ leží striktne medzi~$Q$ a~$Q_1$ a~$|\angle AQ_1Q|=|\angle ACB|$. Dokážte, že $P$, $Q$, $P_1$, $Q_1$ ležia na jednej kružnici.
}{
    Cieľom je pridať niečo, čo objasní zvláštnu uhlovú podmienku pre~$P_1$ a~$Q_1$. Úloha má však priveľa bodov a~rôzne veci, ktoré vyzerajú sľubne, sa dajú pridať bez väčšieho úžitku, napríklad priesečník~$P_1P$ s~$BC$ alebo~$PQ$ s~$AB$ (skúste to). Problém s~týmito priesečníkmi je, že hoci dávajú tetivový štvoruholník, neposkytujú žiadne ďalšie užitočné uhly na prenesenie. Je preto lepšie pretnúť nejaké priamky s~kružnicou tak, aby to do rovnosti~$|\angle PP_1A|=|\angle CBA|$ prinieslo ešte jeden uhol.
}{
    Kľúčovým magickým bodom je druhý priesečník~$C_2$ priamky~$CC_1$ s~kružnicou opísanou trojuholníku~$ABC$. Vďaka nemu dostaneme $|\angle CBA|=|\angle CC_2A|$. Magicky tak vznikne tetivový štvoruholník~$C_1C_2P_1A$. Analogicky definujeme~$B_2$. A~verte či nie, teraz je to už len stredne ťažká úloha na počítanie uhlov.
}{
    Označme $\omega$ kružnicu opísanú trojuholníku $ABC$, $\beta = |\angle CBA|$ a~$\gamma = |\angle ACB|$.

    Uhol $\beta$ z~podmienky pre bod $P_1$ vidno v~obrázku pri vrchole $B$, teda ďaleko od bodov $C_1$ a~$P_1$, ktorých sa podmienka týka. Zároveň je to obvodový uhol nad tetivou $AC$, takže ho vidno z~každého bodu oblúka $AC$ obsahujúceho $B$. Stačí teda pretnúť s~$\omega$ vhodnú priamku. Tou pravou je $CC_1$: ležia na nej $C_1$ aj $Q$. Nech je teda $C_2$ jej druhý priesečník s~$\omega$ a~analogicky nech $B_2$ je druhý priesečník priamky $BB_1$ s~$\omega$.

    Body $B_1$, $C_1$ sú vnútorné body strán $AC$, $AB$ (pri $C_1 = A$ by $P_1$ ležal na priamke $PA$ a~uhol $|\angle PP_1A|$ by bol nulový alebo priamy, pri $C_1 = B$ by $Q$ ležal na priamke $BC$ a~priamka $PQ$ by s~ňou splynula; pre $B_1$ analogicky), takže ležia vnútri $\omega$. Ak by $P = B$, bola by priamka $PQ$ rovnobežkou s~$BC$ vedenou bodom $B$, čiže samotnou priamkou $BC$; preto $P \neq B$ a~rovnako $Q \neq C$. Úsečka $BB_1$ leží okrem koncového bodu $B$ vnútri $\omega$, takže vnútri $\omega$ leží $P$, a~rovnako $Q$.

    Priamka $CC_1$ pretína priamku $AB$ jedine v~jej vnútornom bode $C_1$, ktorý leží vnútri $\omega$; na priamke $CC_1$ je preto poradie bodov $C$, $C_1$, $C_2$ a~bod $C_2$ leží na otvorenom oblúku $AB$ neobsahujúcom $C$. Ten je časťou oblúka $AC$ obsahujúceho $B$, takže z~obvodových uhlov nad tetivou $AC$ dostávame
    $$
    |\angle AC_2C| = |\angle ABC| = \beta.
    $$

    Keďže $B_1$ je vnútorný bod strany $AC$, leží celá úsečka $BB_1$ okrem bodu $B$ v~otvorenej polrovine s~hraničnou priamkou $AB$ obsahujúcej $C$; špeciálne tam leží $P$. Bod $C_1$ leží na priamke $AB$ a~medzi $P$ a~$P_1$, takže $P_1$ leží v~opačnej otvorenej polrovine, a~v~tej leží aj $C_2$, lebo $C_1$ leží medzi $C$ a~$C_2$. Z~toho, že $C_1$ leží na polpriamke $P_1P$ aj na polpriamke $C_2C$, ďalej plynie
    $$
    |\angle AP_1C_1| = |\angle AP_1P| = \beta = |\angle AC_2C| = |\angle AC_2C_1|.
    $$
    Body $P_1$ a~$C_2$ teda vidia úsečku $AC_1$ pod rovnakým uhlom a~ležia na tej istej strane priamky $AC_1$, takže podľa obrátenej vety o~obvodovom uhle ležia $A$, $C_1$, $C_2$, $P_1$ na jednej kružnici. Analogicky, s~uhlom $\gamma$ namiesto $\beta$, ležia na jednej kružnici body $A$, $B_1$, $B_2$, $Q_1$, pričom $B_2$ leží na otvorenom oblúku $AC$ neobsahujúcom $B$.

    Bod $Q$ leží na úsečke $CC_1$, teda v~uzavretej polrovine priamky $AB$ obsahujúcej $C$; to isté platí pre $P$. Bod $P_1$ leží v~opačnej otvorenej polrovine, takže $P_1 \neq P$ i~$P_1 \neq Q$, a~analogicky $Q_1 \neq P$ i~$Q_1 \neq Q$.

    Ďalej už len počítame uhly. Aby sme sa vyhli rozboru vzájomných polôh, prejdeme na orientované uhly modulo $180^\circ$: $\angle(\ell,m)$ je uhol, o~ktorý treba otočiť priamku $\ell$, aby bola rovnobežná s~priamkou $m$. Použijeme dve štandardné vlastnosti: štyri navzájom rôzne body $X$, $Y$, $Z$, $W$ ležia na jednej kružnici alebo priamke práve vtedy, keď $\angle(XZ,XW) = \angle(YZ,YW)$, a~rovnobežné priamky zvierajú s~každou ďalšou priamkou ten istý orientovaný uhol.

    Priamky $BB_1$ a~$CC_1$ pretínajú priamku $BC$, takže nie sú rovnobežné s~$PQ$ a~s~priamkou $PQ$ majú spoločný jediný bod, postupne $P$ a~$Q$. Body $B_2$, $C_2$ ležia na $\omega$, kým $P$ a~$Q$ vnútri nej, takže ani $B_2$, ani $C_2$ neleží na priamke $PQ$; oblúky, na ktorých ležia, sú navyše disjunktné, takže $B_2 \neq C_2$. Označme $\Omega$ kružnicu prechádzajúcu bodmi $P$, $Q$, $C_2$. V~nasledujúcich krokoch je tak vždy možnosť \uv{alebo na priamke} vylúčená.

    Najprv ukážeme, že na $\Omega$ leží aj $B_2$. Bod $P$ leží na priamke $BB_2$, bod $Q$ na priamke $CC_2$ a~$PQ \parallel BC$, takže
    $$
    \eqalign{
        \angle(B_2P, B_2C_2) &= \angle(B_2B, B_2C_2) = \angle(CB, CC_2) \cr
        &= \angle(QP, QC_2),
    }
    $$
    kde prostredná rovnosť sú obvodové uhly nad tetivou $BC_2$ kružnice $\omega$. Body $P$, $Q$, $C_2$, $B_2$ teda ležia na $\Omega$.

    Teraz bod $P_1$; pre $P_1 = C_2$ niet čo dokazovať, nech je teda $P_1 \neq C_2$. Bod $C_1$ leží na priamke $P_1P$ aj na priamke $AB$ a~body $A$, $C_1$, $C_2$, $P_1$ ležia na jednej kružnici, takže
    $$
    \eqalign{
        \angle(P_1P, P_1C_2) &= \angle(P_1C_1, P_1C_2) = \angle(AC_1, AC_2) \cr
        &= \angle(AB, AC_2) = \angle(CB, CC_2) = \angle(QP, QC_2),
    }
    $$
    kde predposledná rovnosť sú obvodové uhly nad tetivou $BC_2$ kružnice $\omega$ a~posledná pochádza z~predošlého výpočtu. Bod $P_1$ preto leží na kružnici $\Omega$.

    S~bodom $Q_1$ naložíme symetricky, len namiesto $C_2$ použijeme $B_2$ (pre $Q_1 = B_2$ opäť niet čo dokazovať). Bod $B_1$ leží na priamke $Q_1Q$ aj na priamke $AC$ a~body $A$, $B_1$, $B_2$, $Q_1$ ležia na jednej kružnici, takže
    $$
    \eqalign{
        \angle(Q_1Q, Q_1B_2) &= \angle(Q_1B_1, Q_1B_2) = \angle(AB_1, AB_2) \cr
        &= \angle(AC, AB_2) = \angle(BC, BB_2) = \angle(PQ, PB_2),
    }
    $$
    kde predposledná rovnosť sú obvodové uhly nad tetivou $CB_2$ kružnice $\omega$ a~posledná využíva $PQ \parallel BC$ a~to, že $P$ leží na priamke $BB_2$. Bod $Q_1$ tak leží na kružnici prechádzajúcej bodmi $P$, $Q$, $B_2$, čo je $\Omega$.

    Na kružnici $\Omega$ teda ležia body $P$, $Q$, $B_2$, $C_2$, $P_1$ i~$Q_1$; špeciálne ležia na jednej kružnici body $P$, $Q$, $P_1$ a~$Q_1$.
}

\EnvId{sIGVitSRoddubnK736EUu-bod-na-vyske-pravouhleho-trojuholnika}
\Problem{0}{IMO 2012, P5}{
    V~pravouhlom trojuholníku~$ABC$ s~pravým uhlom pri vrchole~$A$ nech~$D$ je päta výšky z~$A$. Nech~$X$ je ľubovoľný vnútorný bod úsečky~$AD$. Nech~$K$ je bod na úsečke~$BX$ taký, že $|CK| = |CA|$. Podobne nech~$L$ je bod na úsečke~$CX$ taký, že $|BL| = |BA|$. Označme~$M$ priesečník priamok~$BL$ a~$CK$. Dokážte, že $|MK| = |ML|$.
}{
    V~konfigurácii nevidieť nič zmysluplné. Niečo je tam skryté, no nie je ľahké vidieť, čo presne. Jednou možnosťou je skúsiť urobiť niečo zmysluplné s~priamkami~$BX$ a~$CX$. Body~$K$ a~$L$ na nich ležia a~nevidíme, ako to využiť. Jednou možnosťou je pretnúť ich s~niečím zmysluplným.
}{
    Kľúčovými bodmi sú priesečníky~$E$ a~$F$ priamok~$BX$ a~$CX$ s~kružnicou nad priemerom~$BC$. Ani tie samy osebe nestačia, no ďalší bod je celkom prirodzený: priesečník~$G$ priamok~$BF$ a~$CE$. Zjavne leží na priamke~$AXD$ a~$X$ je ortocentrum trojuholníka~$GBC$. Hoci sa to tak nemusí zdať, úloha sa teraz dá dotiahnuť, a~to kombináciou mocnosti bodu, obvodových uhlov a~tetivových štvoruholníkov. Ani v~tomto štádiu to nie je ľahká úloha.
}{
    Z~predpokladov o~bodoch $K$ a~$L$ využijeme len toľko: $K$ leží na priamke $BX$ a~$|CK| = |CA|$, čiže leží na kružnici $\omega_c$ so stredom $C$ a~polomerom $|CA|$; podobne $L$ leží na priamke $CX$ a~na kružnici $\omega_b$ so stredom $B$ a~polomerom $|BA|$. Obe kružnice prechádzajú bodom $A$.

    Jediné, čo body $K$ a~$L$ v~obrázku drží, sú priamky $BX$ a~$CX$, a~samy osebe sa nedajú uchopiť. Pretnime ich preto s~kružnicou, ktorá o~konfigurácii niečo vie: s~Tálesovou kružnicou $\tau$ nad priemerom $BC$. Leží na nej aj bod $A$, lebo uhol pri $A$ je pravý.

    Keďže uhly pri $B$ a~$C$ sú ostré, bod $D$ leží vnútri úsečky $BC$; bod $X$ teda neleží na priamke $BC$ a~priamky $BX$, $CX$ sú rôzne. Ani jedna z~nich nie je kolmá na $BC$: kolmica na $BC$ vedená bodom $B$ je rovnobežná s~priamkou $AD$ a~od nej rôzna (bod $B$ na $AD$ neleží), takže bod $X$ neobsahuje, a~rovnako pre $C$. Priamka $BX$ preto nie je dotyčnicou $\tau$ v~bode $B$ a~pretína $\tau$ ešte v~bode $E \neq B$; pritom $E \neq C$, lebo inak by $X$ ležal na $BC$. Analogicky nech $F$ je druhý priesečník priamky $CX$ s~$\tau$, opäť $F \neq B$ a~$F \neq C$. Podľa Tálesovej vety je
    $$
    |\angle BEC| = |\angle BFC| = 90^\circ,
    $$
    čiže $CE \perp BX$ a~$BF \perp CX$.

    Priamka $CE$ je kolmica z~$C$ na $BX$, teda os tetivy, ktorú kružnica $\omega_c$ vytína na priamke $BX$; podobne $BF$ je os tetivy kružnice $\omega_b$ na priamke $CX$. Ich priesečník je prirodzený kandidát na stred kružnice prechádzajúcej bodmi $K$ aj $L$, a~to je bod, ktorý úlohu rozlomí. Označme ho $G$. Existuje a~je jediný, lebo priamky $BX$, $CX$ sú rôzne a~pretínajú sa v~$X$, takže nie sú rovnobežné, a~teda nie sú rovnobežné ani kolmice $CE$, $BF$ na ne.

    Body $G$, $B$, $C$ tvoria trojuholník. Bod $E$ leží na $\tau$ a~je rôzny od $B$ i~od $C$, takže neleží na priamke $BC$ a~priamky $CE$, $BC$ sa pretínajú jedine v~$C$; rovnako $F$ neleží na $BC$, čiže $C$ neleží na priamke $BF$. Keby $G$ ležal na $BC$, z~prvého by bolo $G = C$, čo druhé vylučuje.

    Priamka $BX$ prechádza vrcholom $B$ a~je kolmá na priamku $CG = CE$, je to teda výška trojuholníka $GBC$ z~vrcholu $B$; rovnako $CX$ je jeho výška z~vrcholu $C$. Bod $X$ je preto ortocentrum trojuholníka $GBC$ a~tretia výška dáva $GX \perp BC$. Body $G$ a~$X$ sú pritom rôzne: keby $G = X$, ležal by $X$ na priamke $CE$ aj na priamke $BX$, tie sú na seba kolmé a~pretínajú sa jedine v~$E$, čiže by bolo $X = E$; lenže $A$ leží na $\tau$ a~$D$ vnútri nej, takže celá úsečka $AD$ okrem bodu $A$ leží vnútri $\tau$, kým $E$ leží na $\tau$. Priamky $GX$ a~$AD$ sú tak obe kolmé na $BC$ a~obe prechádzajú bodom $X$, čiže splývajú a~bod $G$ leží na priamke $AD$.

    Teraz už len počítame, a~to stále tým istým nástrojom: ak je priamka $p$ kolmá na priamku $q$ v~bode $T$, tak pre ľubovoľné body $Y$ na $p$ a~$Z$ na $q$ platí $|YZ|^2 = |YT|^2 + |TZ|^2$ (pre $Y = T$ alebo $Z = T$ triviálne). Použijeme to v~bodoch $D$, $E$, $F$ a~nakoniec v~$K$ a~$L$.

    Priamka $AD$ je kolmá na $BC$ v~bode $D$ a~ležia na nej body $A$ aj $G$. Rozdiely pre dvojice $(G,B)$, $(G,C)$, resp. $(A,B)$, $(A,C)$ preto dávajú
    $$
    \gather
    |GB|^2 - |GC|^2 = |DB|^2 - |DC|^2 = \\
    = |AB|^2 - |AC|^2.
    \endgather
    $$
    To je kľúčová identita a~má dve tváre. Po prvé, dá sa prepísať ako
    $$
    |GB|^2 - |BA|^2 = |GC|^2 - |CA|^2,
    $$
    čo hovorí, že mocnosť bodu $G$ ku kružnici $\omega_b$ sa rovná jeho mocnosti ku kružnici $\omega_c$, teda že $G$ leží na ich chordále (ktorou je práve priamka $AD$). Spoločnú hodnotu označme $k$. Po druhé, keďže uhol pri $A$ je pravý a~teda $|AB|^2 + |AC|^2 = |BC|^2$, dá sa prepísať aj dvoma tvarmi
    $$
    \gather
    |GB|^2 = |GC|^2 + |BC|^2 - 2|CA|^2, \\
    |GC|^2 = |GB|^2 + |BC|^2 - 2|BA|^2.
    \endgather
    $$

    Priamka $CE$ je kolmá na $BX$ v~bode $E$, pričom $C$ aj $G$ ležia na $CE$ a~$B$ aj $K$ na $BX$. Rozdiely pre dvojice $(G,K)$, $(C,K)$, resp. $(G,B)$, $(C,B)$ dávajú
    $$
    \gather
    |GK|^2 - |CK|^2 = |GE|^2 - |CE|^2 = \\
    = |GB|^2 - |CB|^2,
    \endgather
    $$
    lebo v~prvej dvojici sa odčíta $|EK|^2$ a~v~druhej $|EB|^2$. Dosadíme $|CK| = |CA|$ a~prvý z~dvoch tvarov kľúčovej identity:
    $$
    \gather
    |GK|^2 = |GB|^2 - |CB|^2 + |CA|^2 = \\
    = |GC|^2 - |CA|^2 = k.
    \endgather
    $$
    Bod $F$ dá to isté so zámenou $B \leftrightarrow C$, $K \leftrightarrow L$, $E \leftrightarrow F$, ktorá celú konfiguráciu zachováva: priamka $BF$ je kolmá na $CX$ v~bode $F$, na $BF$ ležia $B$ aj $G$, na $CX$ ležia $C$ aj $L$, a~zo vzťahu $|BL| = |BA|$ a~druhého tvaru kľúčovej identity vyjde
    $$
    |GL|^2 = |GC|^2 - |BC|^2 + |BA|^2 = k.
    $$
    Bod $G$ je teda od $K$ aj od $L$ vzdialený $\sqrt k$.

    Ak $k = 0$, je $K = G = L$ a~tvrdenie platí triviálne. Nech teda $k > 0$ a~nech $\gamma$ je kružnica so stredom $G$ a~polomerom $\sqrt k$; body $K$, $L$ na nej ležia. Z~$|GK|^2 = |GC|^2 - |CA|^2$ a~$|CK| = |CA|$ dostaneme
    $$
    |GK|^2 + |CK|^2 = |GC|^2,
    $$
    takže podľa obrátenej Pytagorovej vety je $|\angle GKC| = 90^\circ$ a~priamka $CK$ je dotyčnicou $\gamma$ v~bode $K$; rovnako je $BL$ dotyčnicou $\gamma$ v~bode $L$. Bod $M$ leží na oboch dotyčniciach, takže posledné použitie tej istej Pytagorovej vety, v~bode $K$ s~priamkami $GK \perp CK$ a~potom v~bode $L$ s~priamkami $GL \perp BL$, dáva
    $$
    |MK|^2 = |MG|^2 - k = |ML|^2.
    $$
}

\secc Moje obľúbené

Zopár konštrukcií, ktoré mi prídu obzvlášť rozkošné. Nie nutne najťažšie úlohy tu \SlightSmile{}

\EnvId{JLEqfNmB2Hw0P5oc4qFwQ-body-na-vyskach-a-obsahy}
\Problem{0}{}{
    Nech $ABC$ je ostrouhlý trojuholník. Body $X$, $Y$, $Z$ ležia postupne na výškach z $A$, $B$, $C$, vo vnútri trojuholníka, tak, že $A_{BCX} + A_{CAY} + A_{ABZ} = S_{ABC}$, kde $A_{PQR}$ označuje obsah trojuholníka $PQR$. Dokážte, že kružnica opísaná trojuholníku $XYZ$ prechádza ortocentrom trojuholníka $ABC$.
}{
    Keď máme súčet obsahov, je lepšie ho pekne interpretovať. Tieto trojuholníky sa však neohrabane prekrývajú. Tak radšej tieto obsahy \uv{presuňme} inam.
}{
    Trik je v~tom, že ak $X'$ je ľubovoľný bod na priamke prechádzajúcej bodom $X$ rovnobežne s $BC$, tak $A_{BCX}=A_{BCX'}$. Na čo nás toto navádza?
}{
    Očividne dáva zmysel uvažovať priamky rovnobežné so stranami, prechádzajúce našimi bodmi, a~pretnúť ich; zvoľme priamku cez $Y$ rovnobežnú s $AC$ a~priamku cez $Z$ rovnobežnú s $AB$ a~nechajme ich pretnúť v $T$. Čo hovorí obsahová podmienka teraz? Sme prakticky pri cieli.
}{
    Symbol $S_{ABC}$ zo zadania je obsah trojuholníka $ABC$, teda $A_{ABC}$; ďalej píšeme dôsledne $A_{PQR}$.

    Pre bod $P$ roviny označme $A_{BCP}$ orientovaný obsah trojuholníka $BCP$: obyčajný obsah so znamienkom $+$, ak $P$ leží v~tej istej polrovine s~hraničnou priamkou $BC$ ako bod $A$, so znamienkom $-$, ak leží v~opačnej, a~nulu, ak $P$ leží na priamke $BC$. Analogicky definujeme $A_{CAP}$ (znamienko podľa $B$) a~$A_{ABP}$ (podľa $C$). Body $X$, $Y$, $Z$ ležia vnútri trojuholníka, takže pre ne orientovaný obsah splýva s~obyčajným a~podmienka zo zadania sa nemení.

    Platí $A_{BCP} = \frac12 |BC| \cdot d(P)$, kde $d(P)$ je orientovaná vzdialenosť bodu $P$ od priamky $BC$. Pre dva body preto $A_{BCP} = A_{BCP'}$ práve vtedy, keď $P$ a~$P'$ ležia na tej istej rovnobežke s~priamkou $BC$; analogicky pre $CA$ a~$AB$. Ďalej pre každý bod $P$ platí
    $$
    A_{BCP} + A_{CAP} + A_{ABP} = A_{ABC}.
    $$
    Orientovaná vzdialenosť od priamky je afinná funkcia bodu, takže afinná je aj ľavá strana; v~bodoch $A$, $B$, $C$ nadobúda hodnotu $A_{ABC}$ a~afinná funkcia je svojimi hodnotami v~troch nekolineárnych bodoch určená jednoznačne.

    Tri trojuholníky z~podmienky sa neprehľadne prekrývajú, no obsah $A_{BCX}$ sa nezmení, keď $X$ posunieme po rovnobežke s~$BC$. Všetky tri obsahy teda vieme presunúť tak, aby prislúchali jedinému bodu: stačí pretnúť vhodné rovnobežky.

    Nech $T$ je priesečník priamky vedenej bodom $Y$ rovnobežne s~$CA$ a~priamky vedenej bodom $Z$ rovnobežne s~$AB$; priamky $CA$ a~$AB$ nie sú rovnobežné, takže $T$ existuje a~je jediný. Z~kritéria vyššie je $A_{CAT} = A_{CAY}$ a~$A_{ABT} = A_{ABZ}$, čiže
    $$
    \eqalign{
        A_{BCT} &= A_{ABC} - A_{CAT} - A_{ABT} \cr
        &= A_{ABC} - A_{CAY} - A_{ABZ} = A_{BCX},\cr
    }
    $$
    kde posledná rovnosť je práve podmienka zo zadania. Body $X$ a~$T$ tak ležia na tej istej rovnobežke s~$BC$.

    Označme $H$ ortocentrum trojuholníka $ABC$. Body $X$, $H$ ležia na výške z~$A$, ktorá je kolmá na $BC$; podobne $Y$, $H$ ležia na výške z~$B$ kolmej na $CA$ a~$Z$, $H$ na výške z~$C$ kolmej na $AB$.

    Keby bolo $T = H$, ležal by bod $H$ s~bodom $X$ na tej istej rovnobežke s~$BC$ a~zároveň na výške z~$A$; tieto dve priamky sú na seba kolmé, takže majú jediný spoločný bod a~$X = H$. Rovnako by bolo $Y = H$ a~$Z = H$, čo je spor s~tým, že $XYZ$ je trojuholník. Preto $T \neq H$ a~kružnica $k$ s~priemerom $HT$ je dobre definovaná.

    Bod $X$ leží na $k$: pre $X = H$ aj $X = T$ je to zrejmé, inak je priamka $XH$ výškou z~$A$ a~priamka $XT$ rovnobežkou s~$BC$, tie sú na seba kolmé, takže $|\angle HXT| = 90^\circ$ a~$X$ leží na $k$ podľa Tálesovej vety. Rovnako ležia na $k$ body $Y$ (výška z~$B$ je kolmá na $CA$) a~$Z$ (výška z~$C$ je kolmá na $AB$).

    Kružnica $k$ prechádza bodmi $X$, $Y$, $Z$, a~keďže tie tvoria trojuholník, je práve $k$ jeho opísanou kružnicou. Tá teda prechádza bodom $H$.
}

\EnvId{V8qpCa8IyPgddyS2NJS8Y-trojuholnik-s-obvodom-styri}
\Problem{0}{}{
    Nech $ABC$ je trojuholník s~obvodom $4$. Body $P$ a $Q$ ležia postupne na polpriamkach $AB$ a $AC$ tak, že $|AP| = |AQ| = 1$ a~úsečky $BC$ a $PQ$ sa pretínajú v~bode $X$. Dokážte, že jeden z~dvoch trojuholníkov, na ktoré $AX$ delí trojuholník $ABC$, má obvod $2$.
}{
    Prvý trik je spomenúť si na jedno obzvlášť pekné miesto, kde sa v~geometrii trojuholníka vyskytuje obvod (alebo azda jeho polovica)?
}{
    Vzdialenosť z $A$ do dotykových bodov kružnice pripísanej oproti vrcholu $A$ je presne polovica obvodu, teda 2. A~$|AP|=|AQ|=1$. Je to náhoda?
}{
    Nie je to náhoda (inak by som sa nepýtal). Priamka $DE$ je obrazom $A$-poláry kružnice pripísanej oproti vrcholu $A$ v~rovnoľahlosti so stredom $A$ a~koeficientom $1/2$. Čo o~tejto priamke vieme?
}{
    Nuž, táto priamka je chordálou $A$ a~kružnice pripísanej oproti vrcholu $A$! A~$F$ na nej leží. Blížime sa.
}{
    Posledným bodom na nakreslenie je dotykový bod pripísanej kružnice a $BC$. Využitím faktu, že $F$ leží na tejto chordále, dostaneme nejaké pekné rovnosti dĺžok.
}{
    Obvod je 4, čiže polobvod je 2, a~$|AP| = |AQ| = 1$ je presne jeho polovica. Polobvod sa v~geometrii trojuholníka objavuje ako vzdialenosť z~vrcholu do dotykových bodov kružnice pripísanej oproti tomuto vrcholu, tak si ju do obrázka pridajme.

    Označme $a = |BC|$, $b = |CA|$, $c = |AB|$, takže $a + b + c = 4$, a~nech $\omega$ je kružnica pripísaná trojuholníku $ABC$ oproti vrcholu $A$. Tá sa dotýka strany $BC$ v~bode $U$ a~predĺžení strán $AB$, $AC$ za bodmi $B$, $C$ postupne v~bodoch $V$, $W$. Mocnosť bodu $M$ vzhľadom na kružnicu $\gamma$ označme $p(M,\gamma)$.

    Z~rovnosti dĺžok dotyčníc z~jedného bodu je $|BV| = |BU|$, $|CW| = |CU|$ a~$|AV| = |AW|$, a~keďže $U$ leží na úsečke $BC$, dostávame
    $$
    |AV| + |AW| = (c + |BU|) + (b + |CU|) = a + b + c = 4.
    $$
    Preto $|AV| = |AW| = 2$, a~teda aj $|BU| = |BV| = 2 - c$ a~$|CU| = |CW| = 2 - b$.

    Body $P$, $V$ ležia na tej istej polpriamke z~$A$ a~$|AP| = 1 = \tfrac12|AV|$, takže $P$ je stredom úsečky $AV$; rovnako $Q$ je stredom úsečky $AW$. Priamka $PQ$ je teda obrazom priamky $VW$, čiže poláry bodu $A$ vzhľadom na $\omega$, v~rovnoľahlosti so stredom $A$ a~koeficientom $\tfrac12$. To nie je náhoda a~dá sa to povedať oveľa užitočnejšie: priamka $PQ$ je chordálou kružnice $\omega$ a~kružnice s~nulovým polomerom a~stredom $A$.

    Mocnosť bodu $M$ vzhľadom na túto degenerovanú kružnicu je $|MA|^2$ a~štandardný dôkaz, že body s~rovnakou mocnosťou vzhľadom na dve nesústredné kružnice tvoria priamku kolmú na spojnicu ich stredov, funguje aj pre nulový polomer. Bod $A$ leží zvonka $\omega$, lebo $p(A,\omega) = |AV|^2 = 4 > 0$, takže nie je stredom $\omega$ a~táto chordála naozaj je priamka. Priamka $AB$ sa dotýka $\omega$ v~bode $V$, preto
    $$
    p(P,\omega) = |PV|^2 = |PA|^2,
    $$
    kde druhá rovnosť platí práve preto, že $P$ je stredom úsečky $AV$; rovnako $p(Q,\omega) = |QW|^2 = |QA|^2$. Body $P$, $Q$ sú rôzne, lebo ležia vo vzdialenosti 1 na rôznych polpriamkach z~$A$, a~tak je hľadanou chordálou práve priamka $PQ$.

    Bod $X$ leží na priamke $PQ$, čiže $|XA|^2 = p(X,\omega)$. Zároveň leží na priamke $BC$, ktorá sa dotýka $\omega$ v~bode $U$, takže $p(X,\omega) = |XU|^2$. Dostávame
    $$
    |XA| = |XU|.
    $$

    Úsečka $AX$ delí trojuholník $ABC$ na trojuholníky $ABX$ a~$ACX$. Body $U$ aj $X$ ležia na úsečke $BC$, takže $X$ leží na úsečke $BU$ alebo na úsečke $UC$. V~prvom prípade je $|BX| + |XU| = |BU|$ a~obvod trojuholníka $ABX$ je
    $$
    |AB| + |BX| + |XA| = c + |BX| + |XU| = c + |BU| = 2.
    $$
    V~druhom prípade je $|CX| + |XU| = |CU|$ a~obvod trojuholníka $ACX$ je
    $$
    |AC| + |CX| + |XA| = b + |CX| + |XU| = b + |CU| = 2.
    $$
}

\EnvId{hWnbfKi7JKW34d1E7tUxF-kruznice-nad-stranami-stvoruholnika}
\Problem{0}{}{
    Štvoruholník $ABCD$ spĺňa $a + b = c + d$, kde $a, b, c, d$ sú dĺžky jeho strán. Nad každou z~jeho strán ako priemerom sú zostrojené kružnice. Dokážte, že existuje kružnica dotýkajúca sa všetkých štyroch z~nich.
}{
    Predstavte si želanú magickú kružnicu a~skúste uhádnuť, kde by mohol byť jej stred. Zamyslite sa nad tým, kde sú stredy našich štyroch Tálesových kružníc.
}{
    Stredom kružnice bude stred jednej z~uhlopriečok; funguje len jedna z~nich, tak skúste obe a~nevzdávajte sa.
}{
    Vzdialenosti z~tohto stredu do stredov Tálesových kružníc vyplývajú z~vety o~strednej priečke. Zostáva vypočítať želaný polomer a~overiť, že naša magická podmienka ho zaručuje.
}{
    Nech $a = |AB|$, $b = |BC|$, $c = |CD|$, $d = |DA|$ a~nech $P$, $Q$, $R$, $S$ sú postupne stredy strán $AB$, $BC$, $CD$, $DA$. Tálesove kružnice zo zadania sú teda $\omega_P$ so stredom $P$ a~polomerom $a/2$, $\omega_Q$ so stredom $Q$ a~polomerom $b/2$, $\omega_R$ so stredom $R$ a~polomerom $c/2$ a~$\omega_S$ so stredom $S$ a~polomerom $d/2$.

    Stred hľadanej kružnice musíme z~niečoho vyrobiť a~jediné, čo o~daných štyroch kružniciach vieme, sú polohy ich stredov: sú to stredy strán. Keď je obrázok plný stredov úsečiek, oplatí sa pridať ďalší stred, aby vznikli stredné priečky, a~prirodzenými kandidátmi sú stredy uhlopriečok. Uhlopriečka $AC$ pritom rozdeľuje strany na dvojice $\{AB, BC\}$ a~$\{CD, DA\}$, čiže presne tak, ako ich zoskupuje predpoklad $a + b = c + d$. Skúsme teda jej stred.

    Nech $O$ je stred úsečky $AC$. V~trojuholníku $ABC$ sú $O$, $P$ stredy strán $CA$, $AB$ a~$O$, $Q$ stredy strán $AC$, $CB$; v~trojuholníku $ACD$ sú $O$, $R$ stredy strán $CA$, $CD$ a~$O$, $S$ stredy strán $AC$, $AD$. Z~vety o~strednej priečke tak dostávame
    $$
    \gather
    |OP| = \tfrac12 |BC| = \tfrac b2, \qquad |OQ| = \tfrac12 |AB| = \tfrac a2, \\
    |OR| = \tfrac12 |DA| = \tfrac d2, \qquad |OS| = \tfrac12 |CD| = \tfrac c2.
    \endgather
    $$
    O~tvare štvoruholníka sme pritom nič nepredpokladali, všetko ďalšie je už len práca s~týmito štyrmi vzdialenosťami.

    Podstatné je kríženie: vzdialenosť z~$O$ do stredu kružnice nad stranou $AB$ je polovica dĺžky $BC$ a~naopak. Preto
    $$
    \gather
    |OP| + \tfrac a2 = |OQ| + \tfrac b2 = \tfrac{a+b}2, \\
    |OR| + \tfrac c2 = |OS| + \tfrac d2 = \tfrac{c+d}2,
    \endgather
    $$
    a~vďaka predpokladu $a + b = c + d$ sú všetky štyri hodnoty rovnaké. Označme túto spoločnú hodnotu $\rho = (a+b)/2 = (c+d)/2$ (je to štvrtina obvodu štvoruholníka) a~nech $\Omega$ je kružnica so stredom $O$ a~polomerom $\rho$.

    Zostáva overiť dotyk. Ak majú dve kružnice polomery $\rho > r > 0$ a~ich stredy sú vo vzdialenosti $\rho - r > 0$, dotýkajú sa vnútorne: menšia leží vnútri väčšej a~jediný spoločný bod je priesečník polpriamky zo stredu väčšej do stredu menšej s~väčšou kružnicou. (Vonkajší dotyk by znamenal vzdialenosť stredov $\rho + r$, o~ten tu nejde.) Predošlé rovnosti hovoria práve toto:
    $$
    |OP| = \rho - \tfrac a2 = \tfrac b2 > 0, \qquad |OQ| = \rho - \tfrac b2 = \tfrac a2 > 0,
    $$
    $$
    |OR| = \rho - \tfrac c2 = \tfrac d2 > 0, \qquad |OS| = \rho - \tfrac d2 = \tfrac c2 > 0,
    $$
    pričom všetky štyri polomery $a/2$, $b/2$, $c/2$, $d/2$ sú kladné a~menšie než $\rho$. Kružnica $\Omega$ sa teda vnútorne dotýka všetkých štyroch Tálesových kružníc, ktoré ležia v~jej vnútri.

    Na záver ešte, prečo nefunguje druhá uhlopriečka. Ak je $O'$ stred úsečky $BD$, tá istá úvaha v~trojuholníkoch $ABD$ a~$CBD$ dá $|O'P| = d/2$, $|O'S| = a/2$, $|O'Q| = c/2$ a~$|O'R| = b/2$, takže tentoraz
    $$
    \gather
    |O'P| + \tfrac a2 = |O'S| + \tfrac d2 = \tfrac{a+d}2, \\
    |O'Q| + \tfrac b2 = |O'R| + \tfrac c2 = \tfrac{b+c}2.
    \endgather
    $$
    Kružnica so stredom $O'$ dotýkajúca sa vnútorne všetkých štyroch by si teda vyžadovala $a + d = b + c$, čo je iná podmienka než tá zadaná. Uhlopriečka $BD$ zoskupuje strany do dvojíc $\{AB, DA\}$ a~$\{BC, CD\}$, a~to je zoskupenie, ktoré nám predpoklad nedáva.
}

\EnvId{KWvpK6SAI-wBt88-OpQjm-sestuholnik-s-vpisanou-kruznicou}
\Problem{0}{}{
    Ak je do šesťuholníka~$ABCDEF$ vpísaná kružnica, tak sa uhlopriečky~$AD$, $BE$, $CF$ pretínajú v~jednom bode.
}{
    Nakreslite tri kružnice tak, aby uhlopriečky $AD$, $BE$, $CF$ boli chordálami. Je to záludné.
}{
    Jedna z~kružníc sa dotýka polpriamok opačných k $BA$ a $DE$. Analogicky ostatné. Jediný trik je umiestniť ich tak, aby chordály fungovali. Na to určite potrebujeme použiť kružnicu vpísanú šesťuholníku $ABCDEF$.
}{
    Tri priamky prechádzajúce jedným bodom volajú po potenčnom strede, hľadáme teda tri kružnice, ktorých chordálami sú práve $AD$, $BE$, $CF$; jediné, čo v~úlohe máme, je vpísaná kružnica, a~mocnosť vrcholu vieme uchopiť vtedy, keď z~neho vedie dotyčnica.

    Nech je $ABCDEF$ konvexný a~nech sa jeho vpísaná kružnica $\omega$ so stredom $O$ a~polomerom $r$ dotýka strán $AB$, $BC$, $CD$, $DE$, $EF$, $FA$ postupne v~bodoch $P$, $Q$, $R$, $S$, $T$, $U$; v~tomto poradí ležia tieto body aj na $\omega$, v~rovnakom smere obehu ako vrcholy. Dĺžky dotyčníc z~vrcholov označme
    $$
    \gather
    a = |AP| = |AU|, \quad b = |BP| = |BQ|, \quad c = |CQ| = |CR|,\\
    d = |DR| = |DS|, \quad e = |ES| = |ET|, \quad f = |FT| = |FU|.
    \endgather
    $$

    Pre bod $M$ na dotyčnici kružnice $\omega'$ s~dotykovým bodom $Z$ je $p(M,\omega') = |MZ|^2$, kde $p(M,\omega')$ značí mocnosť bodu $M$ ku kružnici $\omega'$. Nech sa preto kružnica $\omega_1$ dotýka priamok protiľahlých strán $AB$, $DE$ v~bodoch $P_1$, $S_1$, kružnica $\omega_2$ priamok $BC$, $EF$ v~bodoch $Q_1$, $T_1$ a~kružnica $\omega_3$ priamok $CD$, $FA$ v~bodoch $R_1$, $U_1$. Každý vrchol potom leží na dotyčnici práve dvoch z~nich, napríklad $B$ na dotyčniciach $AB$, $BC$ kružníc $\omega_1$, $\omega_2$, a~priamka $BE$ bude chordálou kružníc $\omega_1$, $\omega_2$, len čo $|BP_1| = |BQ_1|$ a~$|ES_1| = |ET_1|$. Chceme teda, aby v~každom vrchole boli oba tamojšie dotykové body od neho rovnako vzdialené.

    Pôvodné dotykové body to spĺňajú, lenže dávajú $\omega_1 = \omega_2 = \omega_3 = \omega$. Posunieme ich preto po stranách o~tú istú dĺžku, striedavo v~smere obehu a~proti nemu: na každej strane leží práve jeden z~vrcholov $B$, $D$, $F$ a~dotykový bod posunieme smerom k~nemu, a~to až zaň. Zvoľme $t > \max(b,d,f)$ a~nech body $P_1$, $Q_1$ ležia na polpriamkach opačných k~$BA$, $BC$ vo vzdialenosti $t-b$ od $B$, body $R_1$, $S_1$ na polpriamkach opačných k~$DC$, $DE$ vo vzdialenosti $t-d$ od $D$ a~body $T_1$, $U_1$ na polpriamkach opačných k~$FE$, $FA$ vo vzdialenosti $t-f$ od $F$. Potom
    $$
    \gather
    |AP_1| = |AU_1| = a+t, \quad |BP_1| = |BQ_1| = t-b,\\
    |CQ_1| = |CR_1| = c+t, \quad |DR_1| = |DS_1| = t-d,\\
    |ES_1| = |ET_1| = e+t, \quad |FT_1| = |FU_1| = t-f,
    \endgather
    $$
    napríklad $|AP_1| = |AB| + |BP_1| = (a+b)+(t-b)$ a~$|CQ_1| = |CB| + |BQ_1| = (b+c)+(t-b)$.

    Ostáva ukázať, že kružnica dotýkajúca sa $AB$ v~bode $P_1$ a~$DE$ v~bode $S_1$ vôbec existuje; práve tu sa striedanie smerov spotrebuje. Nech $m_1$ je os úsečky $PS$ a~$\sigma$ osová súmernosť podľa nej. Z~$|OP| = |OS| = r$ prechádza $m_1$ bodom $O$, takže $\sigma$ zobrazuje $\omega$ na seba, bod $P$ na bod $S$, a~teda aj dotyčnicu $AB$ na dotyčnicu $DE$. Keďže je šesťuholník konvexný, leží pri obehu $A$, $B$, $C$, $D$, $E$, $F$ stred $O$ pri každej strane po tej istej ruke, povedzme vľavo. Súmernosť je nepriama zhodnosť a~ľavú stranu vymieňa za pravú: smer $A \to B$, ktorý má v~bode $P$ stred $O$ vľavo, sa preto zobrazí na ten smer priamky $DE$, ktorý má v~bode $S$ stred $O$ vpravo, čiže na smer $E \to D$. Body $P_1$, $S_1$ pritom vznikli posunom o~$t$ práve v~týchto dvoch smeroch, takže $\sigma(P_1) = S_1$.

    Priamka $m_1$ nie je kolmá na $AB$: inak by splynula s~priamkou $OP$, jedinou kolmicou na $AB$ vedenou bodom $O$, a~$\sigma$ by nechala $P$ na mieste, čiže $P = S$. Kolmica na $AB$ v~bode $P_1$ preto pretne $m_1$ v~jedinom bode $O_1$; ten $\sigma$ nechá na mieste a~túto kolmicu zobrazí na kolmicu na $DE$ v~bode $S_1$, takže $|O_1P_1| = |O_1S_1|$ a~kružnica $\omega_1$ so stredom $O_1$ a~polomerom $|O_1P_1|$ je hľadaná. Nulový polomer by znamenal $O_1 = P_1 = S_1$, teda že $P_1$ je priesečníkom rôznych priamok $AB$, $DE$ (rôzne sú, lebo sa $\omega$ dotýkajú v~rôznych bodoch); pri danej dvojici protiľahlých strán to vylučuje najviac jednu hodnotu $t$, spolu teda najviac tri, a~$t$ volíme mimo nich. Napokon $O_1 \neq O$, lebo $O$ leží na kolmici na $AB$ v~bode $P$ a~$O_1$ na rovnobežnej kolmici v~bode $P_1 \neq P$. Rovnako zostrojíme $\omega_2$ so stredom $O_2$ na osi $m_2$ úsečky $QT$ a~$\omega_3$ so stredom $O_3$ na osi $m_3$ úsečky $RU$.

    Kružnice sú po dvoch nesústredné. Osi $m_1$, $m_2$, $m_3$ prechádzajú bodom $O$ a~sú navzájom rôzne: body $P$, $Q$, $S$, $T$ ležia na $\omega$ v~tomto cyklickom poradí, takže sa tetivy $PS$ a~$QT$ pretínajú a~nie sú rovnobežné, a~teda nesplývajú ani kolmice na ne vedené bodom $O$, čiže $m_1$ a~$m_2$; analogicky pre dvojice $QT$, $RU$ a~$RU$, $PS$. Priamky $m_1$, $m_2$, $m_3$ tak majú jediný spoločný bod $O$, a~keďže $O_i$ leží na $m_i$ a~$O_i \neq O$, sú stredy $O_1$, $O_2$, $O_3$ navzájom rôzne.

    Vrcholy ležia na dotyčniciach príslušných kružníc, takže
    $$
    \gather
    p(B,\omega_1) = (t-b)^2 = p(B,\omega_2),\\
    p(E,\omega_1) = (e+t)^2 = p(E,\omega_2).
    \endgather
    $$
    Rôzne body $B$, $E$ teda ležia na chordále kružníc $\omega_1$, $\omega_2$, tou je preto priamka $BE$. Analogicky hodnoty $(c+t)^2$ v~$C$ a~$(t-f)^2$ v~$F$ robia $CF$ chordálou kružníc $\omega_2$, $\omega_3$ a~hodnoty $(a+t)^2$ v~$A$ a~$(t-d)^2$ v~$D$ robia $AD$ chordálou kružníc $\omega_1$, $\omega_3$.

    Tri po dvoch nesústredné kružnice majú buď potenčný stred, ktorým prechádzajú všetky tri chordály, alebo ich stredy ležia na jednej priamke a~potom sú všetky tri chordály na túto priamku kolmé, teda rovnobežné, prípadne splynú. Druhá možnosť odpadá: uhlopriečka $AD$ delí konvexný šesťuholník na $ABCD$ a~$ADEF$, takže $B$ a~$E$ ležia v~opačných polrovinách určených priamkou $AD$ a~úsečka $BE$ ju pretína v~jedinom bode. Uhlopriečky $AD$, $BE$, $CF$ teda prechádzajú jedným bodom, potenčným stredom kružníc $\omega_1$, $\omega_2$, $\omega_3$.
}

\DisplayHints

\DisplayProblemSolutions

\bye
